% Elementary Probability — Pure Mathematics with Statistics Upper Sixth — Cours signé OnBuch+
% Official MINESEC syllabus. Compiler avec : tectonic PMS10-elementary-probability.tex
\newcommand{\DOCMATIERE}{Pure Mathematics with Statistics}
\newcommand{\DOCNIVEAU}{Upper Sixth}
\newcommand{\DOCMODULE}{Upper Sixth — Pure mathematics with statistics}
\newcommand{\DOCLECON}{10}
\newcommand{\DOCTITRE}{Elementary Probability}
% OnBuch+ — préambule « cours signés » ENRICHI (compilé avec tectonic / XeLaTeX)
% Système visuel « Pop » d'OnBuch+ : fond crème (jamais blanc), bordures encre
% épaisses, ombres « dures » décalées, titres Archivo Black en capitales.
% Version MATHÉMATIQUES : graphes (pgfplots/tikz), boîtes pédagogiques
% (définition, propriété, méthode, attention, le savais-tu, exercice résolu,
% exercice, TP, bilan), page de garde et signature OnBuch+ ancrée en bas.
%
% Macros attendues dans main.tex AVANT \input{preamble} :
%   \DOCMATIERE  \DOCNIVEAU  \DOCMODULE  \DOCLECON (numéro)  \DOCTITRE
\documentclass[11pt]{article}

\usepackage{fontspec}
\usepackage[english]{babel}
\usepackage[a4paper,margin=2cm,top=3.4cm,bottom=2.8cm,headheight=1.4cm,footskip=1.3cm]{geometry}
\usepackage{amsmath,amssymb,amsfonts}
\usepackage{mathtools} % \xmapsto, \coloneqq... souvent utilisés par les modèles
\usepackage{cancel} % simplifications barrées (bilans de forces)
% \abs{z} (module, valeur absolue) et \norm{u} (norme), utilisés par les modèles
\providecommand{\abs}{}\let\abs\relax\DeclarePairedDelimiter{\abs}{\lvert}{\rvert}
\providecommand{\norm}{}\let\norm\relax\DeclarePairedDelimiter{\norm}{\lVert}{\rVert}
\usepackage[table]{xcolor} % [table] : fournit aussi \rowcolors (alternance de lignes)
\usepackage{pagecolor}
\usepackage{tikz}
\usetikzlibrary{arrows.meta,positioning,calc,shapes.geometric,shapes.misc,shapes.arrows,decorations.pathmorphing,decorations.pathreplacing,decorations.markings,patterns,shadows,mindmap,trees,fit,backgrounds,matrix,intersections,angles,quotes}
\usepackage{pgfplots}
\usepackage[version=4]{mhchem} % équations nucléaires \ce{^{235}_{92}U}
\usepackage[european,siunitx]{circuitikz} % schémas électriques
\pgfplotsset{compat=1.18}
\usepgfplotslibrary{fillbetween,groupplots} % aires entre courbes (calcul intégral)
\usepackage{enumitem}
\usepackage{fancyhdr}
% [most] charge toutes les bibliothèques tcolorbox (listings, minted, raster,
% documentation...) et gonfle inutilement la mémoire TeX sur un document long
% (~300 boîtes) : on ne charge que ce qui est réellement utilisé.
\usepackage[skins,breakable]{tcolorbox}
\usepackage{graphicx}
\usepackage{colortbl}
\usepackage{booktabs}
\usepackage{tabularx}
\usepackage{diagbox} % case d'en-tête diagonale des tableaux à double entrée
% Colonnes extensibles centrée / gauche / droite (C, L, R), souvent utilisées par les modèles
\newcolumntype{C}{>{\centering\arraybackslash}X}
\newcolumntype{L}{>{\raggedright\arraybackslash}X}
\newcolumntype{R}{>{\raggedleft\arraybackslash}X}
\usepackage{array}
\usepackage{multicol}
\usepackage{siunitx}
% parse-numbers=false : accepte aussi \SI{2,5 \times 10^{-3}}{...}
\sisetup{locale=UK,output-decimal-marker={.},group-minimum-digits=4,parse-numbers=false}
\DeclareSIUnit\ohm{\ensuremath{\symup{\Omega}}} % U+2126 absent de la police
\DeclareSIUnit\division{div} % réglages d'oscilloscope (ms/div, V/div)
\DeclareSIUnit\tr{tr} % tours (vitesse de rotation, ex. \SI{1500}{\tr\per\minute})
\usepackage{qrcode}
% \degree : commande fréquemment utilisée par les modèles pour un angle ou une
% température en dehors de \SI{}{} (non fournie par les paquets chargés).
\providecommand{\degree}{^{\circ}}
% \qed (fin de démonstration), utilisé par les modèles sans amsthm
\providecommand{\qed}{\hfill$\square$}
% \barre : double barre des valeurs interdites dans un tableau de variations (tkz-tab)
\providecommand{\barre}{\ensuremath{\|}}
% environnement proof (amsthm non chargé) utilisé par les modèles
\ifdefined\proof\else\newenvironment{proof}[1][Proof]{\par\noindent\textit{#1.}\ }{\qed\par}\fi
\usepackage{lastpage}
\usepackage{hyperref}
\hypersetup{hidelinks}

% ============================================================
% Palette « Pop » OnBuch+ — mêmes hex que la classe P (plus_ui.dart)
% ============================================================
\definecolor{poporange}{HTML}{F59321}
\definecolor{popdark}{HTML}{E07A0C}
\definecolor{popcream}{HTML}{FFF6E8}
\definecolor{popink}{HTML}{1C1714}
\definecolor{poppurple}{HTML}{7A5AE0}
\definecolor{popgreen}{HTML}{2E9E6A}
\definecolor{poppink}{HTML}{E0457B}
\definecolor{popblue}{HTML}{2F7DE1}
\definecolor{popgold}{HTML}{C98A12}
\definecolor{popmuted}{HTML}{8A7F76}
\definecolor{popline}{HTML}{EADFD2}
\definecolor{popgray}{HTML}{8A7F76}
\colorlet{popgrey}{popgray}
% Alias tolérants (le modèle nomme parfois les couleurs différemment)
\colorlet{popwhite}{white}
\colorlet{popblack}{popink}
\colorlet{popred}{poppink}
\colorlet{popyellow}{popgold}
% Teintes claires pour fonds de tableaux / zones
\colorlet{popblueL}{popblue!10!white}
\colorlet{poporangeL}{poporange!14!white}
\colorlet{popgreenL}{popgreen!12!white}
\colorlet{poppurpleL}{poppurple!12!white}
\colorlet{poppinkL}{poppink!12!white}
\colorlet{popgoldL}{popgold!14!white}

\pagecolor{popcream}

% ============================================================
% Typographie
% ============================================================
\defaultfontfeatures{Ligatures=TeX}
\setmainfont{PlusJakartaSans-Medium.ttf}[Path=../../../fonts/,BoldFont=PlusJakartaSans-Bold.ttf]
% Maths en sans-serif (Fira Math) avec chiffres et lettres droites de Plus
% Jakarta, pour que les formules se fondent dans le texte.
\usepackage[math-style=ISO,bold-style=ISO]{unicode-math}
\setmathfont{FiraMath-Regular.otf}[Path=../../../fonts/]
\setmathfont{PlusJakartaSans-Medium.ttf}[Path=../../../fonts/,range={up/num,up/{Latin,latin},\mathup}]
% (icomma retiré : décimales avec point en anglais)
% Caractères mathématiques Unicode absents des polices de texte (Plus Jakarta,
% Archivo Black) : rendus par leur équivalent mathématique, en texte comme en
% formule — sinon ils s'affichent en carré vide (ex. « Arithmétique dans ℤ »).
\usepackage{newunicodechar}
\newunicodechar{ℝ}{\ensuremath{\mathbb{R}}}
\newunicodechar{ℤ}{\ensuremath{\mathbb{Z}}}
\newunicodechar{ℂ}{\ensuremath{\mathbb{C}}}
\newunicodechar{ℕ}{\ensuremath{\mathbb{N}}}
\newunicodechar{ℚ}{\ensuremath{\mathbb{Q}}}
\newunicodechar{𝔻}{\ensuremath{\mathbb{D}}}
\newunicodechar{α}{\ensuremath{\alpha}}
\newunicodechar{β}{\ensuremath{\beta}}
\newunicodechar{γ}{\ensuremath{\gamma}}
\newunicodechar{δ}{\ensuremath{\delta}}
\newunicodechar{ε}{\ensuremath{\varepsilon}}
\newunicodechar{θ}{\ensuremath{\theta}}
\newunicodechar{λ}{\ensuremath{\lambda}}
\newunicodechar{μ}{\ensuremath{\mu}}
\newunicodechar{σ}{\ensuremath{\sigma}}
\newunicodechar{τ}{\ensuremath{\tau}}
\newunicodechar{φ}{\ensuremath{\varphi}}
\newunicodechar{ψ}{\ensuremath{\psi}}
\newunicodechar{ω}{\ensuremath{\omega}}
\newunicodechar{Σ}{\ensuremath{\Sigma}}
% majuscules produites par \MakeUppercase dans les titres (α-aminés -> Α)
\newunicodechar{Α}{\ensuremath{\alpha}}
\newunicodechar{Β}{\ensuremath{\beta}}
\newunicodechar{Γ}{\ensuremath{\Gamma}}
\newunicodechar{Δ}{\ensuremath{\Delta}}
\newunicodechar{Ε}{\ensuremath{\varepsilon}}
\newunicodechar{Θ}{\ensuremath{\Theta}}
\newunicodechar{Λ}{\ensuremath{\Lambda}}
\newunicodechar{Μ}{\ensuremath{\mu}}
\newunicodechar{Τ}{\ensuremath{\tau}}
\newunicodechar{Φ}{\ensuremath{\Phi}}
\newunicodechar{Ψ}{\ensuremath{\Psi}}
\newunicodechar{Ω}{\ensuremath{\Omega}}
\newunicodechar{↦}{\ensuremath{\mapsto}}
\newunicodechar{∈}{\ensuremath{\in}}
\newunicodechar{∉}{\ensuremath{\notin}}
\newunicodechar{∧}{\ensuremath{\wedge}}
\newunicodechar{⇔}{\ensuremath{\Leftrightarrow}}
\newunicodechar{⟺}{\ensuremath{\Longleftrightarrow}}
\newunicodechar{⇒}{\ensuremath{\Rightarrow}}
\newunicodechar{⟹}{\ensuremath{\Longrightarrow}}
\newunicodechar{∘}{\ensuremath{\circ}}
\newunicodechar{∪}{\ensuremath{\cup}}
\newunicodechar{∩}{\ensuremath{\cap}}
\newunicodechar{≡}{\ensuremath{\equiv}}
\newunicodechar{⊂}{\ensuremath{\subset}}
\newunicodechar{⊥}{\ensuremath{\perp}}
\newunicodechar{∃}{\ensuremath{\exists}}
\newunicodechar{∀}{\ensuremath{\forall}}
\newunicodechar{∛}{\ensuremath{\sqrt[3]{\,}}}
\newunicodechar{∜}{\ensuremath{\sqrt[4]{\,}}}
\newunicodechar{⃗}{\ensuremath{{}^{\to}}}
\newunicodechar{ⁿ}{\ifmmode^{n}\else\textsuperscript{\ensuremath{n}}\fi}
\newunicodechar{ˣ}{\ifmmode^{x}\else\textsuperscript{\ensuremath{x}}\fi}
\newunicodechar{ᵇ}{\ifmmode^{b}\else\textsuperscript{\ensuremath{b}}\fi}
\newunicodechar{ʳ}{\ifmmode^{r}\else\textsuperscript{\ensuremath{r}}\fi}
\newunicodechar{ᵘ}{\ifmmode^{u}\else\textsuperscript{\ensuremath{u}}\fi}
\newunicodechar{ʸ}{\ifmmode^{y}\else\textsuperscript{\ensuremath{y}}\fi}
\newunicodechar{ᵃ}{\ifmmode^{a}\else\textsuperscript{\ensuremath{a}}\fi}
\newunicodechar{ᵉ}{\ifmmode^{e}\else\textsuperscript{\ensuremath{e}}\fi}
\newunicodechar{ᵏ}{\ifmmode^{k}\else\textsuperscript{\ensuremath{k}}\fi}
\newunicodechar{ᵖ}{\ifmmode^{p}\else\textsuperscript{\ensuremath{p}}\fi}
\newunicodechar{ᴺ}{\ifmmode^{N}\else\textsuperscript{\ensuremath{N}}\fi}
\newunicodechar{ᶜ}{\ifmmode^{c}\else\textsuperscript{\ensuremath{c}}\fi}
\newunicodechar{ʰ}{\ifmmode^{h}\else\textsuperscript{\ensuremath{h}}\fi}
\newunicodechar{ᵠ}{\ifmmode^{q}\else\textsuperscript{\ensuremath{q}}\fi}
\newunicodechar{ₙ}{\ifmmode_{n}\else\textsubscript{\ensuremath{n}}\fi}
\newunicodechar{ₐ}{\ifmmode_{a}\else\textsubscript{\ensuremath{a}}\fi}
\newunicodechar{ᵢ}{\ifmmode_{i}\else\textsubscript{\ensuremath{i}}\fi}
\newunicodechar{ᵣ}{\ifmmode_{r}\else\textsubscript{\ensuremath{r}}\fi}
\newunicodechar{ₖ}{\ifmmode_{k}\else\textsubscript{\ensuremath{k}}\fi}
\newunicodechar{ᵦ}{\ifmmode_{\beta}\else\textsubscript{\ensuremath{\beta}}\fi}
\newunicodechar{ₓ}{\ifmmode_{x}\else\textsubscript{\ensuremath{x}}\fi}
\newfontfamily\popdisplay{ArchivoBlack-Regular.ttf}[Path=../../../fonts/]
\newfontfamily\poplogo{SpaceGrotesk-Bold.ttf}[Path=../../../fonts/]
\newfontfamily\popsemi{PlusJakartaSans-SemiBold.ttf}[Path=../../../fonts/]
\newfontfamily\popxbold{PlusJakartaSans-ExtraBold.ttf}[Path=../../../fonts/]
\newfontfamily\popxboldls{PlusJakartaSans-ExtraBold.ttf}[Path=../../../fonts/,LetterSpace=3.0]

\setlist[enumerate]{leftmargin=1.7em,itemsep=5pt,topsep=4pt}
\setlist[itemize]{leftmargin=1.7em,itemsep=4pt,topsep=4pt,label={\color{poporange}\textbullet}}
\setlength{\parindent}{0pt}
\setlength{\parskip}{5pt}
\renewcommand{\arraystretch}{1.25}

% pgfplots : style Pop par défaut
\pgfplotsset{
  every axis/.append style={axis line style={popink,line width=1pt},tick style={popink},
    grid style={popline,line width=0.5pt},label style={font=\small},tick label style={font=\footnotesize}},
  popcurve/.style={poporange,line width=1.8pt,no markers,smooth},
}
% Même style disponible dans un \draw TikZ simple (hors pgfplots)
\tikzset{popcurve/.style={poporange,line width=1.8pt}}
\tikzset{popgrid/.style={popline,very thin}}
\colorlet{popcurve}{poporange} % le nom du style est parfois utilisé comme couleur

% ============================================================
% Pill / squiggle
% ============================================================
\newcommand{\poppill}[3]{%
  \tikz[baseline=-0.55ex]\node[draw=popink,line width=1.2pt,rounded corners=2.6pt,%
    fill=#2,inner xsep=6.5pt,inner ysep=3pt]{%
    \popxboldls\fontsize{7}{7}\selectfont\color{#3}\MakeUppercase{#1}};%
}
\newcommand{\popsquiggle}[1]{%
  \tikz[baseline=-0.4ex]{
    \draw[#1,line width=2.6pt,line cap=round]
      plot [smooth] coordinates {(0,0.09) (0.34,0.01) (0.68,0.21) (1.02,0.01) (1.36,0.10)};
  }%
}
% Mot-clé surligné (marqueur orange sous le texte)
\newcommand{\cle}[1]{{\popxbold\color{popdark}#1}}

% ============================================================
% En-tête / pied
% ============================================================
\pagestyle{fancy}
\fancyhf{}
\renewcommand{\headrulewidth}{0pt}
\renewcommand{\footrulewidth}{0pt}
\lhead{\raisebox{-0.15ex}{\poplogo\fontsize{13}{13}\selectfont\color{popink}OnBuch\color{poporange}+}}
\rhead{\poppill{\DOCMATIERE\ \textbullet\ \DOCNIVEAU\ \textbullet\ Chapter \DOCLECON}{white}{popink}}
\lfoot{\popxbold\fontsize{7.6}{7.6}\selectfont\color{popmuted}\MakeUppercase{Signed course OnBuch+}\ \textbullet\ {\popsemi\DOCTITRE}}
\rfoot{\tikz[baseline=-0.6ex]\node[rounded corners=8.5pt,fill=popink,minimum height=17pt,minimum width=17pt,inner xsep=5pt,inner sep=0]{\popxbold\fontsize{8}{8}\selectfont\color{popcream}\thepage\,/\,\pageref*{LastPage}};}

% ============================================================
% Titres
% ============================================================
\newcommand{\chaptitre}[2]{%
  \begin{tcolorbox}[enhanced,colback=poporange,colframe=popink,boxrule=2.6pt,arc=11pt,
      drop shadow=popink,left=15pt,right=15pt,top=12pt,bottom=11pt,before skip=3pt,after skip=18pt]
    \poppill{#2}{popink}{popcream}\par\vspace{9pt}
    {\raggedright\hyphenpenalty=10000\exhyphenpenalty=10000\popdisplay\fontsize{22}{24}\selectfont\color{popcream}\MakeUppercase{#1}\par}
  \end{tcolorbox}
}
% Section numérotée : pastille orange avec le numéro + titre Archivo
\newcounter{popsec}
\newcommand{\coursec}[1]{%
  \stepcounter{popsec}\setcounter{popsub}{0}%
  \par\vspace{16pt}\noindent
  \begin{minipage}{\linewidth}
  \tikz[baseline=-0.7ex]\node[draw=popink,line width=1.6pt,fill=poporange,rounded corners=4pt,minimum size=22pt,inner sep=0]{\popdisplay\fontsize{12}{12}\selectfont\color{popcream}\thepopsec};%
  \hspace{8pt}{\popdisplay\fontsize{15}{17}\selectfont\color{popink}\MakeUppercase{#1}}\par
  \vspace{4pt}\noindent{\color{poporange}\rule{36pt}{3.6pt}}
  \end{minipage}\par\nopagebreak\vspace{8pt}
}
\newcounter{popsub}
\newcommand{\courssub}[1]{%
  \stepcounter{popsub}%
  \par\vspace{9pt}\noindent{\popxbold\fontsize{11.5}{13}\selectfont\color{popink}{\color{poporange}\thepopsec.\thepopsub}\hspace{6pt}#1}\par\nopagebreak\vspace{4pt}
}

% ============================================================
% Boîtes pédagogiques — même recette : fond blanc, bordure épaisse colorée,
% ombre dure encre, badge plein. Titre optionnel [#1].
% ============================================================
\tcbset{popbase/.style={breakable,enhanced,colback=white,boxrule=2.3pt,arc=9pt,
  shadow={3pt}{-3pt}{0pt}{popink},left=14pt,right=14pt,top=11pt,bottom=11pt,before skip=11pt,after skip=15pt,
  fontupper=\popsemi\fontsize{10.6}{14.5}\selectfont\color{popink},
  fontlower=\popsemi\fontsize{10.6}{14.5}\selectfont\color{popink}}}
% Coche dessinée (le glyphe ✓ n'existe pas dans les polices du document)
\newcommand{\popcheck}[1]{\tikz[baseline=-0.5ex]\draw[#1,line width=1.6pt,line cap=round,line join=round](-0.13,0.0)--(-0.04,-0.09)--(0.13,0.1);}
\providecommand{\coche}{\popcheck{popgreen}}
\providecommand{\resultat}[1]{\par\smallskip\noindent\fcolorbox{popgreen}{popgreenL}{\parbox{\dimexpr\linewidth-2\fboxsep-2\fboxrule\relax}{\textbf{Result.} #1}}\par\smallskip}
\newcommand{\popboxhead}[3]{\poppill{#1}{#2}{white}\ifx\relax#3\relax\else\hspace{6pt}{\popxbold\fontsize{10.6}{12}\selectfont\color{popink}#3}\fi\par\vspace{7pt}}

\newtcolorbox{aretenir}[1][]{popbase,colframe=popgreen,before upper={\popboxhead{Key points}{popgreen}{#1}}}
\newtcolorbox{exemplebox}[1][]{popbase,colframe=poporange,before upper={\popboxhead{Exemple}{poporange}{#1}}}
\newtcolorbox{definition}[1][]{popbase,colframe=popblue,before upper={\popboxhead{Definition}{popblue}{#1}}}
\newtcolorbox{propriete}[1][]{popbase,colframe=poppurple,before upper={\popboxhead{Property}{poppurple}{#1}}}
\newtcolorbox{methode}[1][]{popbase,colframe=popgold,before upper={\popboxhead{Method}{popgold}{#1}}}
\newtcolorbox{attention}[1][]{popbase,colframe=poppink,before upper={\popboxhead{Warning}{poppink}{#1}}}
\newtcolorbox{experience}[1][]{popbase,colframe=popblue,colback=popblueL,before upper={\popboxhead{Experiment}{popblue}{#1}}}
% « Le savais-tu ? » : carte violette pleine (contexte, culture, Cameroun)
\newtcolorbox{savaistu}[1][]{popbase,colback=poppurple,colframe=popink,
  fontupper=\popsemi\fontsize{10.4}{14.2}\selectfont\color{white},
  before upper={\poppill{Did you know?}{white}{poppurple}\ifx\relax#1\relax\else\hspace{6pt}{\popxbold\color{white}#1}\fi\par\vspace{7pt}}}
% Exercice résolu : énoncé en haut, \tcblower puis la solution sur fond crème
\newcounter{popexo}\newcounter{popexr}
\newtcolorbox{exoresolu}[1][]{popbase,colframe=poporange,colbacklower=popcream,
  segmentation style={popink,line width=1.2pt,dashed},
  before upper={\stepcounter{popexr}\popboxhead{Worked example \thepopexr}{poporange}{#1}},
  before lower={\poppill{Solution}{popink}{popcream}\par\vspace{6pt}}}
% Exercice (non résolu en place) — la correction est dans la partie Corrigés
\newtcolorbox{exercice}[1][]{popbase,colframe=popink,
  before upper={\stepcounter{popexo}\popboxhead{Exercise \thepopexo}{popink}{#1}}}
% Corrigé
\newtcolorbox{corrige}[1][]{popbase,colframe=popgreen,colback=popgreenL,
  before upper={\popboxhead{Solution}{popgreen}{#1}}}
% Objectifs / prérequis / bilan
\newtcolorbox{objectifs}{popbase,colframe=popink,colback=poporangeL,
  before upper={\popboxhead{Chapter objectives}{popink}{}}}
\newtcolorbox{prerequis}{popbase,colframe=popink,colback=popblueL,
  before upper={\popboxhead{Prerequisites}{popblue}{}}}
\newtcolorbox{bilan}{popbase,colframe=popink,colback=white,boxrule=2.8pt,
  before upper={\popboxhead{Fiche bilan}{poporange}{L'essentiel à connaître pour l'examen}}}
% Figure encadrée avec légende
\newtcolorbox{popfigure}[1][]{enhanced,breakable=false,colback=white,colframe=popline,boxrule=1.4pt,arc=8pt,
  left=10pt,right=10pt,top=10pt,bottom=8pt,before skip=10pt,after skip=12pt,halign=center,
  lower separated=false,halign lower=center,
  fontlower=\popsemi\fontsize{9}{11.5}\selectfont\color{popmuted},
  #1}
\newcommand{\legende}[1]{\par\vspace{6pt}{\popsemi\fontsize{9}{11.5}\selectfont\color{popmuted}#1}}

% ============================================================
% Signature OnBuch+
% ============================================================
\newcommand{\poplogoinline}[1]{{\poplogo\fontsize{#1}{#1}\selectfont\color{popink}OnBuch\color{poporange}+}}
\newcommand{\signatureonbuch}{%
  \begin{tcolorbox}[enhanced,colback=popink,colframe=popink,boxrule=2.2pt,arc=10pt,
      drop shadow=poporange,left=14pt,right=14pt,top=10pt,bottom=10pt,before skip=0pt,after skip=0pt]
    \tikz[baseline=-0.7ex]\node[circle,fill=popgreen,draw=popcream,line width=1.3pt,minimum size=22pt,inner sep=0]{\popcheck{white}};%
    \hspace{9pt}\begin{minipage}[c]{0.62\linewidth}
      {\popxbold\fontsize{12}{13}\selectfont\color{popcream}Signed\ \textbullet\ {\poplogo OnBuch\color{poporange}+}}\par\vspace{2pt}
      {\raggedright\popsemi\fontsize{8}{10.5}\selectfont\color{popline}\MakeUppercase{OnBuch+ teaching team}\\\MakeUppercase{Follows the official MINESEC syllabus}\par}
    \end{minipage}\hfill
    {\poplogo\fontsize{22}{22}\selectfont\color{popcream}OnBuch\color{poporange}+}
  \end{tcolorbox}%
}

% ============================================================
% Page de garde
% #1 titre · #2 matière · #3 niveau · #4 module · #5 n° leçon · #6 durée ·
% #7 description / objectifs (texte ou itemize)
% ============================================================
\newcommand{\pagegarde}[7]{%
  \thispagestyle{empty}
  \newgeometry{a4paper,margin=1.7cm,top=1.6cm,bottom=1.4cm}%
  \begin{tikzpicture}[remember picture,overlay]
    \fill[poporange!18] ([xshift=-3.2cm,yshift=-3.6cm]current page.north east) circle (4.6cm);
    \fill[poppurple!14] ([xshift=2.4cm,yshift=7.6cm]current page.south west) circle (3.8cm);
  \end{tikzpicture}%
  {\poplogo\fontsize{30}{30}\selectfont\color{popink}OnBuch\color{poporange}+}\hfill
  \raisebox{0.6ex}{\poppill{Signed course \textbullet\ Premium}{popink}{popcream}}\par
  \vspace{1pt}\popsquiggle{poppurple}\par
  \vspace{8pt}
  \begin{tcolorbox}[enhanced,colback=poporange,colframe=popink,boxrule=2.8pt,arc=13pt,
      drop shadow=popink,left=18pt,right=18pt,top=12pt,bottom=13pt,after skip=10pt]
    \poppill{#2 \textbullet\ #3}{popink}{popcream}\hspace{5pt}\poppill{#4}{white}{popink}\par\vspace{8pt}
    {\popdisplay\fontsize{14}{14}\selectfont\color{popink}CHAPTER #5}\par\vspace{4pt}
    {\raggedright\hyphenpenalty=10000\exhyphenpenalty=10000\popdisplay\fontsize{25}{27}\selectfont\color{popcream}\MakeUppercase{#1}\par}
    \vspace{8pt}
    {\popxbold\fontsize{9.5}{9.5}\selectfont\color{popink}\MakeUppercase{Suggested time: #6}}
  \end{tcolorbox}
  \begin{tcolorbox}[enhanced,colback=white,colframe=popink,boxrule=2.2pt,arc=10pt,
      drop shadow=popink,left=16pt,right=16pt,top=10pt,bottom=10pt,after skip=8pt]
    \poppill{In this chapter}{poppurple}{white}\par\vspace{5pt}
    \popsemi\fontsize{9.3}{12.6}\selectfont\color{popink}#7
  \end{tcolorbox}
  \noindent\begin{minipage}[t]{0.487\linewidth}
    \begin{tcolorbox}[enhanced,colback=popgreen,colframe=popink,boxrule=2.2pt,arc=10pt,
        drop shadow=popink,left=11pt,right=11pt,top=10pt,bottom=10pt,height=3.1cm,valign=center]
      \tcbox[colback=white,colframe=popink,boxrule=1.3pt,arc=5pt,boxsep=0pt,nobeforeafter,
        left=3pt,right=3pt,top=3pt,bottom=3pt,valign=center]{\qrcode[height=1.9cm]{https://play.google.com/store/apps/details?id=com.luvvix.onbuch&hl=en}}%
      \hspace{8pt}\begin{minipage}[c]{0.5\linewidth}
        {\popdisplay\fontsize{11}{12.5}\selectfont\color{white}\MakeUppercase{Download OnBuch}}\par\vspace{4pt}
        {\raggedright\popsemi\fontsize{8.4}{11}\selectfont\color{white}Free \textbullet\ Google Play\\AI tutor, past papers, results\par}
      \end{minipage}
    \end{tcolorbox}
  \end{minipage}\hfill
  \begin{minipage}[t]{0.487\linewidth}
    \begin{tcolorbox}[enhanced,colback=poppurple,colframe=popink,boxrule=2.2pt,arc=10pt,
        drop shadow=popink,left=11pt,right=11pt,top=10pt,bottom=10pt,height=3.1cm,valign=center]
      \tikz[baseline=-0.7ex]\node[circle,fill=white,draw=popink,line width=1.3pt,minimum size=1.6cm,inner sep=0]{\popdisplay\fontsize{24}{24}\selectfont\color{poppurple}+};%
      \hspace{8pt}\begin{minipage}[c]{0.6\linewidth}
        {\popdisplay\fontsize{11}{12.5}\selectfont\color{white}\MakeUppercase{Go OnBuch+}}\par\vspace{4pt}
        {\raggedright\popsemi\fontsize{8.4}{11}\selectfont\color{white}All signed courses, unlimited AI tutor, PDF sheets — OnBuch+ tab in the app\par}
      \end{minipage}
    \end{tcolorbox}
  \end{minipage}
  \vspace{8pt}
  \popcontenu
  \vfill
  \signatureonbuch
  \restoregeometry
  \newpage
}

% Rangée « Ce cours contient » (page de garde)
\newcommand{\poptile}[3]{% #1 couleur #2 gros texte #3 légende
  \begin{tcolorbox}[enhanced,colback=#1,colframe=popink,boxrule=2pt,arc=9pt,shadow={2.5pt}{-2.5pt}{0pt}{popink},
      left=6pt,right=6pt,top=9pt,bottom=9pt,halign=center,valign=center,height=1.75cm,width=\linewidth,nobeforeafter]
    {\popdisplay\fontsize{12.5}{13}\selectfont\color{white}\MakeUppercase{#2}}\par\vspace{3pt}
    {\centering\popsemi\fontsize{7.8}{9.5}\selectfont\color{white}#3\par}
  \end{tcolorbox}}
\newcommand{\popcontenu}{%
  \par\noindent{\popxboldls\fontsize{7.5}{7.5}\selectfont\color{popmuted}THIS SIGNED COURSE CONTAINS}\par\vspace{7pt}
  \noindent\begin{minipage}[t]{0.235\linewidth}\poptile{popblue}{Course}{complete, illustrated, step by step}\end{minipage}\hfill
  \begin{minipage}[t]{0.235\linewidth}\poptile{poporange}{Methods}{and worked examples}\end{minipage}\hfill
  \begin{minipage}[t]{0.235\linewidth}\poptile{poppink}{Exercises}{exam style + solutions}\end{minipage}\hfill
  \begin{minipage}[t]{0.235\linewidth}\poptile{popgreen}{Summary sheet}{the essentials to remember}\end{minipage}\par}

% Sommaire
\newtcolorbox{sommaire}{enhanced,colback=white,colframe=popink,boxrule=2.2pt,arc=10pt,
  drop shadow=popink,left=18pt,right=18pt,top=13pt,bottom=13pt,before skip=6pt,after skip=20pt,
  before upper={\poppill{Contents}{poppurple}{white}\par\vspace{9pt}\popsemi\fontsize{10.6}{14.5}\selectfont\color{popink}}}

% Signature de fin de cours — poussée tout en bas de la dernière page
\newcommand{\findecours}{%
  \par\vfill\nopagebreak
  \begin{center}\popsquiggle{poporange}\end{center}\vspace{4pt}
  \signatureonbuch
}
\AtBeginDocument{\def\llbracket{\mathopen{[\mkern-3.5mu[}}\def\rrbracket{\mathclose{]\mkern-3.5mu]}}} % intervalles d'entiers (glyphe ⟦ absent de la police)
% Glyphes absents de Fira Math : redéfinis à partir de glyphes disponibles
\AtBeginDocument{%
  \def\setminus{\mathbin{\backslash}}\let\smallsetminus\setminus
  \def\vdots{\mathord{\vbox{\baselineskip3.2pt\lineskiplimit0pt\kern4pt\hbox{.}\hbox{.}\hbox{.}}}}%
  \def\ddots{\mathinner{\mkern1mu\raise6.5pt\hbox{.}\mkern2mu\raise3.6pt\hbox{.}\mkern2mu\raise0.7pt\hbox{.}\mkern1mu}}%
  \def\triangle{\mathord{\scalebox{0.9}{$\symup{\Delta}$}}}%
  \def\nabla{\mathord{\raisebox{\depth}{\scalebox{1}[-1]{$\symup{\Delta}$}}}}%
  \def\checkmark{\popcheck{popdark}}%
  \def\star{\mathbin{\ast}}\def\bigstar{\mathord{\ast}}%
  \def\diamond{\mathbin{\scalebox{0.6}{$\Diamond$}}}%
  \def\bigcup{\mathop{\mathchoice{\raisebox{-0.3em}{\scalebox{1.7}{$\cup$}}}{\scalebox{1.25}{$\cup$}}{\cup}{\cup}}\displaylimits}%
  \def\bigcap{\mathop{\mathchoice{\raisebox{-0.3em}{\scalebox{1.7}{$\cap$}}}{\scalebox{1.25}{$\cap$}}{\cap}{\cap}}\displaylimits}%
  \def\lesseqqgtr{\mathrel{\lessgtr}}\def\bigoplus{\mathop{\oplus}\displaylimits}%
}
\providecommand{\ssi}{if and only if} % abréviation fréquente
\AtBeginDocument{\providecommand{\wideparen}{\overparen}} % arc de cercle

% Fonctions usuelles que les modèles utilisent sans que amsmath les fournisse
\providecommand{\arsinh}{\operatorname{arsinh}}\providecommand{\arcosh}{\operatorname{arcosh}}
\providecommand{\artanh}{\operatorname{artanh}}\providecommand{\arsech}{\operatorname{arsech}}
\providecommand{\arcsch}{\operatorname{arcsch}}\providecommand{\arcoth}{\operatorname{arcoth}}
\providecommand{\arcsinh}{\operatorname{arsinh}}\providecommand{\arccosh}{\operatorname{arcosh}}
\providecommand{\arctanh}{\operatorname{artanh}}\providecommand{\sech}{\operatorname{sech}}
\providecommand{\csch}{\operatorname{csch}}\providecommand{\arcsec}{\operatorname{arcsec}}
\providecommand{\arccsc}{\operatorname{arccsc}}\providecommand{\arccot}{\operatorname{arccot}}
\providecommand{\sgn}{\operatorname{sgn}}\providecommand{\lcm}{\operatorname{lcm}}
\providecommand{\Var}{\operatorname{Var}}\providecommand{\Cov}{\operatorname{Cov}}

%% FIGFIX-BEGIN v3 — correctifs globaux des figures (docs/onbuch-generation/tools/figfix)
\usepackage{adjustbox}\usepackage{etoolbox}
% toute figure TikZ/pgfplots plus large que la ligne est réduite à la largeur disponible
\BeforeBeginEnvironment{tikzpicture}{\begin{adjustbox}{max width=\linewidth}}
\AfterEndEnvironment{tikzpicture}{\end{adjustbox}}
% légendes pgfplots lisibles par-dessus les courbes
\pgfplotsset{every axis/.append style={legend style={fill=white,fill opacity=0.92,text opacity=1,draw=black!15,inner sep=3pt}}}
% étiquettes d'arêtes (syntaxe edge["..."]) avec fond blanc
\tikzset{every edge quotes/.append style={fill=white,inner sep=1.2pt}}
% bulles de cartes mentales : texte à la ligne dans la bulle, bulle un peu plus grande
\tikzset{every concept/.append style={text width=2.0cm,align=center,minimum size=2.3cm,inner sep=2pt}}
%% FIGFIX-END
\begin{document}

\pagegarde{Elementary Probability}{Pure Mathematics with Statistics}{Upper Sixth}{Upper Sixth — Pure mathematics with statistics}{10}{15 periods}{This chapter introduces the language and laws of probability: experiments, sample spaces, events, classical and empirical probability, the addition and multiplication laws, conditional probability, tree diagrams, counting techniques and Bayes' theorem. It provides the essential foundation for the statistics and probability questions of the GCE Advanced Level Paper in Pure Mathematics with Statistics.
\vspace{4pt}
\begin{multicols}{2}\small
\begin{itemize}
\item By the end of this chapter I can define probability and use correctly the basic vocabulary: experiment, outcome, sample space, event.
\item By the end of this chapter I can list the sample space of simple and compound experiments and identify types of events (simple, compound, complementary, mutually exclusive, exhaustive, independent).
\item By the end of this chapter I can compute classical (theoretical) probabilities using equally likely outcomes.
\item By the end of this chapter I can estimate probabilities empirically from relative frequencies and relate them to theoretical probability.
\item By the end of this chapter I can apply the addition law for mutually exclusive and for non-mutually exclusive events.
\item By the end of this chapter I can apply the multiplication law and compute conditional probabilities, testing events for independence.
\end{itemize}\end{multicols}}

\begin{sommaire}
\begin{multicols}{2}
\begin{enumerate}
\item Probability, Experiments and Sample Spaces
\item Types of Events and Classical Probability
\item Empirical Probability and the Addition Laws
\item Multiplication Law, Conditional Probability and Independence
\item Tree Diagrams and Selection With or Without Replacement
\item Counting Techniques in Probability
\item Total Probability and Bayes' Theorem
\item Integration activity
\item Methods and worked examples
\item Exercises
\item Solutions to the exercises
\item Summary sheet
\end{enumerate}
\end{multicols}
\end{sommaire}

\begin{objectifs}
\begin{itemize}
\item By the end of this chapter I can define probability and use correctly the basic vocabulary: experiment, outcome, sample space, event.
\item By the end of this chapter I can list the sample space of simple and compound experiments and identify types of events (simple, compound, complementary, mutually exclusive, exhaustive, independent).
\item By the end of this chapter I can compute classical (theoretical) probabilities using equally likely outcomes.
\item By the end of this chapter I can estimate probabilities empirically from relative frequencies and relate them to theoretical probability.
\item By the end of this chapter I can apply the addition law for mutually exclusive and for non-mutually exclusive events.
\item By the end of this chapter I can apply the multiplication law and compute conditional probabilities, testing events for independence.
\item By the end of this chapter I can construct and use tree diagrams, including selections with and without replacement.
\item By the end of this chapter I can use permutations and combinations to count outcomes in probability problems.
\item By the end of this chapter I can state and apply Bayes' theorem to reverse conditional probabilities in real-life contexts.
\end{itemize}
\end{objectifs}

\begin{prerequis}
\begin{itemize}
\item Set language and operations: union, intersection, complement, Venn diagrams (Lower Sixth).
\item Fractions, decimals, percentages and their simplification.
\item Factorial notation, permutations (nPr) and combinations (nCr) from the chapter on counting techniques.
\item Basic algebraic manipulation: solving simple linear equations and working with fractions of fractions.
\item Interpreting simple statistical tables and relative frequencies (Lower Sixth statistics).
\end{itemize}
\end{prerequis}

\newpage

\coursec{Probability, Experiments and Sample Spaces}

Every rainy season in Douala, commuters at Ndokoti junction gamble on whether the sky will open before they reach Bonabéri; every Saturday, punters at a Yaoundé betting kiosk stake money on football outcomes they cannot control. Behind each of these everyday decisions lies the same question: \emph{how likely is an uncertain event?} Probability gives a precise numerical answer between $0$ and $1$, and it powers insurance premiums, medical testing, election forecasts and the National Lottery. In this chapter you will learn to measure chance rigorously, from tossing a coin to applying Bayes' theorem to a medical diagnosis.

To measure chance, we must first agree on the language we use to describe uncertain situations. That is the goal of this first section: we define what a \cle{random experiment} is, we list all its possible results in a \cle{sample space}, and we learn to describe \cle{events} as sets of results. Everything else in the chapter is built on these foundations.

\courssub{Deterministic and Random Experiments}

In science and in daily life, we perform \emph{experiments}: actions whose results we observe. Some experiments are completely predictable.

\begin{experience}[Two kinds of experiments]
\textbf{Experiment 1.} Heat pure water at sea level. It boils at $100\,^{\circ}\mathrm{C}$, every single time. The result is fixed by the laws of physics.

\textbf{Experiment 2.} Toss a 25 FCFA coin and note which face is up. Sometimes it is Head, sometimes Tail. No matter how carefully you repeat the toss under identical conditions, you cannot predict the result of the next toss.
\end{experience}

Experiment 1 is \cle{deterministic}: the same conditions always produce the same result. Experiment 2 is \cle{random}: the result cannot be predicted in advance, even though we know all the possible results.

\begin{definition}[Random experiment, trial, outcome]
A \cle{random experiment} is an action or process whose result cannot be predicted with certainty before it is performed, but whose set of all possible results is known.

Each performance of the experiment is called a \cle{trial}, and each possible result of a trial is called an \cle{outcome}.
\end{definition}

\begin{exemplebox}[Classifying experiments]
Classify each experiment as deterministic or random.
\begin{enumerate}
\item Dropping a stone from the top of a building in Buea and measuring whether it falls. \emph{Deterministic} — it always falls.
\item Throwing a fair die and recording the number on top. \emph{Random} — any of $1, 2, 3, 4, 5, 6$ may appear.
\item Drawing one card from a well-shuffled pack of $52$ cards. \emph{Random}.
\item Counting the number of days in the month of February 2024. \emph{Deterministic} — the answer, $29$, is fixed.
\end{enumerate}
\end{exemplebox}

\courssub{The Sample Space}

Before we can measure how likely something is, we must know \emph{everything that can happen}. This complete list is called the sample space.

\begin{definition}[Sample space]
The \cle{sample space} of a random experiment, denoted $S$ (or $\Omega$), is the set of \textbf{all possible outcomes} of the experiment.
\end{definition}

\begin{exemplebox}[Sample spaces of one-stage experiments]
\begin{enumerate}
\item \textbf{Tossing one coin:} $S = \{H, T\}$, where $H$ = Head and $T$ = Tail.
\item \textbf{Throwing one die:} $S = \{1, 2, 3, 4, 5, 6\}$.
\item \textbf{Drawing one card and noting its suit:} $S = \{\text{Hearts}, \text{Diamonds}, \text{Clubs}, \text{Spades}\}$.
\item \textbf{Spinning a spinner with four coloured sectors:} $S = \{\text{Red}, \text{Blue}, \text{Green}, \text{Yellow}\}$.
\end{enumerate}
\end{exemplebox}

\courssub{Sample Spaces of Two-Stage Experiments}

Many experiments have \emph{two stages}: two coins are tossed, two dice are thrown, or a coin is tossed and then a die is thrown. Listing the sample space carelessly almost always leads to missing or duplicated outcomes. We therefore use \textbf{systematic enumeration}: ordered pairs, lists, or two-way tables (grids).

\begin{methode}[Systematic listing of a sample space]
\begin{enumerate}
\item Identify the stages of the experiment and the possible results of each stage.
\item Represent each outcome as an ordered pair $(\text{stage 1 result}, \text{stage 2 result})$.
\item For two-stage experiments, draw a two-way table: rows for the first stage, columns for the second stage; each cell is one outcome.
\item Count the outcomes: if stage 1 has $m$ results and stage 2 has $n$ results, the sample space has $m \times n$ outcomes.
\end{enumerate}
\end{methode}

\begin{exemplebox}[Two coins]
Tossing two coins: stage 1 gives $H$ or $T$, stage 2 gives $H$ or $T$. The sample space is
\[ S = \{(H,H),\ (H,T),\ (T,H),\ (T,T)\}, \]
which has $2 \times 2 = 4$ outcomes. Note that $(H,T)$ and $(T,H)$ are \emph{different} outcomes: in the first, coin 1 shows Head; in the second, coin 2 shows Head.
\end{exemplebox}

\begin{exemplebox}[Two dice]
Throwing two dice (one red, one blue): each die has $6$ faces, so the sample space has $6 \times 6 = 36$ outcomes, each written $(\text{red}, \text{blue})$. The grid below displays all $36$ outcomes and highlights the event ``the sum of the two numbers is $7$''.
\end{exemplebox}

\begin{popfigure}
\begin{center}
\begin{tikzpicture}[scale=0.62, every node/.style={font=\small}]
% grid cells
\foreach \i in {1,...,6}{
  \foreach \j in {1,...,6}{
    \pgfmathtruncatemacro{\s}{\i+\j}
    \ifnum\s=7
      \fill[poporange!45] (\j-1,6-\i) rectangle (\j,7-\i);
    \fi
  }
}
\draw[popline, step=1] (0,0) grid (6,6);
% outcomes
\foreach \i in {1,...,6}{
  \foreach \j in {1,...,6}{
    \node at (\j-0.5,6.5-\i) {(\i,\j)};
  }
}
% axis labels
\foreach \i in {1,...,6}{
  \node[left] at (0,6.5-\i) {\i};
  \node[below] at (\i-0.5,0) {\i};
}
\node[left, align=center] at (-0.9,3) {\textbf{Die 1}};
\node[below] at (3,-0.9) {\textbf{Die 2}};
% legend
\fill[poporange!45] (7.2,4.6) rectangle (7.8,5.2);
\node[right, align=left] at (8.0,4.9) {sum $=7$};
\end{tikzpicture}
\end{center}
\legende{The 36 outcomes of throwing two dice. The shaded diagonal cells are the six outcomes whose coordinates sum to 7.}
\end{popfigure}

\begin{exemplebox}[A coin, then a die]
Toss a coin, then throw a die. The sample space is
\[ S = \{(H,1),(H,2),(H,3),(H,4),(H,5),(H,6),(T,1),(T,2),(T,3),(T,4),(T,5),(T,6)\}, \]
with $2 \times 6 = 12$ outcomes.
\end{exemplebox}

\courssub{Events and Types of Events}

Usually we are not interested in one single outcome but in a \emph{collection} of outcomes: ``an even number on the die'', ``at least one Head'', ``a sum of 7''. Such collections are called events.

\begin{definition}[Event, simple event, compound event]
An \cle{event} is any subset of the sample space $S$. We say an event $A$ \emph{occurs} when the outcome of the trial belongs to $A$.
\begin{itemize}
\item A \cle{simple} (or \cle{elementary}) \cle{event} contains exactly one outcome.
\item A \cle{compound event} contains two or more outcomes.
\item The whole sample space $S$ is the \cle{certain event}: it always occurs.
\item The empty set $\varnothing$ is the \cle{impossible event}: it never occurs.
\end{itemize}
\end{definition}

\begin{exemplebox}[Events when throwing one die]
With $S = \{1,2,3,4,5,6\}$:
\begin{itemize}
\item $A = \{5\}$ is a simple event (``score 5'').
\item $B = \{2,4,6\}$ is a compound event (``even score'').
\item $C = \{1,2,3,4,5,6\} = S$ is the certain event (``score less than 7'').
\item $D = \varnothing$ is the impossible event (``score 8'').
\end{itemize}
\end{exemplebox}

\begin{attention}[An outcome is not an event]
Do not confuse an \emph{outcome} with an \emph{event}. The outcome is a single result, e.g.\ $3$. An event is a \textbf{set} of outcomes, e.g.\ $\{3\}$ or $\{1,3,5\}$. Writing ``the event $3$'' is incorrect; write ``the event $\{3\}$''. This distinction becomes essential when we compute probabilities: we always compute the probability of a \emph{set}.
\end{attention}

\courssub{Types of Events: Mutually Exclusive and Exhaustive Events}

Beyond simple and compound, events are classified by how they relate to each other.

\begin{definition}[Mutually exclusive events]
Two events $A$ and $B$ are \cle{mutually exclusive} (or \cle{disjoint}) if they cannot occur simultaneously in the same trial. In set notation:
\[ A \cap B = \varnothing. \]
\end{definition}

\begin{definition}[Exhaustive events]
A set of events $A_1, A_2, \dots, A_n$ is \cle{exhaustive} if at least one of them must occur in every trial. In set notation:
\[ A_1 \cup A_2 \cup \dots \cup A_n = S. \]
If the events are both mutually exclusive and exhaustive, they form a \cle{partition} of the sample space.
\end{definition}

\begin{exemplebox}[Mutually exclusive and exhaustive events]
Throwing one die: $S = \{1,2,3,4,5,6\}$.
\begin{itemize}
\item $A = \{1,2,3\}$ (score $\le 3$) and $B = \{4,5,6\}$ (score $\ge 4$) are mutually exclusive ($A \cap B = \varnothing$) and exhaustive ($A \cup B = S$). They form a partition.
\item $C = \{2,4,6\}$ (even) and $D = \{1,3,5\}$ (odd) are also mutually exclusive and exhaustive.
\item $E = \{1,2\}$ and $F = \{2,3\}$ are \textbf{not} mutually exclusive ($E \cap F = \{2\} \ne \varnothing$).
\item $G = \{1,2\}$ and $H = \{4,5\}$ are mutually exclusive but \textbf{not} exhaustive ($G \cup H = \{1,2,4,5\} \ne S$).
\end{itemize}
\end{exemplebox}

\begin{propriete}[Probability of mutually exclusive events (preview)]
If $A$ and $B$ are mutually exclusive, then $P(A \cap B) = 0$. The full Addition Law is developed in Section 3.
\end{propriete}

\courssub{The Complement of an Event}

\begin{definition}[Complement]
The \cle{complement} of an event $A$, written $A'$ (or $\bar{A}$), is the event that $A$ does \textbf{not} occur. It contains every outcome of $S$ that is not in $A$:
\[ A' = \{\text{outcomes in } S \text{ but not in } A\}. \]
\end{definition}

For example, throwing one die with $A = \{2,4,6\}$ (even score) gives $A' = \{1,3,5\}$ (odd score). Since every trial produces exactly one outcome, either $A$ occurs or $A'$ occurs, never both and never neither. Consequently, $A$ and $A'$ are always mutually exclusive and exhaustive: $A \cap A' = \varnothing$ and $A \cup A' = S$.

\begin{popfigure}
\begin{center}
\begin{tikzpicture}[scale=1.0]
% sample space rectangle
\fill[popblueL] (0,0) rectangle (7,4.4);
\draw[popdark, thick] (0,0) rectangle (7,4.4);
\node[popdark, font=\bfseries] at (0.55,3.95) {$S$};
% event A
\fill[poporange!55] (2.5,2.2) circle (1.35);
\draw[poporange, very thick] (2.5,2.2) circle (1.35);
\node[popdark, font=\bfseries] at (2.5,2.2) {$A$};
% complement label
\node[popdark, font=\bfseries] at (5.5,2.2) {$A'$};
\node[popmuted, align=center, font=\small] at (5.5,1.4) {shaded region:\\ $A$ does not occur};
\end{tikzpicture}
\end{center}
\legende{The sample space $S$ (rectangle) containing an event $A$ (circle). The region outside the circle but inside the rectangle is the complement $A'$.}
\end{popfigure}

\courssub{The Probability Scale (Intuitive Introduction)}

We now attach a number to each event. Intuitively, probability measures \emph{how likely} an event is, on a scale from $0$ to $1$:
\begin{itemize}
\item $P = 0$: the event is \textbf{impossible} (it can never happen);
\item $P = 1$: the event is \textbf{certain} (it must happen);
\item values in between measure increasing likelihood: $P = 0.5$ means ``as likely as not''.
\end{itemize}

\begin{propriete}[Fundamental properties of probability]
For any event $A$ in a sample space $S$:
\[ 0 \le P(A) \le 1, \qquad P(S) = 1, \qquad P(\varnothing) = 0, \qquad P(A) + P(A') = 1. \]
(These properties follow directly from the definition of probability as a proportion of outcomes; they are not examinable as proofs but must be known and used fluently.)
\end{propriete}

\begin{attention}[{Probabilities outside $[0,1]$}]
A probability can \textbf{never} be negative or greater than $1$. If your calculation gives $P(A) = 1.4$ or $P(A) = -0.2$, you have made an error — go back and check. Also, $P(A) = 0$ means the event is impossible; it does \emph{not} mean ``very unlikely''.
\end{attention}

\courssub{Equally Likely Outcomes and Fairness}

To assign probabilities we need one more ingredient: the idea that some outcomes are \emph{equally likely}.

\begin{definition}[Equally likely outcomes]
Outcomes of a random experiment are \cle{equally likely} when no outcome is more likely to occur than any other. This is the case for a \textbf{fair} coin, an \textbf{unbiased} die, or a well-shuffled pack of cards.
\end{definition}

When a fair coin is tossed, $P(H) = P(T) = 0.5$. When an unbiased die is thrown, each face has probability $\tfrac{1}{6}$. The assumption of fairness is a \emph{modelling choice}: a real coin may be slightly imperfect, but the fair model is an excellent approximation — and the GCE will always tell you when a coin or die is \emph{biased}.

\begin{exoresolu}[Sample space, events and complements]
Two fair coins are tossed.
\begin{enumerate}
\item Write down the sample space $S$.
\item Write as sets the events: $A$ = ``exactly one Head'', $B$ = ``at least one Head'', $C$ = ``both coins show the same face''.
\item Write down $B'$ and verify that $B \cap B' = \varnothing$ and $B \cup B' = S$.
\item Which of $A$, $B$, $C$ are simple events and which are compound?
\item Are $A$ and $C$ mutually exclusive? Are they exhaustive?
\end{enumerate}
\tcblower
\textbf{Solution.}
\begin{enumerate}
\item $S = \{(H,H), (H,T), (T,H), (T,T)\}$.
\item $A = \{(H,T), (T,H)\}$;\quad $B = \{(H,H), (H,T), (T,H)\}$;\quad $C = \{(H,H), (T,T)\}$.
\item $B'$ = ``no Head'' $= \{(T,T)\}$. Then $B \cap B' = \varnothing$ (no outcome is in both) and $B \cup B' = \{(H,H),(H,T),(T,H),(T,T)\} = S$, as required for complementary events.
\item $A$, $B$ and $C$ each contain more than one outcome, so all three are \textbf{compound events}. The simple events here are $\{(H,H)\}$, $\{(H,T)\}$, $\{(T,H)\}$ and $\{(T,T)\}$.
\item $A \cap C = \varnothing$ (no outcome has exactly one Head and both coins the same), so $A$ and $C$ are \textbf{mutually exclusive}. $A \cup C = \{(H,T),(T,H),(H,H),(T,T)\} = S$, so they are also \textbf{exhaustive}. Hence $A$ and $C$ form a partition of $S$.
\end{enumerate}
\end{exoresolu}

\begin{exercice}[Building sample spaces and identifying events]
\begin{enumerate}
\item A fair die is thrown and a fair coin is tossed. List the sample space and state the number of outcomes.
\item A bag contains 3 balls: one red (R), one green (G), one blue (B). Two balls are drawn, one after the other, the first being replaced before the second draw. List the sample space.
\item For the experiment in (2), write as sets the events $E$ = ``both balls have the same colour'' and $F$ = ``the red ball is drawn at least once''. Then write $E'$ and $F'$.
\item In the experiment of (1), let $G$ = ``the die shows an even number'' and $H$ = ``the coin shows Tail''. Are $G$ and $H$ mutually exclusive? Are they exhaustive?
\end{enumerate}
\end{exercice}

\begin{corrige}[Exercise 1]
\begin{enumerate}
\item $S = \{(1,H),(2,H),(3,H),(4,H),(5,H),(6,H),(1,T),(2,T),(3,T),(4,T),(5,T),(6,T)\}$; there are $6 \times 2 = 12$ outcomes.
\item $S = \{(R,R),(R,G),(R,B),(G,R),(G,G),(G,B),(B,R),(B,G),(B,B)\}$; $3 \times 3 = 9$ outcomes.
\item $E = \{(R,R),(G,G),(B,B)\}$; $F = \{(R,R),(R,G),(R,B),(G,R),(B,R)\}$; $E' = \{(R,G),(R,B),(G,R),(G,B),(B,R),(B,G)\}$ (the two balls have different colours); $F' = \{(G,G),(G,B),(B,G),(B,B)\}$ (the red ball is not drawn).
\item $G = \{(2,H),(4,H),(6,H),(2,T),(4,T),(6,T)\}$; $H = \{(1,T),(2,T),(3,T),(4,T),(5,T),(6,T)\}$. $G \cap H = \{(2,T),(4,T),(6,T)\} \ne \varnothing$, so \textbf{not mutually exclusive}. $G \cup H = \{(1,T),(2,H),(2,T),(3,T),(4,H),(4,T),(5,T),(6,H),(6,T)\} \ne S$ (missing $(1,H),(3,H),(5,H)$), so \textbf{not exhaustive}.
\end{enumerate}
\end{corrige}

\begin{aretenir}[Key points of this section]
\begin{itemize}
\item A \cle{random experiment} has an unpredictable result; a \cle{trial} is one performance; an \cle{outcome} is one possible result.
\item The \cle{sample space} $S$ is the set of all outcomes. For two-stage experiments, use ordered pairs and two-way tables: $m$ results then $n$ results give $m \times n$ outcomes.
\item An \cle{event} is a subset of $S$: simple (one outcome) or compound (several). $S$ is certain, $\varnothing$ is impossible.
\item \cle{Mutually exclusive} events cannot occur together ($A \cap B = \varnothing$). \cle{Exhaustive} events cover all possibilities ($A \cup B = S$). A partition is both.
\item The \cle{complement} $A'$ is ``$A$ does not occur'', with $A \cap A' = \varnothing$, $A \cup A' = S$, and later $P(A) + P(A') = 1$.
\item For every event: $0 \le P(A) \le 1$, $P(S) = 1$, $P(\varnothing) = 0$.
\item \cle{Equally likely} outcomes (fair coin, unbiased die) are the basis of classical probability, developed in the next section.
\end{itemize}
\end{aretenir}

\coursec{Types of Events and Classical Probability}

In the previous section we introduced the language of probability: experiments, outcomes, sample spaces and events. We now move to the heart of the matter — \emph{how to assign a numerical value to the likelihood of an event}. The oldest and most intuitive approach is \textbf{classical (theoretical) probability}, which applies when all outcomes of an experiment are \cle{equally likely}. We shall also classify events according to how they relate to each other (mutually exclusive, exhaustive, complementary) and see how these relationships simplify calculations.

\courssub{Classical Probability — Equally Likely Outcomes}

\begin{experience}[Observing symmetry in games of chance]
Take a fair six-sided die. There is no physical reason why the face $3$ should appear more often than the face $5$: the cube is symmetric, the mass is evenly distributed, and each face has the same shape and size. Common sense therefore tells us that each of the six faces has the same chance, namely $\frac{1}{6}$. The same reasoning applies to a fair coin (two symmetric faces) or to a well-shuffled pack of cards (no card is favoured). This \emph{symmetry principle} is the foundation of classical probability.
\end{experience}

\begin{definition}[Classical Probability]
If a random experiment has a finite sample space $S$ in which all $n(S)$ outcomes are \textbf{equally likely}, and $A$ is an event consisting of $n(A)$ outcomes, then the probability of $A$ is
\[
P(A) = \frac{n(A)}{n(S)} = \frac{\text{number of favourable outcomes}}{\text{total number of possible outcomes}}.
\]
\end{definition}

\begin{propriete}[Elementary properties of probability]
For any experiment with finite sample space $S$ and equally likely outcomes:
\begin{itemize}
\item $0 \le P(A) \le 1$ for every event $A$;
\item $P(S) = 1$ (the certain event) and $P(\varnothing) = 0$ (the impossible event);
\item $P(A) + P(A') = 1$, where $A'$ is the complement of $A$.
\end{itemize}
These follow immediately from the definition, since $0 \le n(A) \le n(S)$, $n(S)/n(S)=1$, $n(\varnothing)=0$ and $n(A)+n(A')=n(S)$.
\end{propriete}

\begin{attention}[Equally likely outcomes only!]
The formula $P(A)=n(A)/n(S)$ is \textbf{valid only when every outcome in $S$ has the same chance of occurring}. Do not apply it to biased coins, loaded dice, or situations where outcomes have different probabilities (e.g. the weather tomorrow). In such cases we must use empirical probability (next section) or another model. A frequent GCE error is to count outcomes that are \emph{not} equally likely: for example, when two coins are tossed, the three events ``two heads'', ``one head and one tail'', ``two tails'' are \textbf{not} equally likely — the correct sample space is $\{HH, HT, TH, TT\}$ with four equally likely outcomes.
\end{attention}

\begin{exemplebox}[Rolling a fair die]
A fair six-sided die is rolled once. The sample space is $S=\{1,2,3,4,5,6\}$ with $n(S)=6$ equally likely outcomes.
\begin{itemize}
\item Event $A$: ``an even number appears'' $\Rightarrow A=\{2,4,6\},\ n(A)=3$. $P(A)=\frac{3}{6}=\frac{1}{2}$.
\item Event $B$: ``a number greater than 4'' $\Rightarrow B=\{5,6\},\ n(B)=2$. $P(B)=\frac{2}{6}=\frac{1}{3}$.
\item Event $C$: ``a 7 appears'' $\Rightarrow C=\varnothing,\ n(C)=0$. $P(C)=0$ (impossible event).
\item Event $D$: ``a number less than 7'' $\Rightarrow D=S,\ n(D)=6$. $P(D)=1$ (certain event).
\end{itemize}
\end{exemplebox}

\begin{popfigure}
\begin{tikzpicture}[scale=0.9]
% 52-card grid: 4 suits (rows) x 13 ranks (columns)
\foreach \s/\suit/\col in {0/S/black, 1/H/poporange, 2/D/poporange, 3/C/black}{
  \foreach \r/\rank in {0/A, 1/2, 2/3, 3/4, 4/5, 5/6, 6/7, 7/8, 8/9, 9/10, 10/J, 11/Q, 12/K}{
    \pgfmathsetmacro{\x}{\r*1.1}
    \pgfmathsetmacro{\y}{-\s*1.1}
    \draw[thin, fill=white] (\x,\y) rectangle (\x+1,\y+1);
    \node[font=\tiny, text=\col] at (\x+0.5,\y+0.7) {\rank};
    \node[font=\small\bfseries, text=\col] at (\x+0.5,\y+0.3) {\suit};
  }
}
% Suit labels on left
\foreach \s/\suit/\col in {0/S/black, 1/H/poporange, 2/D/poporange, 3/C/black}{
  \pgfmathsetmacro{\y}{-\s*1.1+0.5}
  \node[anchor=east, font=\small\bfseries, text=\col] at (-0.3,\y) {\suit};
}
% Rank labels on top
\foreach \r/\rank in {0/A, 1/2, 2/3, 3/4, 4/5, 5/6, 6/7, 7/8, 8/9, 9/10, 10/J, 11/Q, 12/K}{
  \pgfmathsetmacro{\x}{\r*1.1+0.5}
  \node[anchor=south, font=\tiny\bfseries] at (\x,0.3) {\rank};
}
% Title
\node[font=\bfseries] at (6.6,1.2) {Standard 52-Card Pack};
\end{tikzpicture}
\legende{Structure of a standard 52-card deck: 4 suits $\times$ 13 ranks. S = Spades, C = Clubs (black cards); H = Hearts, D = Diamonds (red cards). A = Ace, J = Jack, Q = Queen, K = King.}
\end{popfigure}

\begin{exemplebox}[Drawing a card from a well-shuffled pack]
A standard pack has $52$ cards. One card is drawn at random.
\begin{itemize}
\item $P(\text{Ace}) = \frac{4}{52} = \frac{1}{13}$.
\item $P(\text{Heart}) = \frac{13}{52} = \frac{1}{4}$.
\item $P(\text{Ace of Hearts}) = \frac{1}{52}$.
\item $P(\text{Face card: J, Q or K}) = \frac{12}{52} = \frac{3}{13}$.
\item $P(\text{Red card}) = \frac{26}{52} = \frac{1}{2}$.
\end{itemize}
\end{exemplebox}

\begin{exemplebox}[Balls in a bag]
A bag contains $5$ red, $3$ blue and $2$ green balls (total $10$). One ball is drawn at random.
\[
P(\text{Red}) = \frac{5}{10}=\frac{1}{2},\quad
P(\text{Blue}) = \frac{3}{10},\quad
P(\text{Green}) = \frac{2}{10}=\frac{1}{5},\quad
P(\text{Not red}) = \frac{5}{10}=\frac{1}{2}.
\]
Note that $P(\text{Red})+P(\text{Blue})+P(\text{Green}) = \frac{1}{2}+\frac{3}{10}+\frac{1}{5} = 1$: the three events are mutually exclusive and exhaustive.
\end{exemplebox}

\begin{exemplebox}[Lottery-type draw]
In a lottery, $6$ numbers are drawn from $\{1,2,\dots,49\}$ without replacement, and the order does not matter. The total number of equally likely combinations is $\binom{49}{6}=13\,983\,816$. If you buy one ticket with a specific combination, your chance of winning the jackpot is
\[
P(\text{jackpot}) = \frac{1}{\binom{49}{6}} = \frac{1}{13\,983\,816} \approx 7.15\times 10^{-8}.
\]
This tiny number explains why jackpot winners are so rare!
\end{exemplebox}

\courssub{Relationships Between Events}

Events are subsets of the sample space $S$. We classify them by how they overlap.

\begin{definition}[Mutually Exclusive (Disjoint) Events]
Two events $A$ and $B$ are \textbf{mutually exclusive} if they cannot occur simultaneously, i.e. $A \cap B = \varnothing$. On a Venn diagram their regions do not overlap. For example, when a die is rolled once, ``obtaining a 2'' and ``obtaining a 5'' are mutually exclusive.
\end{definition}

\begin{definition}[Exhaustive Events]
A collection of events $A_1, A_2, \dots, A_k$ is \textbf{exhaustive} if their union is the whole sample space: $A_1 \cup A_2 \cup \dots \cup A_k = S$. In other words, at least one of them must occur. If they are also pairwise mutually exclusive, they form a \textbf{partition} of $S$, and then
\[
P(A_1) + P(A_2) + \dots + P(A_k) = 1.
\]
\end{definition}

\begin{definition}[Complementary Event]
For any event $A$, its \textbf{complement} $A'$ (also written $\overline{A}$ or $A^c$) is the event ``$A$ does not occur''. $A$ and $A'$ are mutually exclusive and exhaustive: $A \cap A' = \varnothing$ and $A \cup A' = S$. Hence $P(A) + P(A') = 1$, so
\[
P(A') = 1 - P(A).
\]
\end{definition}

\begin{definition}[Independent Events (informal)]
Two events $A$ and $B$ are \textbf{independent} if the occurrence of one does not affect the probability of the other. The formal definition, $P(A \cap B) = P(A)P(B)$, will be given in a later section. For now, think of ``drawing a card, replacing it, shuffling, and drawing again'' — the two draws are independent.
\end{definition}

\begin{attention}[Mutually exclusive $\neq$ Independent]
This is a classic trap!
\begin{itemize}
\item \textbf{Mutually exclusive}: $A$ and $B$ cannot happen together $\Rightarrow P(A \cap B)=0$. If $P(A)>0$ and $P(B)>0$, they are in fact \textbf{dependent}: knowing that $A$ occurred tells you that $B$ certainly did not.
\item \textbf{Independent}: $P(A \cap B)=P(A)P(B)$. Two independent events with non-zero probabilities \emph{can} happen together.
\end{itemize}
Never write $P(A\cap B)=P(A)P(B)$ for mutually exclusive events, and never write $P(A\cap B)=0$ for independent events.
\end{attention}

\begin{popfigure}
\begin{tikzpicture}[scale=1.1]
% (a) Mutually exclusive
\begin{scope}
  \draw[thick, popblue] (0,0) circle (0.9);
  \draw[thick, poporange] (2.4,0) circle (0.9);
  \node at (0,0) {$A$};
  \node at (2.4,0) {$B$};
  \draw[thin] (-1.4,-1.3) rectangle (3.8,1.3);
  \node[anchor=north west] at (-1.4,1.3) {$S$};
  \node[align=center, font=\small] at (1.2,-1.8) {(a) Mutually exclusive:\\ $A\cap B=\varnothing$};
\end{scope}
% (b) Overlapping
\begin{scope}[xshift=7.5cm]
  \draw[thick, popblue] (0,0) circle (0.9);
  \draw[thick, poporange] (1.3,0) circle (0.9);
  \node at (-0.5,0) {$A$};
  \node at (1.8,0) {$B$};
  \node[font=\small] at (0.65,0) {$A\cap B$};
  \draw[thin] (-1.4,-1.3) rectangle (2.7,1.3);
  \node[anchor=north west] at (-1.4,1.3) {$S$};
  \node[align=center, font=\small] at (0.65,-1.8) {(b) Not mutually exclusive:\\ $A\cap B\neq\varnothing$};
\end{scope}
% (c) Exhaustive partition into three
\begin{scope}[yshift=-5.2cm, xshift=2cm]
  \draw[thin] (0,0) rectangle (6,2);
  \node[anchor=north west] at (0,2) {$S$};
  \draw[thick, popblue] (2,0) -- (2,2);
  \draw[thick, popblue] (4,0) -- (4,2);
  \node at (1,1) {$A_1$};
  \node at (3,1) {$A_2$};
  \node at (5,1) {$A_3$};
  \node[align=center, font=\small] at (3,-0.6) {(c) Exhaustive partition: $A_1\cup A_2\cup A_3=S$, pairwise disjoint};
\end{scope}
\end{tikzpicture}
\legende{Venn diagrams: (a) mutually exclusive events; (b) overlapping events; (c) three events forming an exhaustive partition of $S$.}
\end{popfigure}

\courssub{Using the Complement — ``At Least One'' Strategy}

\begin{methode}[Complement Rule for ``At Least One'']
When a problem asks for the probability of ``at least one'' success in several trials, it is often much easier to compute the probability of \textbf{none} (the complement) and subtract from $1$.
\begin{enumerate}
\item Identify the event ``no success'' (all trials fail).
\item Compute $P(\text{no success})$ — usually a product if the trials are independent.
\item Then $P(\text{at least one success}) = 1 - P(\text{no success})$.
\end{enumerate}
\end{methode}

\begin{exemplebox}[At least one six in two rolls of a fair die]
Roll a fair die twice. Find $P(\text{at least one 6})$.
\begin{align*}
P(\text{no 6 in one roll}) &= \frac{5}{6}. \\
P(\text{no 6 in two rolls}) &= \left(\frac{5}{6}\right)^2 = \frac{25}{36}. \\
P(\text{at least one 6}) &= 1 - \frac{25}{36} = \frac{11}{36}.
\end{align*}
\end{exemplebox}

\begin{exemplebox}[At least one ace in a 5-card hand]
A hand of $5$ cards is dealt from a standard $52$-card pack. Find $P(\text{at least one ace})$.
\[
P(\text{no ace}) = \frac{\binom{48}{5}}{\binom{52}{5}} = \frac{1\,712\,304}{2\,598\,960} \approx 0.6588.
\]
\[
P(\text{at least one ace}) = 1 - 0.6588 = 0.3412.
\]
Direct counting would require four separate cases (exactly 1, 2, 3 or 4 aces); the complement gives the answer in one line.
\end{exemplebox}

\courssub{Odds in Favour and Odds Against (GCE Speed Tip)}

\begin{definition}[Odds]
If $P(A)=p$ (with $0<p<1$), then
\begin{itemize}
\item \textbf{Odds in favour of $A$} $= p : (1-p)$, usually simplified to a ratio $a:b$ of integers;
\item \textbf{Odds against $A$} $= (1-p) : p$.
\end{itemize}
Conversely, if the odds in favour of $A$ are $a:b$, then $P(A)=\frac{a}{a+b}$.
\end{definition}

\begin{exemplebox}[Odds examples]
\begin{itemize}
\item Fair die: $P(6)=\frac{1}{6}$. Odds in favour of a 6 $= 1:5$. Odds against $= 5:1$.
\item Drawing a heart from a pack: $P(\text{heart})=\frac{1}{4}$. Odds in favour $= 1:3$.
\item If the odds against an event are $7:2$, then $P(\text{event}) = \frac{2}{7+2} = \frac{2}{9}$.
\end{itemize}
\end{exemplebox}

\courssub{Worked Examination-Style Examples}

\begin{exoresolu}[Two dice — sum and product]
Two fair six-sided dice are rolled. Find the probability that:
\begin{enumerate}
\item the sum of the scores is $8$;
\item the product of the scores is $12$;
\item the sum is $8$ \textbf{or} the product is $12$.
\end{enumerate}
\tcblower
\textbf{Solution.}
The sample space has $n(S)=6\times6=36$ equally likely ordered pairs $(d_1,d_2)$.

\textbf{1.} Sum $=8$: favourable pairs $(2,6),(3,5),(4,4),(5,3),(6,2)$ $\Rightarrow n=5$.
\[
P(\text{sum}=8)=\frac{5}{36}.
\]

\textbf{2.} Product $=12$: favourable pairs $(2,6),(3,4),(4,3),(6,2)$ $\Rightarrow n=4$.
\[
P(\text{product}=12)=\frac{4}{36}=\frac{1}{9}.
\]

\textbf{3.} Let $A$ = ``sum is 8'' and $B$ = ``product is 12''. Then
$A\cap B = \{(2,6),(6,2)\}$ $\Rightarrow n(A\cap B)=2$.
By the addition law (studied in detail in the next section):
\[
P(A\cup B)=P(A)+P(B)-P(A\cap B)=\frac{5}{36}+\frac{4}{36}-\frac{2}{36}=\frac{7}{36}.
\]
\end{exoresolu}

\begin{exoresolu}[Balls drawn without replacement]
A box contains $4$ red and $6$ blue balls. Two balls are drawn \textbf{without replacement}. Find the probability that:
\begin{enumerate}
\item both are red;
\item one is red and one is blue (in any order);
\item at least one is red.
\end{enumerate}
\tcblower
\textbf{Solution.}
Total balls $=10$. The draws are without replacement, so the second draw depends on the first; we use conditional reasoning (or, equivalently, combinations).

\textbf{1.} $P(\text{both red}) = \dfrac{4}{10}\times\dfrac{3}{9} = \dfrac{12}{90} = \dfrac{2}{15}$.

Check with combinations: $\dfrac{\binom{4}{2}}{\binom{10}{2}} = \dfrac{6}{45} = \dfrac{2}{15}$. \checkmark

\textbf{2.} $P(\text{one red, one blue}) = P(RB)+P(BR) = \dfrac{4}{10}\times\dfrac{6}{9} + \dfrac{6}{10}\times\dfrac{4}{9} = \dfrac{24}{90}+\dfrac{24}{90}=\dfrac{48}{90}=\dfrac{8}{15}$.

Check with combinations: $\dfrac{\binom{4}{1}\binom{6}{1}}{\binom{10}{2}} = \dfrac{24}{45} = \dfrac{8}{15}$. \checkmark

\textbf{3.} Using the complement:
\[
P(\text{at least one red}) = 1 - P(\text{both blue}) = 1 - \frac{6}{10}\times\frac{5}{9} = 1 - \frac{30}{90} = \frac{60}{90} = \frac{2}{3}.
\]
(The complement is faster than adding the three cases RR, RB, BR.)
\end{exoresolu}

\courssub{Consolidation Exercises}

\begin{exercice}[Classical Probability and Event Types]
\begin{enumerate}
\item A fair coin is tossed three times. List the sample space. Find the probability of: (a) exactly two heads; (b) at least two heads; (c) at most one tail.
\item A card is drawn from a standard $52$-card pack. Find the probability that it is: (a) a king or a queen; (b) a spade or a face card; (c) neither a heart nor a king.
\item A bag contains $5$ white, $4$ black and $3$ red balls. Three balls are drawn \textbf{with replacement}. Find the probability that: (a) all three are white; (b) exactly two are black; (c) at least one is red.
\item The odds against a student passing a test are $3:2$. What is the probability that the student passes? If two independent students take the test, what is the probability that both pass?
\item Events $A$ and $B$ are such that $P(A)=0.4$, $P(B)=0.5$ and $P(A\cup B)=0.7$. Determine whether $A$ and $B$ are: (a) mutually exclusive; (b) independent.
\end{enumerate}
\end{exercice}

\begin{corrige}[Exercise 1]
\begin{enumerate}
\item Sample space $S=\{HHH,HHT,HTH,THH,HTT,THT,TTH,TTT\}$, $n(S)=8$.
  \begin{itemize}
  \item (a) Exactly two heads: $\{HHT,HTH,THH\}$ $\Rightarrow P=\frac{3}{8}$.
  \item (b) At least two heads: $\{HHH,HHT,HTH,THH\}$ $\Rightarrow P=\frac{4}{8}=\frac{1}{2}$.
  \item (c) At most one tail means at least two heads $\Rightarrow P=\frac{1}{2}$.
  \end{itemize}
\item \begin{itemize}
  \item (a) $P(K\cup Q)=\frac{4+4}{52}=\frac{8}{52}=\frac{2}{13}$ (mutually exclusive).
  \item (b) $P(\text{spade}\cup\text{face}) = \frac{13}{52}+\frac{12}{52}-\frac{3}{52}=\frac{22}{52}=\frac{11}{26}$ (3 face cards are spades).
  \item (c) $P(\text{heart}\cup\text{king}) = \frac{13}{52}+\frac{4}{52}-\frac{1}{52}=\frac{16}{52}=\frac{4}{13}$. Complement: $1-\frac{4}{13}=\frac{9}{13}$.
  \end{itemize}
\item With replacement $\Rightarrow$ independent draws; the bag always contains $12$ balls.
  \begin{itemize}
  \item (a) $P(WWW)=\left(\frac{5}{12}\right)^3=\frac{125}{1728}$.
  \item (b) Exactly two black: $\binom{3}{2}\left(\frac{4}{12}\right)^2\left(\frac{8}{12}\right) = 3\times\frac{1}{9}\times\frac{2}{3}=\frac{2}{9}$.
  \item (c) $P(\text{at least one red}) = 1 - P(\text{no red}) = 1 - \left(\frac{9}{12}\right)^3 = 1 - \frac{27}{64} = \frac{37}{64}$.
  \end{itemize}
\item Odds against $=3:2 \Rightarrow P(\text{pass})=\frac{2}{3+2}=\frac{2}{5}=0.4$. Two independent students: $P(\text{both pass})=(0.4)^2=0.16$.
\item \begin{itemize}
  \item (a) If $A$ and $B$ were mutually exclusive, we would have $P(A\cup B)=P(A)+P(B)=0.9 \neq 0.7$ $\Rightarrow$ they are \textbf{not} mutually exclusive.
  \item (b) First compute $P(A\cap B)=P(A)+P(B)-P(A\cup B)=0.4+0.5-0.7=0.2$. Since $P(A)P(B)=0.4\times0.5=0.2=P(A\cap B)$, the events \textbf{are} independent.
  \end{itemize}
\end{enumerate}
\end{corrige}

\begin{aretenir}[Key Points of This Section]
\begin{itemize}
\item Classical probability $P(A)=\frac{n(A)}{n(S)}$ applies \textbf{only} when all outcomes are equally likely.
\item $0\le P(A)\le 1$; $P(S)=1$; $P(\varnothing)=0$.
\item Mutually exclusive events cannot occur together ($A\cap B=\varnothing$).
\item Exhaustive events cover the whole sample space; if they are also pairwise disjoint, their probabilities sum to $1$.
\item Complementary events are mutually exclusive and exhaustive: $P(A')=1-P(A)$.
\item Use the complement for ``at least one'' problems: $P(\text{at least one})=1-P(\text{none})$.
\item Odds in favour $= p:(1-p)$; odds against $= (1-p):p$; odds $a:b$ in favour means $P(A)=\frac{a}{a+b}$.
\item \textbf{Mutually exclusive $\neq$ independent} — this distinction is crucial for the coming sections.
\end{itemize}
\end{aretenir}

\begin{savaistu}[Did You Know?]
The classical definition of probability, stated by Pierre-Simon Laplace in 1812, was the first rigorous attempt to quantify uncertainty. It works beautifully for games of chance — dice, cards, lotteries — because symmetry guarantees equally likely outcomes. However, it cannot handle questions like ``What is the probability that it will rain tomorrow in Yaoundé?'' That requires the \emph{frequentist} (empirical) approach, which we meet in the next section, or the \emph{Bayesian} approach, which we shall meet at the end of this chapter through Bayes' theorem.
\end{savaistu}

\coursec{Empirical Probability and the Addition Laws}
\courssub{From Classical to Empirical Probability}

In the previous section we studied \cle{classical (theoretical) probability}: when all outcomes of an experiment are equally likely, $P(A)=\dfrac{n(A)}{n(S)}$. But life is not always a fair die. What is the probability that a drawing pin lands point up? That a newborn baby in Bamenda is a girl? That a candidate passes the GCE at first attempt? No symmetry argument can answer these questions. We must \emph{experiment} and observe. This leads us to \cle{empirical probability}, and then to two fundamental laws — the \cle{addition laws} — that let us combine probabilities of events.

\courssub{Empirical (Experimental) Probability}

\begin{experience}[Tossing a coin 100 times]
Take a coin and toss it 100 times, recording the number of heads after 10, 20, 30, \dots, 100 tosses. A class in Buea obtained: 7 heads in 10 tosses, then 11, 16, 23, 26, 34, 38, 44, 47, and finally 51 heads in 100 tosses. The \emph{relative frequencies} of heads are:
\[ \frac{7}{10}=0.70,\quad \frac{11}{20}=0.55,\quad \frac{16}{30}\approx 0.53,\quad \dots,\quad \frac{51}{100}=0.51. \]
The values fluctuate wildly at first, then settle down near $0.5$.
\end{experience}

\begin{definition}[Empirical probability (relative frequency)]
If an experiment is repeated $n$ times under identical conditions and an event $A$ occurs $f$ times, the \cle{relative frequency} of $A$ is
\[ f_n(A)=\frac{f}{n}=\frac{\text{number of times }A\text{ occurs}}{\text{number of trials}}. \]
The \cle{empirical (experimental) probability} of $A$ is the value towards which $f_n(A)$ stabilises as $n$ becomes large; in practice we estimate $P(A)\approx f_n(A)$ for large $n$.
\end{definition}

\begin{propriete}[Stability of relative frequency — informal law of large numbers]
As the number of trials $n$ increases, the relative frequency $f_n(A)$ tends to stabilise around a fixed number, which is the true probability $P(A)$. For a fair coin, $f_n(\text{heads})\to 0.5$. This is why empirical and theoretical probability agree when the theoretical model is correct. (Proof not examinable.)
\end{propriete}

\begin{popfigure}
\begin{center}
\begin{tikzpicture}
\begin{axis}[width=11cm,height=6.2cm,xlabel={Number of tosses $n$},ylabel={Relative frequency of heads},xmin=0,xmax=100,ymin=0.2,ymax=0.8,grid=both,grid style={popline},xtick={0,20,40,60,80,100},ytick={0.2,0.3,0.4,0.5,0.6,0.7,0.8}]
\addplot[color=popblue,thick,mark=*,mark size=1.5pt] coordinates {(10,0.70) (20,0.55) (30,0.533) (40,0.575) (50,0.52) (60,0.567) (70,0.543) (80,0.55) (90,0.522) (100,0.51)};
\addplot[color=poporange,dashed,thick,domain=0:100] {0.5};
\node[color=poporange,anchor=south west] at (axis cs:55,0.5) {\small $y=0.5$};
\end{axis}
\end{tikzpicture}
\end{center}
\legende{Relative frequency of heads against the number of tosses: the fluctuations die out and the frequency converges to the theoretical value $0.5$ (dashed line).}
\end{popfigure}

\begin{exemplebox}[The drawing pin experiment]
A drawing pin is thrown 200 times and lands point up 118 times. There is no symmetry here, so classical probability is useless. The empirical estimate is
\[ P(\text{point up})\approx \frac{118}{200}=0.59. \]
Similarly, a health centre in Yaoundé records 520 girls among 1000 births in a year: the empirical probability that a newborn is a girl is $\frac{520}{1000}=0.52$.
\end{exemplebox}

\begin{attention}[Small samples mislead]
With only 10 tosses, a relative frequency of $0.7$ tells us almost nothing. Empirical probability is reliable only for a \cle{large} number of trials. Never quote an empirical probability from a handful of trials.
\end{attention}

\courssub{Addition Law for Mutually Exclusive Events}

Recall from Section 2: two events $A$ and $B$ are \cle{mutually exclusive} when they cannot occur together, i.e. $A\cap B=\varnothing$.

\begin{propriete}[Addition law for mutually exclusive events]
If $A$ and $B$ are mutually exclusive events of a sample space $S$ with equally likely outcomes, then
\[ P(A\cup B)=P(A)+P(B). \]
\textbf{Proof (examinable).} Since $A\cap B=\varnothing$, the outcomes of $A$ and $B$ are distinct, so $n(A\cup B)=n(A)+n(B)$. Dividing by $n(S)$:
\[ P(A\cup B)=\frac{n(A\cup B)}{n(S)}=\frac{n(A)}{n(S)}+\frac{n(B)}{n(S)}=P(A)+P(B). \]
\textbf{Extension.} If $A_1,A_2,\dots,A_k$ are pairwise mutually exclusive, then
\[ P(A_1\cup A_2\cup\cdots\cup A_k)=P(A_1)+P(A_2)+\cdots+P(A_k). \]
\end{propriete}

\begin{exemplebox}[One card from a deck]
A card is drawn from a standard pack of 52. Let $A=$ ``draw a king'' and $B=$ ``draw a queen''. These are mutually exclusive, so
\[ P(A\cup B)=\frac{4}{52}+\frac{4}{52}=\frac{8}{52}=\frac{2}{13}. \]
\end{exemplebox}

\courssub{General Addition Law for Non-Mutually Exclusive Events}

What if the events overlap? Suppose 30 students of an Upper Sixth class study French ($F$) and 25 study English ($E$), but 12 study \emph{both}. Then $P(F)+P(E)$ counts the 12 bilingual students \emph{twice}. We must subtract the intersection once.

\begin{propriete}[General addition law]
For any two events $A$ and $B$,
\[ P(A\cup B)=P(A)+P(B)-P(A\cap B). \]
\textbf{Proof (examinable).} In the count $n(A)+n(B)$, every element of $A\cap B$ is counted twice, so
\[ n(A\cup B)=n(A)+n(B)-n(A\cap B). \]
Dividing by $n(S)$ gives the result. When $A\cap B=\varnothing$, $P(A\cap B)=0$ and we recover the mutually exclusive law.
\end{propriete}

\begin{methode}[Filling a Venn diagram or two-way table from the inside out]
\begin{enumerate}
\item Identify the \cle{innermost region} first: place $n(A\cap B)$ in the overlap.
\item Complete each circle: $n(A\text{ only})=n(A)-n(A\cap B)$, and similarly for $B$.
\item Place $n(\text{neither})=n(S)-n(A\cup B)$ outside the circles.
\item Check that all regions sum to $n(S)$, then read off any probability.
\end{enumerate}
\end{methode}

\begin{exoresolu}[French and English students]
In a class of 60 students, 30 study French, 25 study English and 12 study both. A student is chosen at random. Find the probability that the student studies (a) French or English; (b) neither.
\tcblower
\textbf{Solution.} Fill the Venn diagram from the inside out: $n(F\cap E)=12$; French only: $30-12=18$; English only: $25-12=13$; neither: $60-(18+12+13)=17$.
\begin{align*}
\text{(a)}\quad P(F\cup E)&=P(F)+P(E)-P(F\cap E)=\frac{30}{60}+\frac{25}{60}-\frac{12}{60}=\frac{43}{60}.\\
\text{(b)}\quad P(\text{neither})&=\frac{17}{60}.
\end{align*}
Check with the diagram: $\dfrac{18+12+13}{60}=\dfrac{43}{60}$. \checkmark
\end{exoresolu}

\begin{popfigure}
\begin{center}
\begin{tikzpicture}[scale=1.05]
\draw[thick,popdark] (-0.4,-0.3) rectangle (5.4,3.5);
\node[popdark] at (5.0,3.2) {$S$};
\draw[thick,popblue,fill=popblueL,opacity=0.85] (1.7,1.6) circle (1.3);
\draw[thick,poporange,fill=popgold,opacity=0.55] (3.3,1.6) circle (1.3);
\node at (1.15,1.6) {\textbf{18}};
\node at (2.5,1.6) {\textbf{12}};
\node at (3.85,1.6) {\textbf{13}};
\node at (4.75,0.25) {\textbf{17}};
\node[popblue] at (1.15,3.05) {$F$};
\node[poporange] at (3.85,3.05) {$E$};
\end{tikzpicture}
\end{center}
\legende{Venn diagram of the 60 students: $18+12+13+17=60$, and $P(F\cup E)=\frac{18+12+13}{60}=\frac{43}{60}$, confirming $P(F)+P(E)-P(F\cap E)$.}
\end{popfigure}

The same information can be organised in a \cle{two-way table}:

\begin{center}
\begin{tabularx}{\linewidth}{|l|X|X|X|}
\hline
\rowcolor{popblueL} & English & Not English & Total \\
\hline
French & 12 & 18 & 30 \\
\hline
Not French & 13 & 17 & 30 \\
\hline
Total & 25 & 35 & 60 \\
\hline
\end{tabularx}
\end{center}

\begin{attention}[The double-counting trap]
The most common error at the GCE is writing $P(A\cup B)=P(A)+P(B)$ when $A$ and $B$ \emph{overlap}. This counts $A\cap B$ twice and can even give a probability greater than 1 — an immediate sign of error. Always ask first: \emph{can the two events happen together?} If yes, subtract $P(A\cap B)$.
\end{attention}

\begin{exercice}[Survey at a market]
Of 80 shoppers interviewed in Douala, 45 bought plantains, 38 bought tomatoes and 20 bought both. A shopper is chosen at random. Find the probability that the shopper bought (a) plantains or tomatoes; (b) tomatoes only; (c) neither.
\end{exercice}

\begin{corrige}[Exercise 1]
Both: 20; plantains only: $45-20=25$; tomatoes only: $38-20=18$; neither: $80-(25+20+18)=17$.
(a) $P=\dfrac{45}{80}+\dfrac{38}{80}-\dfrac{20}{80}=\dfrac{63}{80}$. \quad (b) $P=\dfrac{18}{80}=\dfrac{9}{40}$. \quad (c) $P=\dfrac{17}{80}$.
\end{corrige}

\begin{aretenir}[Key points]
\begin{itemize}
\item Empirical probability: $P(A)\approx \dfrac{\text{number of times }A\text{ occurs}}{\text{number of trials}}$, reliable only for many trials; relative frequency stabilises towards the true probability (law of large numbers).
\item Mutually exclusive: $P(A\cup B)=P(A)+P(B)$; this extends to any number of pairwise mutually exclusive events.
\item General case: $P(A\cup B)=P(A)+P(B)-P(A\cap B)$.
\item Organise overlapping data with a Venn diagram or two-way table, filled from the intersection outwards.
\end{itemize}
\end{aretenir}

\coursec{Multiplication Law, Conditional Probability and Independence}

In the previous section we studied empirical probability and the \emph{addition} laws, which answer the question: ``What is the probability that $A$ \textbf{or} $B$ occurs?'' In real life, however, we very often need to answer a different question: ``What is the probability that $A$ \textbf{and} $B$ occur?'' or even ``What is the probability of $A$ \emph{knowing that} $B$ has already happened?'' Think of a doctor in Bamenda who knows a patient has fever ($B$) and wants the probability of malaria ($A$); or a student who has already drawn one red pen from a box and asks for the probability of drawing a second red pen. These questions lead us to \cle{conditional probability}, the \cle{multiplication law} and the fundamental notion of \cle{independence}.

\courssub{Conditional Probability: Probability with Extra Information}

\textbf{Intuition first.}\par
A school in Buea has $200$ Upper Sixth students, classified by sex and by the field they intend to study:

\begin{center}
\begin{tabularx}{\linewidth}{|l|X|X|X|}
\hline
\rowcolor{popblueL} & Science ($A$) & Arts ($A'$) & Total \\
\hline
Boys ($B$) & $60$ & $30$ & $90$ \\
\hline
Girls ($B'$) & $50$ & $60$ & $110$ \\
\hline
Total & $110$ & $90$ & $200$ \\
\hline
\end{tabularx}
\end{center}

A student is chosen at random. The probability that the student intends to study science is $P(A)=\dfrac{110}{200}=0.55$. But suppose we are \emph{told} that the chosen student is a boy. This extra information changes everything: the whole school is no longer our universe. We now pick from the $90$ boys only, of whom $60$ choose science, so the probability becomes $\dfrac{60}{90}=\dfrac{2}{3}\approx 0.667$. We write this as $P(A\mid B)$, read ``the probability of $A$ \textbf{given} $B$''.

The key idea is that knowing $B$ has occurred \textbf{reduces the sample space} from $S$ to $B$. Inside this reduced sample space, the part where $A$ also occurs is exactly $A\cap B$.

\begin{popfigure}
\begin{center}
\begin{tikzpicture}[scale=1.0]
% Sample space
\draw[thick, popdark] (0,0) rectangle (8,4.6);
\node[popdark, anchor=north west] at (0.1,4.55) {$S$};
% Event B (shaded) and event A
\fill[popblueL, opacity=0.85] (2.2,0.7) circle (1.7);
\fill[poporange, opacity=0.45] (4.4,2.6) circle (1.7);
\draw[thick, popblue] (2.2,0.7) circle (1.7);
\draw[thick, poporange] (4.4,2.6) circle (1.7);
\node[popblue] at (1.1,0.9) {$B$};
\node[poporange] at (5.7,3.4) {$A$};
\node[popdark] at (3.35,1.6) {$A\cap B$};
\node[align=center, popmuted] at (6.7,0.6) {reduced sample\\ space $=B$};
\draw[->, popmuted] (6.2,0.75) -- (3.9,1.1);
\end{tikzpicture}
\end{center}
\legende{Reduced sample space: once $B$ is known to have occurred, only the shaded disc $B$ remains possible, and $P(A\mid B)$ is the proportion of $B$ covered by $A\cap B$.}
\end{popfigure}

For equally likely outcomes, counting inside the reduced sample space gives
\[
P(A\mid B)=\frac{n(A\cap B)}{n(B)}=\frac{n(A\cap B)/n(S)}{n(B)/n(S)}=\frac{P(A\cap B)}{P(B)}.
\]
This motivates the general definition.

\begin{definition}[Conditional probability]
Let $A$ and $B$ be two events of a sample space $S$ with $P(B)>0$. The \cle{conditional probability} of $A$ given $B$ is
\[
P(A\mid B)=\frac{P(A\cap B)}{P(B)}.
\]
It is the probability that $A$ occurs \emph{knowing that} $B$ has occurred. Geometrically, it is the proportion of the event $B$ that is covered by $A\cap B$.
\end{definition}

\begin{exemplebox}[Computing from the contingency table]
Using the table of $200$ students above, with $A=$ ``intends to study science'' and $B=$ ``is a boy'':
\[
P(A\mid B)=\frac{P(A\cap B)}{P(B)}=\frac{60/200}{90/200}=\frac{60}{90}=\frac{2}{3}.
\]
Similarly, the probability that a student is a boy \emph{given} that he or she intends to study science is
\[
P(B\mid A)=\frac{P(A\cap B)}{P(A)}=\frac{60/200}{110/200}=\frac{60}{110}=\frac{6}{11}.
\]
Note that $P(A\mid B)=\tfrac{2}{3}\neq \tfrac{6}{11}=P(B\mid A)$.
\end{exemplebox}

\begin{attention}[$P(A\mid B)$ and $P(B\mid A)$ are generally different]
The word ``given'' tells you which event becomes the \emph{new sample space} (the denominator). $P(A\mid B)$ uses $P(B)$ in the denominator; $P(B\mid A)$ uses $P(A)$. Confusing the two is one of the most frequent errors at the GCE. Always read the sentence slowly and underline what is \emph{known}.
\end{attention}

\begin{exoresolu}[Conditional probability from a Venn diagram]
In a class, $P(A)=0.5$, $P(B)=0.4$ and $P(A\cap B)=0.2$. Find (a) $P(A\mid B)$; (b) $P(B\mid A)$; (c) $P(A\mid B')$.
\tcblower
\textbf{Solution.}
\begin{enumerate}
\item $P(A\mid B)=\dfrac{P(A\cap B)}{P(B)}=\dfrac{0.2}{0.4}=0.5$.
\item $P(B\mid A)=\dfrac{P(A\cap B)}{P(A)}=\dfrac{0.2}{0.5}=0.4$.
\item First $P(A\cap B')=P(A)-P(A\cap B)=0.5-0.2=0.3$, and $P(B')=1-0.4=0.6$. Hence
\[
P(A\mid B')=\frac{P(A\cap B')}{P(B')}=\frac{0.3}{0.6}=0.5.
\]
\end{enumerate}
\end{exoresolu}

\courssub{The Multiplication Law}

Rearranging the definition of conditional probability (multiplying both sides by $P(B)$) immediately gives the rule for the probability of an \emph{intersection}.

\begin{propriete}[Multiplication law of probability]
For any two events $A$ and $B$ with $P(A)>0$ and $P(B)>0$,
\[
P(A\cap B)=P(A)\times P(B\mid A)=P(B)\times P(A\mid B).
\]
(\emph{Proof examinable}: it is a direct rearrangement of the definition $P(A\mid B)=P(A\cap B)/P(B)$.)
\end{propriete}

The law extends naturally to three or more events, for example
\[
P(A\cap B\cap C)=P(A)\,P(B\mid A)\,P(C\mid A\cap B).
\]

\begin{exemplebox}[Drawing without replacement]
A box contains $5$ red and $3$ blue biros. Two biros are drawn at random \emph{without replacement}. What is the probability that both are red?

Let $R_1=$ ``first biro red'' and $R_2=$ ``second biro red''. Then
\[
P(R_1\cap R_2)=P(R_1)\times P(R_2\mid R_1)=\frac{5}{8}\times\frac{4}{7}=\frac{20}{56}=\frac{5}{14}.
\]
After the first red biro is removed, only $4$ red biros remain out of $7$: this is exactly the conditional probability $P(R_2\mid R_1)$.
\end{exemplebox}

\courssub{Independent Events}

\textbf{Intuition.}\par
Sometimes knowing that $B$ has occurred gives \emph{no new information} about $A$: the probability of $A$ stays the same. Tossing a coin and rolling a die: knowing the coin shows Head does not change the probability of obtaining a $6$. Such events are called \cle{independent}.

If $A$ and $B$ are independent, then $P(A\mid B)=P(A)$, and the multiplication law becomes $P(A\cap B)=P(A)\,P(B)$. We take this simpler identity as the definition, because it is symmetric and works even when one probability is zero.

\begin{definition}[Independent events]
Two events $A$ and $B$ are \cle{independent} if and only if
\[
P(A\cap B)=P(A)\times P(B).
\]
When $P(A)>0$ and $P(B)>0$, this is equivalent to each of the conditions
\[
P(A\mid B)=P(A)\qquad\text{and}\qquad P(B\mid A)=P(B).
\]
Otherwise the events are \cle{dependent}.
\end{definition}

\begin{propriete}[Independence of complements]
If $A$ and $B$ are independent, then so are the pairs $(A,B')$, $(A',B)$ and $(A',B')$.
\end{propriete}

\textbf{Proof (instructive).} Since $A=(A\cap B)\cup(A\cap B')$, a union of disjoint events,
\[
P(A\cap B')=P(A)-P(A\cap B)=P(A)-P(A)P(B)=P(A)\bigl(1-P(B)\bigr)=P(A)P(B').
\]
Hence $A$ and $B'$ are independent. Applying the same argument to $A'$ and $B$, and then to $A'$ and $B'$, completes the proof. \hfill$\blacksquare$

\begin{methode}[Testing independence numerically]
To decide whether two events $A$ and $B$ are independent from data or given probabilities:
\begin{enumerate}
\item Compute (or read off) $P(A)$, $P(B)$ and $P(A\cap B)$.
\item Compute the product $P(A)\times P(B)$.
\item Compare: if $P(A\cap B)=P(A)P(B)$, the events are independent; otherwise they are dependent.
\item Equivalently, check whether $P(A\mid B)=P(A)$.
\end{enumerate}
\end{methode}

\begin{exoresolu}[Are the events independent?]
Using the table of $200$ students: $A=$ ``intends to study science'', $B=$ ``is a boy''. Test whether $A$ and $B$ are independent.
\tcblower
\textbf{Solution.} From the table:
\[
P(A)=\frac{110}{200}=0.55,\qquad P(B)=\frac{90}{200}=0.45,\qquad P(A\cap B)=\frac{60}{200}=0.30.
\]
Now $P(A)\times P(B)=0.55\times 0.45=0.2475\neq 0.30=P(A\cap B)$.

Therefore $A$ and $B$ are \textbf{not independent}: the intention to study science depends on sex in this school. Equivalently, $P(A\mid B)=\tfrac{2}{3}\approx 0.667\neq 0.55=P(A)$.
\end{exoresolu}

\begin{exemplebox}[An independent situation]
A fair coin is tossed and a fair die is rolled. Let $A=$ ``Head'' and $B=$ ``a six''. Then
\[
P(A)=\tfrac12,\quad P(B)=\tfrac16,\quad P(A\cap B)=\tfrac{1}{12}=\tfrac12\times\tfrac16.
\]
The events are independent, as expected: the coin cannot influence the die.
\end{exemplebox}

\courssub{Mutually Exclusive versus Independent: a Classic Trap}

These two notions are constantly confused by candidates, yet they are almost \emph{opposite} ideas:
\begin{itemize}
\item \textbf{Mutually exclusive} (disjoint): $A\cap B=\varnothing$, i.e.\ $P(A\cap B)=0$. The occurrence of one \emph{prevents} the other.
\item \textbf{Independent}: $P(A\cap B)=P(A)P(B)$. The occurrence of one gives \emph{no information} about the other.
\end{itemize}

\begin{attention}[Disjoint events with positive probabilities are never independent]
If $A$ and $B$ are mutually exclusive with $P(A)>0$ and $P(B)>0$, then
\[
P(A\cap B)=0\neq P(A)P(B)>0,
\]
so $A$ and $B$ are \textbf{dependent}. Indeed, if $B$ occurs then $A$ cannot occur: $P(A\mid B)=0\neq P(A)$. Knowing $B$ completely determines $A$ — the strongest possible dependence!

\textbf{Counter-example to keep in mind.} Drawing one card from a pack: $A=$ ``king'', $B=$ ``queen''. They are mutually exclusive but clearly not independent. Conversely, $A=$ ``king'' and $C=$ ``heart'' are independent ($P(A\cap C)=\tfrac{1}{52}=\tfrac{4}{52}\times\tfrac{13}{52}$) but not mutually exclusive (the king of hearts exists). Never write $P(A\cap B)=P(A)P(B)$ for disjoint events, and never write $P(A\cup B)=P(A)+P(B)$ for events that overlap.
\end{attention}

\courssub{Exhaustive Events and a Preview of Total Probability}

Recall that events $B_1,B_2,\dots,B_n$ form a \cle{partition} of $S$ (they are \emph{exhaustive} and mutually exclusive) when every outcome belongs to exactly one of them. Then any event $A$ can be split into disjoint pieces:
\[
A=(A\cap B_1)\cup(A\cap B_2)\cup\dots\cup(A\cap B_n).
\]
Applying the addition law for mutually exclusive events, then the multiplication law to each piece:
\[
P(A)=\sum_{i=1}^{n}P(A\cap B_i)=\sum_{i=1}^{n}P(B_i)\,P(A\mid B_i).
\]

\begin{propriete}[Law of total probability — preview]
If $B_1,B_2,\dots,B_n$ form a partition of $S$ with $P(B_i)>0$ for all $i$, then for any event $A$,
\[
P(A)=P(B_1)P(A\mid B_1)+P(B_2)P(A\mid B_2)+\dots+P(B_n)P(A\mid B_n).
\]
This result will be proved and used fully in Section 7, where it becomes the engine of \emph{Bayes' theorem}.
\end{propriete}

\begin{exemplebox}[Splitting over a partition]
A market seller in Douala buys tomatoes from two farms: $60\%$ from farm $B_1$ and $40\%$ from farm $B_2$. $5\%$ of farm $B_1$'s tomatoes are spoilt and $10\%$ of farm $B_2$'s are spoilt. Let $A=$ ``a randomly picked tomato is spoilt''. Then
\[
P(A)=P(B_1)P(A\mid B_1)+P(B_2)P(A\mid B_2)=(0.6)(0.05)+(0.4)(0.10)=0.03+0.04=0.07.
\]
So $7\%$ of the tomatoes are spoilt overall — computed by splitting $A$ over the partition $\{B_1,B_2\}$.
\end{exemplebox}

\begin{aretenir}[Key points of this section]
\begin{itemize}
\item Conditional probability: $P(A\mid B)=\dfrac{P(A\cap B)}{P(B)}$, $P(B)>0$ — the sample space is reduced to $B$.
\item Multiplication law: $P(A\cap B)=P(A)\,P(B\mid A)=P(B)\,P(A\mid B)$.
\item Independence: $P(A\cap B)=P(A)P(B)$, equivalently $P(A\mid B)=P(A)$; complements of independent events remain independent.
\item Mutually exclusive $\neq$ independent: disjoint events with positive probabilities are always dependent.
\item Total probability: if $B_1,\dots,B_n$ partition $S$, then $P(A)=\sum_i P(B_i)P(A\mid B_i)$.
\end{itemize}
\end{aretenir}

\begin{exercice}[Practice]
A bag contains $4$ white and $6$ black balls. Two balls are drawn at random without replacement. Let $A=$ ``first ball white'' and $B=$ ``second ball white''.
\begin{enumerate}
\item Compute $P(A)$, $P(B\mid A)$ and $P(A\cap B)$.
\item Using the law of total probability over the partition $\{A,A'\}$, compute $P(B)$.
\item Are $A$ and $B$ independent? Justify.
\end{enumerate}
\end{exercice}

\begin{corrige}[Exercise 1]
\begin{enumerate}
\item $P(A)=\dfrac{4}{10}=\dfrac{2}{5}$; $P(B\mid A)=\dfrac{3}{9}=\dfrac{1}{3}$; $P(A\cap B)=\dfrac{2}{5}\times\dfrac{1}{3}=\dfrac{2}{15}$.
\item $P(B)=P(A)P(B\mid A)+P(A')P(B\mid A')=\dfrac{4}{10}\times\dfrac{3}{9}+\dfrac{6}{10}\times\dfrac{4}{9}=\dfrac{12+24}{90}=\dfrac{2}{5}$.
\item $P(A)P(B)=\dfrac{2}{5}\times\dfrac{2}{5}=\dfrac{4}{25}\neq\dfrac{2}{15}=P(A\cap B)$, so $A$ and $B$ are \textbf{dependent} (drawing without replacement creates dependence).
\end{enumerate}
\end{corrige}

\coursec{Tree Diagrams and Selection With or Without Replacement}

In the previous section we met the multiplication law $P(A\cap B)=P(A)\times P(B\mid A)$ and the notion of independence. These ideas are powerful, but when an experiment has \emph{several stages} — two coins tossed one after another, two balls drawn from a bag, a machine tested on three consecutive days — it becomes difficult to keep track of all the cases in one's head. The \cle{tree diagram} is the picture that organises the multiplication law for us: it displays every possible route through the experiment, with the probability of each step written on the branch, so that no case is forgotten and no probability is misapplied.

\courssub{Constructing a Tree Diagram}

\begin{definition}[Tree diagram]
A \cle{tree diagram} for a multi-stage experiment is a diagram in which:
\begin{itemize}
\item each \emph{stage} of the experiment is represented by a level of branching;
\item from each \emph{node}, one branch is drawn for \textbf{each possible outcome} of that stage;
\item the \textbf{probability of that outcome} (given the history so far) is written on the branch;
\item the \textbf{final outcome} of each complete route is written at the end of the path, together with its probability.
\end{itemize}
\end{definition}

\begin{propriete}[Rules of the tree]
\begin{enumerate}
\item \textbf{Node rule:} the probabilities on all branches leaving the same node sum to $1$.
\item \textbf{Path rule (multiplication):} the probability of a complete path is the \emph{product} of the probabilities on its branches.
\item \textbf{Event rule (addition):} the probability of an event is the \emph{sum} of the probabilities of all final outcomes that satisfy the event.
\item \textbf{Check:} the probabilities of \emph{all} final outcomes of the tree sum to $1$.
\end{enumerate}
The path rule is exactly the multiplication law $P(A\cap B)=P(A)\,P(B\mid A)$; the event rule is the addition law for mutually exclusive final outcomes. (Justification examinable as an application of those two laws.)
\end{propriete}

\begin{methode}[Solving a probability problem with a tree]
\begin{enumerate}
\item Identify the \textbf{stages} of the experiment and the possible \textbf{outcomes} at each stage.
\item Draw the tree; write on each branch the probability of that outcome \emph{given what has already happened} (conditional probabilities on second and later levels).
\item \textbf{Multiply} along each path to obtain the probability of each final outcome.
\item \textbf{Check:} the final probabilities sum to $1$.
\item \textbf{Add} the probabilities of exactly those final outcomes that satisfy the required event.
\end{enumerate}
\end{methode}

\begin{attention}[Multiply along a path — never add]
A very frequent error is to \emph{add} the probabilities written along a single path. Along a path, events happen \textbf{one after another} (``and''), so we \textbf{multiply}. We only \textbf{add} when combining \emph{different} final outcomes (``or'').
\end{attention}

\courssub{Selection With and Without Replacement}

Many GCE problems concern objects drawn from a collection: balls from a bag, cards from a pack, items from a production line. The crucial question is: \textbf{is the first object put back before the second draw?}

\begin{definition}[With and without replacement]
\begin{itemize}
\item \textbf{With replacement:} the object drawn is returned before the next draw. The composition of the collection is unchanged, so the probabilities are \emph{the same at every stage} and the stages are \cle{independent}.
\item \textbf{Without replacement:} the object drawn is kept out. The composition changes, so the probabilities at the next stage must be \emph{updated}; the stages are \cle{dependent}, and second-level branch probabilities are conditional probabilities $P(\text{second}\mid\text{first})$.
\end{itemize}
\end{definition}

\begin{methode}[Updating probabilities without replacement]
After a draw without replacement, at the next stage:
\begin{enumerate}
\item decrease by $1$ the number of objects in the \textbf{category just drawn};
\item decrease by $1$ the \textbf{total} number of objects;
\item leave all other categories unchanged.
\end{enumerate}
\end{methode}

\begin{attention}[The most common GCE error]
\textbf{With replacement the denominators stay the same; without replacement they change.} Drawing from a bag of $8$ balls: with replacement, every denominator is $8$; without replacement, the second-stage denominators are all $7$. Mixing these up destroys the whole calculation — always ask first: \emph{is it replaced?}
\end{attention}

\courssub{Worked Example: Two Balls Without Replacement}

\begin{exemplebox}[Bag of 5 red and 3 blue balls]
A bag contains $5$ red and $3$ blue balls. Two balls are drawn at random, one after the other, \textbf{without replacement}. First draw: $P(R)=\tfrac{5}{8}$, $P(B)=\tfrac{3}{8}$. If red was drawn first, $4$ red and $3$ blue remain, so $P(R\mid R)=\tfrac{4}{7}$, $P(B\mid R)=\tfrac{3}{7}$. If blue was drawn first, $5$ red and $2$ blue remain, so $P(R\mid B)=\tfrac{5}{7}$, $P(B\mid B)=\tfrac{2}{7}$.
\end{exemplebox}

\begin{popfigure}
\begin{center}
\begin{tikzpicture}[grow=right, level distance=3.4cm,
  level 1/.style={sibling distance=3.2cm},
  level 2/.style={sibling distance=1.6cm},
  every node/.style={font=\small},
  edge from parent/.style={draw=popblue, thick}]
\node {Start}
  child { node {B}
    child { node {BB \quad $\frac{3}{8}\times\frac{2}{7}=\frac{6}{56}$}
      edge from parent node[above] {$\frac{2}{7}$} }
    child { node {BR \quad $\frac{3}{8}\times\frac{5}{7}=\frac{15}{56}$}
      edge from parent node[above] {$\frac{5}{7}$} }
    edge from parent node[above] {$\frac{3}{8}$} }
  child { node {R}
    child { node {RB \quad $\frac{5}{8}\times\frac{3}{7}=\frac{15}{56}$}
      edge from parent node[above] {$\frac{3}{7}$} }
    child { node {RR \quad $\frac{5}{8}\times\frac{4}{7}=\frac{20}{56}$}
      edge from parent node[above] {$\frac{4}{7}$} }
    edge from parent node[above] {$\frac{5}{8}$} };
\end{tikzpicture}
\end{center}
\legende{Two draws without replacement from 5 red and 3 blue balls. Check: $\frac{20+15+15+6}{56}=1$.}
\end{popfigure}

From the tree: $P(\text{both red})=\tfrac{20}{56}=\tfrac{5}{14}$; $P(\text{one of each colour})=\tfrac{15}{56}+\tfrac{15}{56}=\tfrac{15}{28}$; $P(\text{at least one blue})=1-P(RR)=1-\tfrac{20}{56}=\tfrac{9}{14}$ (complement shortcut).

The tree also answers \emph{conditional} questions directly. For example, ``given that the second ball is blue, what is the probability that the first was red?'' The second ball is blue on paths RB and BB, with total probability $\tfrac{15}{56}+\tfrac{6}{56}=\tfrac{21}{56}$; only RB has first red, so
\[ P(\text{first red}\mid\text{second blue})=\frac{P(RB)}{P(RB)+P(BB)}=\frac{15/56}{21/56}=\frac{15}{21}=\frac{5}{7}. \]

\courssub{Three-Stage Experiments: Three Coin Tosses}

\begin{exoresolu}[Exactly two heads in three tosses]
A fair coin is tossed three times. Find the probability of obtaining exactly two heads.
\tcblower
Each toss gives H or T with probability $\tfrac12$; the tosses are independent, so every complete path has probability $\left(\tfrac12\right)^3=\tfrac18$. The paths with exactly two heads are HHT, HTH, THH — three paths. Hence
\[ P(\text{exactly two heads})=\tfrac18+\tfrac18+\tfrac18=\tfrac38. \]
\end{exoresolu}

\begin{popfigure}
\begin{center}
\begin{tikzpicture}[grow=right, level distance=2.6cm,
  level 1/.style={sibling distance=4.4cm},
  level 2/.style={sibling distance=2.2cm},
  level 3/.style={sibling distance=1.1cm},
  every node/.style={font=\scriptsize},
  edge from parent/.style={draw=popmuted}]
\node {Start}
  child { node {T}
    child { node {T}
      child { node {TTT} edge from parent node[above] {$\frac12$} }
      child { node {TTH} edge from parent node[above] {$\frac12$} }
      edge from parent node[above] {$\frac12$} }
    child { node {H}
      child { node[popgreen] {THT} edge from parent[draw=popgreen, thick] node[above] {$\frac12$} }
      child { node[popgreen] {THH} edge from parent[draw=popgreen, thick] node[above] {$\frac12$} }
      edge from parent[draw=popgreen, thick] node[above] {$\frac12$} }
    edge from parent node[above] {$\frac12$} }
  child { node {H}
    child { node {T}
      child { node {HTT} edge from parent node[above] {$\frac12$} }
      child { node[popgreen] {HTH} edge from parent[draw=popgreen, thick] node[above] {$\frac12$} }
      edge from parent[draw=popgreen, thick] node[above] {$\frac12$} }
    child { node {H}
      child { node[popgreen] {HHT} edge from parent[draw=popgreen, thick] node[above] {$\frac12$} }
      child { node {HHH} edge from parent node[above] {$\frac12$} }
      edge from parent[draw=popgreen, thick] node[above] {$\frac12$} }
    edge from parent[draw=popgreen, thick] node[above] {$\frac12$} };
\end{tikzpicture}
\end{center}
\legende{Three coin tosses. The green paths (HHT, HTH, THH) give exactly two heads: $P=\frac{3}{8}$.}
\end{popfigure}

\courssub{``At least one'' and ``exactly one'' from Trees}

\begin{methode}[The complement shortcut for ``at least one'']
To find $P(\text{at least one success})$:
\begin{enumerate}
\item locate the \textbf{single path} with \emph{no} success (e.g.\ all failures);
\item compute its probability by multiplying along that path;
\item subtract from $1$: $\;P(\text{at least one})=1-P(\text{none})$.
\end{enumerate}
This avoids adding many paths and is much faster in the examination.
\end{methode}

\begin{exemplebox}[Quality control]
A machine produces items of which $10\%$ are defective. Three items are tested independently (equivalent to sampling \emph{with} replacement). Then
\[ P(\text{at least one defective})=1-P(\text{none defective})=1-(0.9)^3=1-0.729=0.271. \]
For ``exactly one defective'', add the three paths DGG, GDG, GGD: $\;3\times(0.1)(0.9)^2=0.243$.
\end{exemplebox}

\begin{attention}[Check your tree before using it]
Before reading off any answer, verify the two quick checks: (i) branches from each node sum to $1$; (ii) all final outcomes sum to $1$. If either fails, a branch probability is wrong — usually a denominator that was not updated after a draw without replacement.
\end{attention}

\begin{savaistu}[Did you know?]
Tree diagrams are the everyday tool of \emph{decision analysis}: doctors use them to combine the probabilities of successive test results, and engineers use them (as ``fault trees'') to estimate the reliability of systems such as the power grid. The same multiply-along, add-across logic you use for balls in a bag drives software that assesses risks in aviation and medicine.
\end{savaistu}

\begin{exercice}
A box contains $4$ white and $6$ black balls. Two balls are drawn at random.
\begin{enumerate}
\item Draw the tree and find $P(\text{both black})$ when the draw is \textbf{with replacement}.
\item Repeat when the draw is \textbf{without replacement}.
\item In each case, find $P(\text{the two balls are of different colours})$.
\end{enumerate}
\end{exercice}

\begin{corrige}[Exercise 1]
\textbf{With replacement:} denominators stay $10$.
$P(BB)=\tfrac{6}{10}\times\tfrac{6}{10}=\tfrac{9}{25}$.
$P(\text{different})=\tfrac{6}{10}\cdot\tfrac{4}{10}+\tfrac{4}{10}\cdot\tfrac{6}{10}=\tfrac{12}{25}$.

\textbf{Without replacement:} denominators become $9$.
$P(BB)=\tfrac{6}{10}\times\tfrac{5}{9}=\tfrac{1}{3}$.
$P(\text{different})=\tfrac{6}{10}\cdot\tfrac{4}{9}+\tfrac{4}{10}\cdot\tfrac{6}{9}=\tfrac{48}{90}=\tfrac{8}{15}$.
\end{corrige}

\begin{exercice}
In a certain school, $60\%$ of the candidates enter for Mathematics and $40\%$ for Biology only at a given session; independently, each candidate passes with probability $0.7$. A candidate is chosen at random.
\begin{enumerate}
\item Draw a two-stage tree (subject, then result) and find the probability that the candidate entered for Mathematics and passed.
\item Find the probability that the candidate passed.
\item Given that the candidate passed, find the probability that she entered for Mathematics.
\end{enumerate}
\end{exercice}

\begin{corrige}[Exercise 2]
First level: $P(M)=0.6$, $P(B)=0.4$. Second level: $P(\text{pass}\mid M)=P(\text{pass}\mid B)=0.7$, $P(\text{fail})=0.3$ on each branch.
\begin{enumerate}
\item $P(M\cap\text{pass})=0.6\times0.7=0.42$.
\item $P(\text{pass})=0.6\times0.7+0.4\times0.7=0.42+0.28=0.70$ (as expected, since the pass rate is the same for both subjects).
\item $P(M\mid\text{pass})=\dfrac{0.42}{0.70}=0.6$.
\end{enumerate}
\end{corrige}

\begin{aretenir}[Key points]
\begin{itemize}
\item One branch per outcome; probability on the branch; outcome at the end.
\item Multiply \emph{along} a path; add \emph{across} final outcomes; all final probabilities sum to $1$.
\item With replacement: probabilities constant, stages independent. Without replacement: decrease the drawn category \emph{and} the total by $1$; second-level branches are conditional probabilities.
\item ``At least one'': use $1-P(\text{none})$. ``Exactly one'': add the relevant paths.
\item Conditional questions from a tree: $P(A\mid B)=\dfrac{P(\text{paths in }A\cap B)}{P(\text{paths in }B)}$.
\end{itemize}
\end{aretenir}

\coursec{Counting Techniques in Probability}

In the previous section we used tree diagrams to analyse selections made with or without replacement. Trees are wonderful for two or three stages, but they collapse under their own weight very quickly: drawing a tree for \emph{choose 3 balls from 10} already requires $10 \times 9 \times 8 = 720$ terminal branches. Yet the GCE examiner will happily ask you for the probability that a committee of 5 chosen from 12 people contains exactly 2 women, or that a sample of 4 items from a crate of 50 contains at least one defective. For such problems we need a faster way to \emph{count} outcomes without listing them. That is exactly what \cle{counting techniques} — the multiplication principle, permutations and combinations — provide. This section shows how they turn impossible listing problems into one-line calculations.

\courssub{Why Counting Matters}

Recall the classical definition of probability from earlier in this chapter: when all outcomes of a sample space $S$ are equally likely,
\[
P(A) = \frac{n(A)}{n(S)} = \frac{\text{number of favourable outcomes}}{\text{total number of outcomes}}.
\]
Everything therefore reduces to \emph{counting two sets}: the favourable outcomes and all outcomes. When $n(S) = 6$ (one die) or $n(S) = 36$ (two dice), listing is feasible. But consider:

\begin{itemize}
\item choosing 3 balls from a bag of 10: the number of possible selections is $120$;
\item forming a committee of 5 from 12 candidates: $792$ possible committees;
\item dealing a hand of 5 cards from a pack of 52: $2\,598\,960$ possible hands.
\end{itemize}

No candidate can list these in an examination. The whole art of this section is to compute $n(A)$ and $n(S)$ \emph{by formula}, and — crucially — to count both with the \emph{same rule}.

\begin{attention}[Count consistently]
Whatever counting rule you use for the numerator $n(A)$, you must use the same rule for the denominator $n(S)$. Mixing ordered counting above with unordered counting below (or vice versa) is the single most common source of wrong answers in this topic.
\end{attention}

\courssub{The Multiplication Principle}

\begin{experience}[Counting routes from Bamenda to Buea]
Suppose a student travels from Bamenda to Buea, passing through Douala. There are $3$ possible coaches from Bamenda to Douala, and $4$ possible coaches from Douala to Buea. How many complete journeys are possible? For each of the $3$ first-leg choices there are $4$ second-leg choices, giving $3 \times 4 = 12$ journeys. No list is needed: the choices \emph{multiply}.
\end{experience}

\begin{propriete}[Multiplication principle of counting]
If an experiment is performed in stages, where the first stage can occur in $n_1$ ways, the second in $n_2$ ways, \dots, and the $k$-th in $n_k$ ways, then the total number of possible outcomes is
\[
n_1 \times n_2 \times \cdots \times n_k.
\]
This holds whether the choices at each stage are independent or not, provided the number of options at each stage does not depend on the specific choices made earlier (only the set of available options may change, but its size is fixed).
\end{propriete}

\begin{exemplebox}[Number plates]
A number plate consists of 2 letters followed by 3 digits. Assuming repetition is allowed, the number of plates is
\[
26 \times 26 \times 10 \times 10 \times 10 = 676\,000.
\]
If repetition is \emph{not} allowed among the digits, the count becomes $26 \times 26 \times 10 \times 9 \times 8 = 486\,720$.
\end{exemplebox}

\courssub{Permutations: When Order Matters}

Suppose 8 athletes contest the 100 m final at the school sports day in Yaoundé. In how many ways can the gold, silver and bronze medals be awarded? There are 8 choices for gold, then 7 for silver, then 6 for bronze: $8 \times 7 \times 6 = 336$ ways. Here \emph{order matters}: "Ngong gold, Etame silver" is a different outcome from "Etame gold, Ngong silver". Such an ordered arrangement is called a \cle{permutation}.

\begin{definition}[Permutation]
A \textbf{permutation} of $r$ objects chosen from $n$ distinct objects is an \emph{ordered} selection. The number of such permutations is
\[
{}^{n}P_{r} = \frac{n!}{(n-r)!} = n(n-1)(n-2)\cdots(n-r+1).
\]
In particular, the number of arrangements of all $n$ objects is ${}^{n}P_{n} = n!$.
\end{definition}

The formula is a direct consequence of the multiplication principle: $n$ choices for the first position, $n-1$ for the second, down to $n-r+1$ for the $r$-th position.

\begin{exemplebox}[Arranging digits]
How many 3-digit numbers can be formed from the digits $1, 2, 3, 4, 5$ without repetition? This is an ordered selection of 3 digits from 5:
\[
{}^{5}P_{3} = \frac{5!}{2!} = 5 \times 4 \times 3 = 60.
\]
Hence, if one such number is picked at random, the probability that it is a \emph{specific} number, say 352, is $\dfrac{1}{60}$.
\end{exemplebox}

\begin{exemplebox}[Seating arrangements]
In how many ways can 6 students be seated in a row? ${}^{6}P_{6} = 6! = 720$. If two particular friends, Amina and Bih, must sit together, treat the pair as one block: $5!$ arrangements of the block with the other 4 students, times $2!$ internal arrangements of the pair, giving $5! \times 2! = 240$. Therefore, if the 6 students sit at random,
\[
P(\text{Amina and Bih together}) = \frac{240}{720} = \frac{1}{3}.
\]
\end{exemplebox}

\courssub{Combinations: When Order Does Not Matter}

Now change the question: from 8 athletes, in how many ways can we \emph{select} 3 to represent the school — with no ranking among them? Choosing "Ngong, Etame, Ali" is the \emph{same} selection as "Ali, Ngong, Etame". Order no longer matters. Each unordered group of 3 corresponds to $3! = 6$ ordered permutations, so we must divide the permutation count by $3!$.

\begin{definition}[Combination]
A \textbf{combination} of $r$ objects chosen from $n$ distinct objects is an \emph{unordered} selection. The number of such combinations is
\[
{}^{n}C_{r} = \binom{n}{r} = \frac{n!}{r!\,(n-r)!} = \frac{{}^{n}P_{r}}{r!}.
\]
\end{definition}

\begin{propriete}[Symmetry of combinations]
\[
{}^{n}C_{r} = {}^{n}C_{n-r}.
\]
Choosing the $r$ objects to \emph{take} is the same as choosing the $n-r$ objects to \emph{leave behind}.
\end{propriete}

\begin{attention}[Use the symmetry to save time]
Never compute ${}^{50}C_{47}$ directly: write ${}^{50}C_{47} = {}^{50}C_{3} = \dfrac{50 \times 49 \times 48}{3 \times 2 \times 1} = 19\,600$. In the examination, always replace $r$ by $n-r$ when $r > \tfrac{n}{2}$. This avoids long factorial arithmetic and reduces the risk of slips.
\end{attention}

\begin{methode}[Permutation or combination?]
\begin{enumerate}
\item Read the question and ask: \emph{does the order of selection distinguish two outcomes?}
\item If \textbf{yes} (rankings, arrangements, codes, positions, first/second/third prizes): use ${}^{n}P_{r}$.
\item If \textbf{no} (committees, groups, hands of cards, samples of items): use ${}^{n}C_{r}$.
\item Count the favourable outcomes $n(A)$ and the total outcomes $n(S)$ \emph{with the same rule}, then form $P(A) = \dfrac{n(A)}{n(S)}$.
\end{enumerate}
\end{methode}

\begin{attention}[Never mix the two rules]
If you count favourable committees with combinations (unordered) but count the total with permutations (ordered), the ratio is meaningless. For example, for "3 balls from 10", $\dfrac{{}^{6}C_{3}}{{}^{10}P_{3}} = \dfrac{20}{720}$ is \textbf{wrong}; the correct value is $\dfrac{{}^{6}C_{3}}{{}^{10}C_{3}} = \dfrac{20}{120} = \dfrac{1}{6}$. Both counts must treat order in the same way.
\end{attention}

\begin{popfigure}
\begin{tikzpicture}[scale=1, every node/.style={font=\small}]
% Bar chart comparing nPr and nCr for n = 10
\begin{axis}[
    ybar, bar width=6pt,
    width=12.5cm, height=7cm,
    xlabel={$r$}, ylabel={number of selections},
    xtick={0,1,2,3,4,5,6,7,8,9,10},
    ymin=0, ymax=200000,
    ymode=log, log basis y=10,
    legend style={at={(0.02,0.98)}, anchor=north west, font=\small},
    ymajorgrids=true, grid style={popline!40},
]
\addplot+[ybar, fill=popblue, draw=popdark] coordinates {(0,1) (1,10) (2,90) (3,720) (4,5040) (5,30240) (6,151200) (7,604800) (8,1814400) (9,3628800) (10,3628800)};
\addplot+[ybar, fill=poporange, draw=popdark] coordinates {(0,1) (1,10) (2,45) (3,120) (4,210) (5,252) (6,210) (7,120) (8,45) (9,10) (10,1)};
\legend{${}^{10}P_{r}$ (ordered), ${}^{10}C_{r}$ (unordered)}
\end{axis}
\end{tikzpicture}
\legende{Comparison of ${}^{10}P_{r}$ and ${}^{10}C_{r}$ for $r = 0, 1, \dots, 10$ (logarithmic scale). Ordered counts are always at least as large as unordered counts, and exceed them by the factor $r!$ as soon as $r \geq 2$.}
\end{popfigure}

\courssub{Computing Probabilities with Combinations}

The standard model is a population split into categories (red/black balls, men/women, defective/good items) from which a sample is drawn \emph{without replacement} and \emph{without order}. The total number of samples is a combination; the favourable samples are counted category by category using the multiplication principle, and the counts are then multiplied.

\begin{exoresolu}[Balls from a bag]
A bag contains 6 red balls and 4 black balls. Three balls are drawn at random, without replacement. Find the probability that all three are red.
\tcblower
Order does not matter, so we use combinations throughout.
\textbf{Total outcomes:} the number of ways of choosing 3 balls from 10 is
\[
n(S) = {}^{10}C_{3} = \frac{10 \times 9 \times 8}{3 \times 2 \times 1} = 120.
\]
\textbf{Favourable outcomes:} the number of ways of choosing 3 red balls from the 6 red ones is
\[
n(A) = {}^{6}C_{3} = \frac{6 \times 5 \times 4}{3 \times 2 \times 1} = 20.
\]
Therefore
\[
P(\text{all red}) = \frac{{}^{6}C_{3}}{{}^{10}C_{3}} = \frac{20}{120} = \frac{1}{6}.
\]
\textbf{Check with conditional probability} (as in the previous section): $\dfrac{6}{10} \times \dfrac{5}{9} \times \dfrac{4}{8} = \dfrac{120}{720} = \dfrac{1}{6}$. The two methods agree, as they must.
\end{exoresolu}

\begin{exoresolu}[A committee with exactly 2 women]
A committee of 3 people is to be chosen at random from 4 men and 6 women. Find the probability that the committee contains exactly 2 women.
\tcblower
\textbf{Total committees:} $\;n(S) = {}^{10}C_{3} = 120$.

\textbf{Favourable committees:} choose 2 women from 6 \emph{and} 1 man from 4. By the multiplication principle,
\[
n(A) = {}^{6}C_{2} \times {}^{4}C_{1} = 15 \times 4 = 60.
\]
Hence
\[
P(\text{exactly 2 women}) = \frac{{}^{6}C_{2}\,{}^{4}C_{1}}{{}^{10}C_{3}} = \frac{60}{120} = \frac{1}{2}.
\]
\end{exoresolu}

\begin{popfigure}
\begin{tikzpicture}[font=\small]
% Box of men
\draw[thick, popblue, fill=popblueL] (0,0) rectangle (3.4,2.2);
\node[popdark] at (1.7,1.7) {\textbf{4 men}};
\node[popdark] at (1.7,0.9) {choose 1};
\node[popdark] at (1.7,0.35) {${}^{4}C_{1} = 4$ ways};
% Box of women
\draw[thick, poppink, fill=poppink!15] (5.6,0) rectangle (9,2.2);
\node[popdark] at (7.3,1.7) {\textbf{6 women}};
\node[popdark] at (7.3,0.9) {choose 2};
\node[popdark] at (7.3,0.35) {${}^{6}C_{2} = 15$ ways};
% Arrows to favourable
\draw[-{Stealth}, very thick, poporange] (1.7,0) -- (1.7,-0.7) -- (4.5,-0.7);
\draw[-{Stealth}, very thick, poporange] (7.3,0) -- (7.3,-0.7) -- (4.5,-0.7);
\node[popdark, align=center] at (4.5,-1.35) {favourable: ${}^{4}C_{1} \times {}^{6}C_{2} = 4 \times 15 = 60$};
% Total
\node[popdark, align=center] at (4.5,-2.15) {total committees: ${}^{10}C_{3} = 120$};
\draw[-{Stealth}, very thick, popgreen] (4.5,-1.7) -- (4.5,-2.55);
\node[popdark, align=center] at (4.5,-3.0) {$P(\text{exactly 2 women}) = \dfrac{60}{120} = \dfrac{1}{2}$};
\end{tikzpicture}
\legende{The committee-selection problem: favourable selections are built category by category (${}^{4}C_{1}$ from the men, ${}^{6}C_{2}$ from the women) and multiplied; the total number of committees is ${}^{10}C_{3}$.}
\end{popfigure}

\courssub{"At Least One": Use the Complement}

Phrases such as \emph{at least one}, \emph{at least two}, \emph{no fewer than} signal that direct counting is tedious (several cases to add) but the \emph{complement} is a single easy case.

\begin{methode}[Complement counting for "at least one"]
\begin{enumerate}
\item Identify the complementary event: "at least one defective" becomes "\textbf{no} defective" (all good).
\item Count the complement with combinations: $n(\text{none}) = {}^{g}C_{r}$, where $g$ is the number of good items.
\item Compute $P(\text{at least one}) = 1 - \dfrac{n(\text{none})}{n(S)}$.
\end{enumerate}
\end{methode}

\begin{exoresolu}[At least one defective item]
A box contains 20 electronic components, of which 3 are defective. A sample of 4 components is selected at random for testing. Find the probability that the sample contains at least one defective component.
\tcblower
The total number of samples is
\[
n(S) = {}^{20}C_{4} = \frac{20 \times 19 \times 18 \times 17}{4 \times 3 \times 2 \times 1} = 4845.
\]
The complement of "at least one defective" is "no defective", i.e. all 4 chosen from the 17 good components:
\[
n(\text{no defective}) = {}^{17}C_{4} = \frac{17 \times 16 \times 15 \times 14}{4 \times 3 \times 2 \times 1} = 2380.
\]
Therefore
\[
P(\text{at least one defective}) = 1 - \frac{{}^{17}C_{4}}{{}^{20}C_{4}} = 1 - \frac{2380}{4845} = \frac{2465}{4845} = \frac{29}{57} \approx 0.509.
\]
Note how much shorter this is than adding the three cases "exactly 1, exactly 2, exactly 3 defective".
\end{exoresolu}

\courssub{Linking Counting Back to Tree Diagrams}

Counting formulas are not magic: on small cases they can be \emph{verified} with the tree diagrams of the previous section, and this cross-check is excellent examination practice.

\begin{experience}[Verifying ${}^{4}C_{2}$ with a tree]
A bag holds 2 red balls $R_1, R_2$ and 2 black balls $B_1, B_2$. Two balls are drawn without replacement. The combination formula predicts ${}^{4}C_{2} = 6$ possible (unordered) pairs: $R_1R_2,\ R_1B_1,\ R_1B_2,\ R_2B_1,\ R_2B_2,\ B_1B_2$. Drawing the full two-stage tree gives $4 \times 3 = 12$ ordered outcomes; grouping each pair of reversed duplicates ($R_1R_2$ with $R_2R_1$, etc.) recovers exactly the 6 unordered pairs, since $12 \div 2! = 6$. This is precisely why ${}^{n}C_{r} = {}^{n}P_{r}/r!$.

Moreover, the probability of drawing one ball of each colour can be computed two ways:
\begin{itemize}
\item by combinations: $\dfrac{{}^{2}C_{1}\,{}^{2}C_{1}}{{}^{4}C_{2}} = \dfrac{4}{6} = \dfrac{2}{3}$;
\item by the tree: $P(RB) + P(BR) = \dfrac{2}{4}\cdot\dfrac{2}{3} + \dfrac{2}{4}\cdot\dfrac{2}{3} = \dfrac{1}{3} + \dfrac{1}{3} = \dfrac{2}{3}$.
\end{itemize}
The agreement confirms that combinations are simply trees with the folding done algebraically.
\end{experience}

\begin{aretenir}[Key points of this section]
\begin{itemize}
\item Classical probability is a counting problem: $P(A) = \dfrac{n(A)}{n(S)}$; when listing is impossible, count by formula.
\item Multiplication principle: stages multiply, $n_1 \times n_2 \times \cdots \times n_k$.
\item Order matters $\Rightarrow$ permutations: ${}^{n}P_{r} = \dfrac{n!}{(n-r)!}$.
\item Order does not matter $\Rightarrow$ combinations: ${}^{n}C_{r} = \dfrac{n!}{r!(n-r)!}$.
\item Use the \emph{same} counting rule in the numerator and the denominator.
\item For "at least one", count the complement: $P = 1 - \dfrac{{}^{g}C_{r}}{{}^{n}C_{r}}$.
\item Simplify with ${}^{n}C_{r} = {}^{n}C_{n-r}$.
\end{itemize}
\end{aretenir}

\begin{exercice}[Practice]
\begin{enumerate}
\item A committee of 4 is chosen at random from 5 teachers and 7 students. Find the probability that the committee contains (a) exactly 2 teachers; (b) at least one teacher.
\item Six books are arranged at random on a shelf. Find the probability that two particular books are at the two ends.
\item A lot of 25 phones contains 4 defective ones. Three phones are tested at random. Find the probability that (a) none is defective; (b) at least one is defective.
\item From the digits $1,2,3,4,5,6$, a 4-digit code is formed without repetition. Find the probability that the code is even.
\end{enumerate}
\end{exercice}

\begin{corrige}[Exercise 1]
\begin{enumerate}
\item $n(S) = {}^{12}C_{4} = 495$.
(a) $P(\text{exactly 2 teachers}) = \dfrac{{}^{5}C_{2}\,{}^{7}C_{2}}{{}^{12}C_{4}} = \dfrac{10 \times 21}{495} = \dfrac{210}{495} = \dfrac{14}{33}$.
(b) $P(\text{at least one teacher}) = 1 - \dfrac{{}^{7}C_{4}}{{}^{12}C_{4}} = 1 - \dfrac{35}{495} = \dfrac{460}{495} = \dfrac{92}{99}$.
\item Total arrangements: $6! = 720$. Favourable: place the two books at the ends ($2!$ ways) and arrange the rest ($4!$ ways): $2! \times 4! = 48$. Hence $P = \dfrac{48}{720} = \dfrac{1}{15}$.
\item (a) $P(\text{none defective}) = \dfrac{{}^{21}C_{3}}{{}^{25}C_{3}} = \dfrac{1330}{2300} = \dfrac{133}{230}$. (b) $P(\text{at least one}) = 1 - \dfrac{133}{230} = \dfrac{97}{230}$.
\item Total codes: ${}^{6}P_{4} = 360$. Even codes end in 2, 4 or 6: $3$ choices for the last digit, then ${}^{5}P_{3} = 60$ for the rest, giving $180$. Hence $P = \dfrac{180}{360} = \dfrac{1}{2}$.
\end{enumerate}
\end{corrige}

\begin{savaistu}[Counting and the National Lottery]
Combination counting explains why jackpots are so hard to win. In a lottery where you must match 5 numbers out of 49, the number of possible tickets is ${}^{49}C_{5} = 1\,906\,884$: a single ticket has probability about $5.2 \times 10^{-7}$ of winning. The same mathematics, run in reverse, is used by quality-control engineers in factories — including those bottling drinks in Douala — to decide how large a sample must be so that "at least one defective" is detected with high probability.
\end{savaistu}

\coursec{Total Probability and Bayes' Theorem}

In the previous section we used permutations and combinations to count favourable outcomes when all outcomes are equally likely. But many real problems are not about counting: they are about \emph{combining} information that arrives in stages. A patient takes a test \emph{because} she might be ill; an item is defective \emph{because} it came from a particular machine; a student is late \emph{because} he chose a particular means of transport. In all these situations there is a hidden \cle{cause} (disease, machine, transport) and an observed \cle{result} (positive test, defective item, lateness). This final section gives the two master formulas for such problems: the \cle{law of total probability}, which gathers all the ways a result can happen, and \cle{Bayes' theorem}, which reverses the question: given the result, what is the probability of each cause?

\courssub{Partition of the Sample Space}

\begin{definition}[Partition of a sample space]
Events $B_1, B_2, \dots, B_n$ form a \cle{partition} of the sample space $S$ if:
\begin{enumerate}
\item they are \textbf{mutually exclusive}: $B_i \cap B_j = \emptyset$ for $i \neq j$;
\item they are \textbf{exhaustive}: $B_1 \cup B_2 \cup \dots \cup B_n = S$;
\item each has non-zero probability: $P(B_i) > 0$.
\end{enumerate}
In particular $P(B_1) + P(B_2) + \dots + P(B_n) = 1$.
\end{definition}

Intuitively, a partition is a complete, non-overlapping list of ``cases'': exactly one of the $B_i$ must occur. Typical examples: \{disease, no disease\}; \{item from machine 1, machine 2, machine 3\}; \{travel by bus, by taxi, on foot\}. The simplest partition of all is $\{B, B'\}$, an event and its complement — and indeed most examination questions use just two or three cases.

\begin{popfigure}
\begin{tikzpicture}[scale=1]
\fill[popblueL] (0,0) rectangle (2.7,4);
\fill[popgreenL] (2.7,0) rectangle (6.3,4);
\fill[popgold!40] (6.3,0) rectangle (9,4);
\draw[thick, popdark] (0,0) rectangle (9,4);
\draw[thick, popdark] (2.7,0) -- (2.7,4);
\draw[thick, popdark] (6.3,0) -- (6.3,4);
\node[popdark] at (0.35,3.7) {$S$};
\node at (1.35,3.3) {$B_1$};
\node at (4.5,3.3) {$B_2$};
\node at (7.65,3.3) {$B_3$};
\fill[poporange, opacity=0.55] (1.2,0.9) ellipse (2.6 and 1.0);
\fill[poporange, opacity=0.55] (6.9,1.1) ellipse (1.6 and 0.9);
\draw[thick, poporange] (1.2,0.9) ellipse (2.6 and 1.0);
\draw[thick, poporange] (6.9,1.1) ellipse (1.6 and 0.9);
\node[poporange] at (4.3,0.35) {$A$};
\end{tikzpicture}
\legende{A partition $B_1, B_2, B_3$ of $S$. The event $A$ crosses every strip: $A = (A\cap B_1)\cup(A\cap B_2)\cup(A\cap B_3)$, a disjoint union.}
\end{popfigure}

The key observation from the figure: whatever the event $A$, it can be split into disjoint pieces $A\cap B_1$, $A\cap B_2$, \dots, one piece inside each case. Since the pieces are mutually exclusive, their probabilities add.

\courssub{The Law of Total Probability}

\begin{propriete}[Law of total probability]
If $B_1, B_2, \dots, B_n$ form a partition of $S$, then for any event $A$,
\[ P(A) = P(A\mid B_1)P(B_1) + P(A\mid B_2)P(B_2) + \dots + P(A\mid B_n)P(B_n) = \sum_{i=1}^{n} P(A\mid B_i)\,P(B_i). \]
\emph{Proof (examinable and instructive).} Since the $B_i$ are exhaustive, $A = (A\cap B_1)\cup(A\cap B_2)\cup\dots\cup(A\cap B_n)$. Since the $B_i$ are mutually exclusive, so are the sets $A\cap B_i$, hence by the addition law $P(A) = \sum_i P(A\cap B_i)$. Now apply the multiplication law $P(A\cap B_i) = P(A\mid B_i)P(B_i)$ to each term. $\blacksquare$
\end{propriete}

On a two-stage tree (first the cause $B_i$, then the result $A$ or $A'$), the law of total probability says simply: \emph{add the probabilities of all paths that end at $A$}.

\begin{exemplebox}[Lateness and transport]
A student at a school in Bamenda travels by bus with probability $0.5$, by taxi with probability $0.3$ and on foot with probability $0.2$. The probabilities of arriving late are $0.1$, $0.05$ and $0.3$ respectively. Then
\[ P(\text{late}) = 0.1\times 0.5 + 0.05\times 0.3 + 0.3\times 0.2 = 0.05 + 0.015 + 0.06 = 0.125. \]
The three products are the three paths leading to ``late''; the sum is the total probability.
\end{exemplebox}

\courssub{Bayes' Theorem: Reversing the Conditioning}

Now the deeper question. The student \emph{is} late. What is the probability that he came on foot? The tree gives $P(\text{late}\mid\text{foot})$ directly, but we want $P(\text{foot}\mid\text{late})$ — the conditioning is reversed.

\begin{propriete}[Bayes' theorem]
If $B_1, \dots, B_n$ form a partition of $S$ and $P(A) > 0$, then for each $i$,
\[ P(B_i \mid A) = \frac{P(A \mid B_i)\,P(B_i)}{\displaystyle\sum_{j=1}^{n} P(A \mid B_j)\,P(B_j)} = \frac{P(A\cap B_i)}{P(A)}. \]
\emph{Proof.} By definition of conditional probability, $P(B_i\mid A) = \dfrac{P(A\cap B_i)}{P(A)}$. Write the numerator as $P(A\mid B_i)P(B_i)$ (multiplication law) and expand the denominator by the law of total probability. $\blacksquare$
\end{propriete}

The vocabulary is worth remembering: $P(B_i)$ is called the \cle{prior} probability (before the evidence), and $P(B_i\mid A)$ the \cle{posterior} probability (after observing $A$). Bayes' theorem is the rule that updates the prior into the posterior.

\begin{attention}[$P(A\mid B)$ is not $P(B\mid A)$]
``The probability of a positive test given the disease'' and ``the probability of the disease given a positive test'' are \textbf{different numbers}, often dramatically different. Confusing them is the single most common error in this topic — and in real medical reasoning. Bayes' theorem exists precisely to pass from one to the other.
\end{attention}

\begin{methode}[The two-stage tree method for Bayes problems]
\begin{enumerate}
\item Draw the tree in the \textbf{natural order}: first the cause $B_1, \dots, B_n$ (with their prior probabilities), then the result $A$ or $A'$ on each branch (conditional probabilities).
\item Compute every \textbf{path probability} by multiplying along the branches.
\item The denominator of Bayes' formula is the \textbf{sum of all paths ending at $A$} (total probability).
\item The answer is the \textbf{ratio}: (path through $B_i$ ending at $A$) $\div$ (sum of all paths ending at $A$).
\end{enumerate}
This is exactly how the GCE examiner expects the solution to be presented: a labelled tree, then a clear ratio.
\end{methode}

\courssub{Application 1: Medical Testing}

\begin{exoresolu}[A disease test]
A disease affects $1\%$ of a population. A test detects the disease with probability $0.95$ when it is present (\cle{sensitivity}), and gives a false positive with probability $0.04$ when the disease is absent. A person tests positive. What is the probability that he actually has the disease?
\tcblower
Let $D$ = ``has the disease'', $T$ = ``tests positive''. Then $P(D) = 0.01$, $P(D') = 0.99$, $P(T\mid D) = 0.95$, $P(T\mid D') = 0.04$.

\textbf{Step 1 — path probabilities.}
\begin{align*}
P(T\cap D) &= P(T\mid D)\,P(D) = 0.95 \times 0.01 = 0.0095,\\
P(T\cap D') &= P(T\mid D')\,P(D') = 0.04 \times 0.99 = 0.0396.
\end{align*}

\textbf{Step 2 — total probability of a positive test.}
\[ P(T) = 0.0095 + 0.0396 = 0.0491. \]

\textbf{Step 3 — Bayes' ratio.}
\[ P(D\mid T) = \frac{P(T\cap D)}{P(T)} = \frac{0.0095}{0.0491} \approx 0.193. \]
Despite testing positive, the person has only about a $19\%$ chance of having the disease!
\end{exoresolu}

\begin{popfigure}
\begin{tikzpicture}[grow=right, level distance=3.4cm,
level 1/.style={sibling distance=3.2cm},
level 2/.style={sibling distance=1.6cm},
every node/.style={font=\small, align=center},
edge from parent/.style={draw=popdark, thick}]
\node {Start}
child { node[align=center]{$D'$ \\ $0.99$}
  child { node[align=center]{$T'$ \\ $0.9504$} edge from parent node[below] {$0.96$} }
  child { node[poporange, align=center]{$T$ \\ $0.0396$} edge from parent node[above] {$0.04$} }
  edge from parent node[below] {$0.99$} }
child { node[align=center]{$D$ \\ $0.01$}
  child { node[align=center]{$T'$ \\ $0.0005$} edge from parent node[below] {$0.05$} }
  child { node[poporange, align=center]{$T$ \\ $0.0095$} edge from parent node[above] {$0.95$} }
  edge from parent node[above] {$0.01$} };
\end{tikzpicture}
\legende{Two-stage tree for the medical test. The two orange paths lead to a positive test $T$; Bayes' ratio is $\dfrac{0.0095}{0.0095+0.0396} \approx 0.193$.}
\end{popfigure}

\begin{attention}[The denominator must contain ALL paths to $A$]
A frequent exam error is to write $P(D\mid T) = \dfrac{0.0095}{0.01}$ or $\dfrac{0.0095}{0.95}$. The denominator is $P(T)$, the sum of \textbf{every} path ending at $T$ — including the false-positive path $0.0396$. Forgetting the false positives is exactly the mistake that makes people overestimate the meaning of a positive test.
\end{attention}

\begin{savaistu}[Why rare diseases give surprising results]
The result $P(D\mid T)\approx 0.19$ shocks most people: the test is ``$95\%$ accurate'', yet a positive result means only about a $19\%$ chance of illness. The reason is that the disease is \emph{rare}. Out of $10\,000$ people, about $100$ are ill (about $95$ test positive), while $9\,900$ are healthy (about $396$ false positives). So among the roughly $491$ positive results, only about $95$ are truly ill. When the prior probability $P(D)$ is tiny, even a small false-positive rate floods the positive group. This is why doctors confirm a positive screening test with a second, different test.
\end{savaistu}

\courssub{Application 2: Defective Items from Different Machines}

\begin{exoresolu}[Two machines]
A factory in Douala produces $60\%$ of its items on machine $M_1$ and $40\%$ on machine $M_2$. Machine $M_1$ produces $2\%$ defective items and $M_2$ produces $5\%$ defective items. An item chosen at random is found to be defective. Find the probability that it came from machine $M_2$.
\tcblower
Let $F$ = ``item is defective''. Given: $P(M_1) = 0.6$, $P(M_2) = 0.4$, $P(F\mid M_1) = 0.02$, $P(F\mid M_2) = 0.05$.

\textbf{Path probabilities to $F$:}
\[ P(F\cap M_1) = 0.02 \times 0.6 = 0.012, \qquad P(F\cap M_2) = 0.05 \times 0.4 = 0.020. \]

\textbf{Total probability:} $P(F) = 0.012 + 0.020 = 0.032$.

\textbf{Bayes' ratio:}
\[ P(M_2\mid F) = \frac{P(F\cap M_2)}{P(F)} = \frac{0.020}{0.032} = 0.625. \]
Although $M_2$ produces only $40\%$ of the items, it accounts for $62.5\%$ of the defective ones — because its defect rate is higher.
\end{exoresolu}

\courssub{Application 3: Reversing the Transport Problem}

Return to the Bamenda student: $P(\text{late}) = 0.125$, and the path through ``on foot'' has probability $0.3 \times 0.2 = 0.06$. Given that he is late,
\[ P(\text{foot}\mid\text{late}) = \frac{0.06}{0.125} = 0.48. \]
Before knowing anything, the probability that he walked was only $0.2$; the information ``he is late'' raises it to $0.48$. Bayes' theorem is precisely the mathematics of \emph{updating a probability in the light of new evidence}: the prior $P(B_i)$ becomes the posterior $P(B_i\mid A)$.

\begin{exoresolu}[Three causes — full Bayes with a three-branch partition]
Urn $U_1$ contains 3 red and 7 black balls, urn $U_2$ contains 6 red and 4 black balls, and urn $U_3$ contains 8 red and 2 black balls. An urn is chosen at random (each with probability $\tfrac13$) and a ball is drawn. The ball is red. Find the probability that it came from urn $U_3$.
\tcblower
Let $R$ = ``red ball''. The conditional probabilities are
\[ P(R\mid U_1) = \tfrac{3}{10} = 0.3, \qquad P(R\mid U_2) = \tfrac{6}{10} = 0.6, \qquad P(R\mid U_3) = \tfrac{8}{10} = 0.8. \]

\textbf{Paths to $R$:}
\[ P(R\cap U_1) = 0.3\times\tfrac13 = 0.1, \quad P(R\cap U_2) = 0.6\times\tfrac13 = 0.2, \quad P(R\cap U_3) = 0.8\times\tfrac13 = \tfrac{0.8}{3}. \]

\textbf{Total probability:}
\[ P(R) = \tfrac13(0.3 + 0.6 + 0.8) = \tfrac{1.7}{3}. \]

\textbf{Bayes' ratio:}
\[ P(U_3\mid R) = \frac{0.8/3}{1.7/3} = \frac{0.8}{1.7} = \frac{8}{17} \approx 0.471. \]
The prior probability of $U_3$ was $\tfrac13 \approx 0.333$; observing a red ball raises it to $\tfrac{8}{17} \approx 0.471$, because $U_3$ is the urn richest in red balls.
\end{exoresolu}

\courssub{Exam Technique and Consolidation}

\begin{methode}[GCE strategy for Bayes questions]
\begin{enumerate}
\item Define events clearly in words ($D$, $T$, $M_1$, \dots) and list every given probability with correct conditioning.
\item Draw the labelled two-stage tree and write the path probabilities at the tips.
\item For ``find $P(A)$'', add the relevant paths (total probability).
\item For ``given $A$, find $P(B_i\mid A)$'', write the ratio \emph{path}/\emph{sum of paths to $A$}, showing both numbers explicitly.
\item Give the final answer as a fraction or a decimal to 3 significant figures, and add one sentence of interpretation.
\end{enumerate}
\end{methode}

\begin{exercice}[]
A bag is chosen from three bags $B_1, B_2, B_3$ with probabilities $0.5, 0.3, 0.2$. The bags contain red balls in proportions $0.4, 0.5, 0.75$ respectively. A ball is drawn and is red. Find the probability that it came from bag $B_3$.
\end{exercice}

\begin{corrige}[Exercise 1]
Let $R$ = ``red ball''. Paths to $R$: $0.4\times0.5 = 0.20$; $0.5\times0.3 = 0.15$; $0.75\times0.2 = 0.15$. Total: $P(R) = 0.20+0.15+0.15 = 0.50$. Hence
\[ P(B_3\mid R) = \frac{0.15}{0.50} = 0.3. \]
\end{corrige}

\begin{exercice}[]
In a school, $55\%$ of the candidates study science and $45\%$ study arts. The probability of passing the GCE is $0.8$ for a science candidate and $0.7$ for an arts candidate. (a) Find the probability that a randomly chosen candidate passes. (b) A candidate passed; find the probability that she is a science candidate.
\end{exercice}

\begin{corrige}[Exercise 2]
Let $S$ = ``science'', $A$ = ``arts'', $P$ = ``passes''. (a) Total probability:
\[ P(P) = 0.8\times0.55 + 0.7\times0.45 = 0.44 + 0.315 = 0.755. \]
(b) Bayes' ratio:
\[ P(S\mid P) = \frac{0.44}{0.755} \approx 0.583. \]
\end{corrige}

\begin{aretenir}[Key points of the section]
\begin{itemize}
\item A \cle{partition} splits $S$ into mutually exclusive, exhaustive cases $B_1, \dots, B_n$.
\item \cle{Total probability}: $P(A) = \sum_i P(A\mid B_i)P(B_i)$ — add all paths leading to $A$.
\item \cle{Bayes' theorem}: $P(B_i\mid A) = \dfrac{P(A\mid B_i)P(B_i)}{\sum_j P(A\mid B_j)P(B_j)}$ — the ratio of one path to all paths through $A$.
\item $P(A\mid B) \neq P(B\mid A)$; Bayes reverses the conditioning, updating the \cle{prior} into the \cle{posterior}.
\item With a rare cause, even a very accurate test gives mostly false positives: always compare the true-positive path with \emph{all} positive paths.
\end{itemize}
\end{aretenir}

\newpage

\coursec{Integration activity}

\coursec{Activities of Integration, Assessment and Remediation}

\courssub{Aims of the activity}

This integration activity mobilises, in one authentic investigation, all the skills developed in this chapter:

\begin{itemize}
\item identifying a random experiment, its outcomes and its sample space;
\item distinguishing events and computing \cle{classical probabilities} using \cle{combinations};
\item constructing and reading a \cle{tree diagram} with conditional probabilities;
\item applying the \cle{multiplication law}, the \cle{law of total probability} and \cle{Bayes' theorem};
\item using the \cle{complement rule} for ``at least one'' problems;
\item distinguishing \cle{selection with replacement} from \cle{selection without replacement};
\item communicating a probabilistic conclusion in clear, non-technical language for a decision maker.
\end{itemize}

By the end of the activity, each group should be able to convert a real-life security problem into a probability model, compute the relevant probabilities rigorously, and defend a practical recommendation before the class.

\courssub{Situation}

\textbf{The Mobile Money Fraud Investigation.}\par

Madam Neba runs a busy mobile money kiosk at Commercial Avenue in Bamenda. Every day she processes hundreds of deposits and withdrawals. Recently, her operator's fraud-monitoring unit sent her a worrying report. After a careful audit of transaction records across her region, the unit established the following facts, which we shall take as exact for this investigation:

\begin{itemize}
\item fraudulent withdrawal attempts occur on $2\%$ of all transactions;
\item the operator's automatic security system \emph{flags} (marks as suspicious) $95\%$ of all fraudulent transactions;
\item unfortunately, the same system also \emph{wrongly flags} $3\%$ of all genuine transactions.
\end{itemize}

Whenever a transaction is flagged, Madam Neba must decide what to do. Blocking a genuine customer's transaction causes anger and loss of business; letting a fraudulent transaction through causes direct financial loss. She asks your class, as trainee data analysts, to answer the following questions.

\begin{enumerate}
\item If a transaction is flagged by the system, what is the probability that it is \emph{actually} fraudulent? Madam Neba guesses ``about $95\%$, since the system catches $95\%$ of frauds''. Is she right?
\item On a typical morning, the system flags $5$ transactions. Assuming these flagged transactions behave independently, what is the probability that \emph{at least one} of them is actually fraudulent?
\item As a final check, an auditor examines a batch of $20$ transactions known to contain exactly $2$ fraudulent ones. The auditor selects $4$ transactions at random (without replacement) for a detailed manual review. What is the probability that the review detects at least one of the two frauds?
\end{enumerate}

Finally, each group must write a short \textbf{advisory note} (at most ten lines) to Madam Neba recommending whether flagged transactions should be \emph{blocked automatically} or \emph{verified manually first}, justifying the recommendation with the computed probabilities.

\courssub{Tasks}

Work in groups of three or four. Show all reasoning; a numerical answer without justification earns little credit.

\begin{enumerate}
\item \textbf{Modelling.} Define the experiment and the two events $F$ = ``the transaction is fraudulent'' and $S$ = ``the transaction is flagged by the system''. Write down, in probability notation, the three numerical data of the situation. What is the sample space if we classify one transaction according to both criteria?
\item \textbf{Tree diagram.} Draw a two-stage tree diagram (first stage: fraudulent or genuine; second stage: flagged or not flagged). Label every branch with its probability and compute the probability of each of the four terminal outcomes.
\item \textbf{Total probability.} Using the tree, compute $P(S)$, the probability that a transaction is flagged.
\item \textbf{Bayes' theorem.} Compute $P(F \mid S)$, the probability that a flagged transaction is actually fraudulent. Compare with Madam Neba's guess of $95\%$ and explain, in two or three sentences, why the true value is so much lower even though the system catches $95\%$ of frauds.
\item \textbf{Several flagged transactions.} Assuming independence, compute the probability that at least one of $5$ flagged transactions is fraudulent. (Hint: use the complement rule.)
\item \textbf{Counting techniques.} For the auditor's sample of $4$ transactions chosen from $20$ containing $2$ frauds, compute:
\begin{enumerate}
\item the total number of possible samples;
\item the number of samples containing \emph{no} fraud;
\item the probability that the sample contains at least one fraud.
\end{enumerate}
Explain clearly why combinations (and not permutations) are appropriate here, and why this is selection \emph{without} replacement.
\item \textbf{Advisory note.} Write your recommendation to Madam Neba. It must mention the value of $P(F \mid S)$, the consequence of blocking genuine customers, and a concrete proposal (for example: manual verification, a second confirmation step, or a phone call to the customer).
\item \textbf{Plenary.} Each group presents its tree, its value of $P(F \mid S)$ and its recommendation in three minutes. Compare the recommendations across groups.
\end{enumerate}

\begin{attention}[The two classic confusions]
\begin{itemize}
\item $P(S \mid F) = 0.95$ is \textbf{not} the same as $P(F \mid S)$. The first is given; the second is what Bayes' theorem computes. Confusing them is the most frequent error in this type of problem.
\item The auditor samples \emph{without} replacement: use combinations, not powers. Writing $\left(\tfrac{18}{20}\right)^4$ would only be correct if each transaction were replaced before the next draw.
\end{itemize}
\end{attention}

\courssub{Model answer}

\textbf{1. Modelling.} The experiment is: ``observe one transaction and record whether it is fraudulent and whether it is flagged''. Let $F$ = ``the transaction is fraudulent'' and $S$ = ``the transaction is flagged''. The data are
\[ P(F) = 0.02, \qquad P(S \mid F) = 0.95, \qquad P(S \mid \bar{F}) = 0.03, \]
where $\bar{F}$ = ``the transaction is genuine'', with $P(\bar{F}) = 1 - 0.02 = 0.98$. The sample space, classifying by both criteria, is
\[ \Omega = \{F \cap S,\; F \cap \bar{S},\; \bar{F} \cap S,\; \bar{F} \cap \bar{S}\}. \]

\textbf{2. Tree diagram.}

\begin{popfigure}
\begin{tikzpicture}[grow=right, level distance=3.2cm,
level 1/.style={sibling distance=2.6cm},
level 2/.style={sibling distance=1.3cm},
every node/.style={font=\small}]
\node {Transaction}
child { node {$\bar{F}$} 
  child { node {$\bar{S}$} edge from parent node[below] {$0.97$} }
  child { node {$S$} edge from parent node[above] {$0.03$} }
  edge from parent node[below] {$0.98$} }
child { node {$F$} 
  child { node {$\bar{S}$} edge from parent node[below] {$0.05$} }
  child { node {$S$} edge from parent node[above] {$0.95$} }
  edge from parent node[above] {$0.02$} };
\end{tikzpicture}
\legende{Two-stage tree for the fraud investigation}
\end{popfigure}

The four terminal probabilities (multiplication law) are:
\begin{align*}
P(F \cap S) &= 0.02 \times 0.95 = 0.019,\\
P(F \cap \bar{S}) &= 0.02 \times 0.05 = 0.001,\\
P(\bar{F} \cap S) &= 0.98 \times 0.03 = 0.0294,\\
P(\bar{F} \cap \bar{S}) &= 0.98 \times 0.97 = 0.9506.
\end{align*}
Check: $0.019 + 0.001 + 0.0294 + 0.9506 = 1$. \checkmark

\textbf{3. Total probability.} A transaction is flagged either because it is fraudulent and caught, or genuine and wrongly flagged:
\[ P(S) = P(F \cap S) + P(\bar{F} \cap S) = 0.019 + 0.0294 = 0.0484. \]
So about $4.84\%$ of all transactions are flagged.

\textbf{4. Bayes' theorem.}
\[ P(F \mid S) = \frac{P(F \cap S)}{P(S)} = \frac{P(S \mid F)\,P(F)}{P(S)} = \frac{0.95 \times 0.02}{0.0484} = \frac{0.019}{0.0484} \approx 0.3926. \]
Only about $39.3\%$ of flagged transactions are actually fraudulent — far below Madam Neba's guess of $95\%$. The reason is that fraud is \emph{rare} ($2\%$): the genuine transactions are so numerous ($98\%$) that even a small false-alarm rate of $3\%$ produces $0.0294$ worth of false alarms, which \emph{outnumber} the $0.019$ of true alarms. In other words, most flagged transactions are innocent.

\textbf{5. At least one fraud among 5 flagged transactions.} For a flagged transaction, $P(\bar{F} \mid S) = 1 - 0.3926 = 0.6074$. Assuming independence,
\[ P(\text{at least one fraud}) = 1 - (0.6074)^5 \approx 1 - 0.0827 \approx 0.9173. \]
There is about a $91.7\%$ chance that at least one of the five flagged transactions is genuinely fraudulent — so the flags must be taken seriously even though each individual flag is more likely than not to be a false alarm.

\textbf{6. Auditor's sample (without replacement).}
\begin{enumerate}
\item Total number of samples of $4$ from $20$:
\[ \binom{20}{4} = \frac{20 \times 19 \times 18 \times 17}{4 \times 3 \times 2 \times 1} = 4845. \]
\item Samples containing no fraud: choose all $4$ from the $18$ genuine ones:
\[ \binom{18}{4} = \frac{18 \times 17 \times 16 \times 15}{24} = 3060. \]
\item By the complement rule:
\[ P(\text{at least one fraud}) = 1 - \frac{3060}{4845} = \frac{1785}{4845} \approx 0.3684. \]
\end{enumerate}
Combinations are appropriate because the \emph{order} in which the auditor picks the transactions is irrelevant — only the \emph{set} selected matters. The draws are without replacement since a transaction cannot be reviewed twice. (With replacement, the answer would have been $1 - (0.9)^4 = 0.3439$, a different value.)

\textbf{7. Advisory note (example).} \emph{``Dear Madam Neba, our analysis shows that only about $39\%$ of flagged transactions are actually fraudulent: more than $6$ out of $10$ flags are false alarms on honest customers. We therefore advise you \textbf{not} to block flagged transactions automatically. Instead, verify each flag manually (call the customer, check identity) before blocking. This is worthwhile because among any $5$ flagged transactions there is about a $92\%$ chance that at least one is a real fraud.''}

\begin{attention}[Remediation targets]
If your group wrote $P(F \mid S) = 0.95$, revisit Task 4: you confused a conditional probability with its converse. If in Task 6 you used powers of $\tfrac{18}{20}$, revisit the with/without replacement distinction: sampling several items from a finite batch at once is always without replacement.
\end{attention}

\newpage

\coursec{Methods and worked examples}

\courssub{Skill 1: Computing a Classical (Theoretical) Probability}

\begin{methode}[Calculating $P(A)=\dfrac{n(A)}{n(S)}$]
\begin{enumerate}
\item Identify the \cle{experiment} and check that all outcomes are \cle{equally likely} (fair coin, fair die, well-shuffled cards, random draw).
\item List or count the \cle{sample space} $S$; write $n(S)$.
\item List or count the outcomes favourable to the event $A$; write $n(A)$.
\item Compute $P(A)=\dfrac{n(A)}{n(S)}$ and simplify the fraction. Give the answer as a fraction, decimal or percentage as requested.
\item Check: the answer must satisfy $0\le P(A)\le 1$. If not, recount.
\end{enumerate}
\textbf{Reflexes:} a standard die has $6$ faces; a pack of cards has $52$ cards ($4$ suits, $13$ ranks); when two dice are thrown $n(S)=36$ and a table of sums avoids omissions.
\end{methode}

\begin{attention}[Equally likely outcomes]
The formula $P(A)=\dfrac{n(A)}{n(S)}$ is valid \textbf{only} when all outcomes of the sample space are equally likely. For two dice, the $36$ ordered pairs are equally likely, but the $11$ possible sums are \emph{not}: a sum of $2$ occurs once while a sum of $7$ occurs six times. Never apply the formula to a list of non-equiprobable outcomes.
\end{attention}

\begin{exoresolu}[Two dice at a school fair]
At a fund-raising game in a college in Bamenda, two fair dice are thrown once. Find the probability that (a) the sum of the scores is $9$; (b) the sum is at least $10$; (c) the two dice show the same number.
\tcblower
The experiment is throwing two fair dice; all $36$ ordered pairs are equally likely, so $n(S)=6\times 6=36$.

\textbf{(a) Sum equal to 9.} The favourable outcomes are
\[(3,6),\ (4,5),\ (5,4),\ (6,3),\quad\text{so } n(A)=4.\]
Hence $P(\text{sum}=9)=\dfrac{4}{36}=\dfrac{1}{9}$.

\textbf{(b) Sum at least 10.} This means sum $=10, 11$ or $12$:
\begin{align*}
\text{sum }10 &: (4,6),(5,5),(6,4) \quad (3 \text{ outcomes})\\
\text{sum }11 &: (5,6),(6,5) \quad (2 \text{ outcomes})\\
\text{sum }12 &: (6,6) \quad (1 \text{ outcome})
\end{align*}
So $n(B)=3+2+1=6$ and $P(B)=\dfrac{6}{36}=\dfrac{1}{6}$.

\textbf{(c) Same number on both dice.} The doubles are $(1,1),(2,2),\dots,(6,6)$, giving $n(C)=6$, so
\[P(C)=\frac{6}{36}=\frac{1}{6}.\]
All three answers lie between $0$ and $1$, as required.
\end{exoresolu}

\courssub{Skill 2: Estimating an Empirical Probability}

\begin{methode}[Using relative frequency]
\begin{enumerate}
\item Identify the total number of \cle{trials} $N$ (throws, observations, sampled items).
\item Count the number $f$ of trials in which the event occurred.
\item Compute the \cle{relative frequency} $P(A)\approx \dfrac{f}{N}$.
\item Interpret: the larger $N$ is, the more reliable the estimate (law of large numbers).
\item Use the estimate to predict: expected number of occurrences in $n$ future trials $\approx n\times \dfrac{f}{N}$.
\end{enumerate}
\textbf{Reflex:} empirical probability comes from \emph{data}; theoretical probability comes from \emph{reasoning} about equally likely outcomes.
\end{methode}

\begin{exoresolu}[Defective loaves in a bakery]
A bakery in Douala inspected $500$ loaves and found that $35$ were underweight. (a) Estimate the probability that a randomly chosen loaf is underweight. (b) On a day when $2\,000$ loaves are produced, how many underweight loaves should be expected? (c) Is this a theoretical or an empirical probability? Justify.
\tcblower
\textbf{(a)} The number of trials is $N=500$ and the frequency of ``underweight'' is $f=35$. Hence
\[P(\text{underweight})\approx \frac{f}{N}=\frac{35}{500}=0.07.\]

\textbf{(b)} Expected number $\approx n\times 0.07 = 2\,000\times 0.07 = 140$ loaves.

\textbf{(c)} This is an \cle{empirical} (experimental) probability: it is obtained from observed data, not from a symmetry argument. The loaves are not ``equally likely outcomes'' of a known sample space; the estimate $0.07$ could change if a larger sample were taken, and it becomes more reliable as $N$ grows.
\end{exoresolu}

\courssub{Skill 3: Applying the Addition Laws}

\begin{methode}[Mutually exclusive or not?]
\begin{enumerate}
\item Define the events $A$ and $B$ clearly in words.
\item Ask: \emph{can $A$ and $B$ happen at the same time?} If no, they are \cle{mutually exclusive} and $P(A\cap B)=0$.
\item If mutually exclusive: $P(A\cup B)=P(A)+P(B)$.
\item If not: use the general law $P(A\cup B)=P(A)+P(B)-P(A\cap B)$.
\item For ``neither'', use the complement: $P(\text{neither})=1-P(A\cup B)$.
\end{enumerate}
\textbf{Reflex:} never add probabilities blindly; subtract the overlap or you count it twice.
\end{methode}

\begin{exoresolu}[Drawing a card from a pack]
A card is drawn at random from a standard pack of $52$ cards. Find the probability that it is (a) a king or a queen; (b) a heart or a king; (c) neither a heart nor a king.
\tcblower
Each card is equally likely, so probabilities are counted out of $52$.

\textbf{(a)} Let $K=$ ``king'', $Q=$ ``queen''. A card cannot be both, so the events are mutually exclusive:
\[P(K\cup Q)=P(K)+P(Q)=\frac{4}{52}+\frac{4}{52}=\frac{8}{52}=\frac{2}{13}.\]

\textbf{(b)} Let $H=$ ``heart''. A card \emph{can} be both a heart and a king (the king of hearts), so the events are \textbf{not} mutually exclusive. We have $P(H)=\frac{13}{52}$, $P(K)=\frac{4}{52}$ and $P(H\cap K)=\frac{1}{52}$. By the general addition law:
\[P(H\cup K)=\frac{13}{52}+\frac{4}{52}-\frac{1}{52}=\frac{16}{52}=\frac{4}{13}.\]

\textbf{(c)} ``Neither a heart nor a king'' is the complement of $H\cup K$:
\[P(\text{neither})=1-P(H\cup K)=1-\frac{4}{13}=\frac{9}{13}.\]
Check: there are $13$ hearts plus $3$ non-heart kings $=16$ ``bad'' cards, leaving $52-16=36$ good cards, and $\frac{36}{52}=\frac{9}{13}$. $\checkmark$
\end{exoresolu}

\courssub{Skill 4: Multiplication Law, Conditional Probability and Independence}

\begin{methode}[Using $P(A\cap B)=P(A)\,P(B\mid A)$]
\begin{enumerate}
\item Translate the question: ``probability of $B$ \emph{given that} $A$ has occurred'' means $P(B\mid A)=\dfrac{P(A\cap B)}{P(A)}$, provided $P(A)\ne 0$.
\item For successive draws, update the sample space after each draw (the composition of the bag/pack changes if there is no replacement).
\item Multiply along the chain: $P(A\cap B)=P(A)\times P(B\mid A)$.
\item To test \cle{independence}, check whether $P(B\mid A)=P(B)$, equivalently $P(A\cap B)=P(A)P(B)$.
\item For \cle{exhaustive} events, check that the probabilities sum to $1$.
\end{enumerate}
\textbf{Reflex:} ``given that'' $\Rightarrow$ conditional probability; ``and'' $\Rightarrow$ intersection (multiply); ``or'' $\Rightarrow$ union (add carefully).
\end{methode}

\begin{attention}[Mutually exclusive is not independent]
Do not confuse the two notions. \cle{Mutually exclusive} means $A\cap B=\emptyset$ (they cannot occur together). \cle{Independent} means $P(A\cap B)=P(A)P(B)$ (knowing one occurred does not change the probability of the other). Two non-impossible mutually exclusive events are \emph{never} independent, since $P(A\cap B)=0\ne P(A)P(B)$.
\end{attention}

\begin{exoresolu}[Students studying French and Computer Science]
In an Upper Sixth class of $40$ students, $25$ study French, $18$ study Computer Science and $10$ study both. A student is chosen at random. (a) Find the probability that the student studies Computer Science, given that the student studies French. (b) Are the events ``studies French'' and ``studies Computer Science'' independent? (c) Show whether these two events are exhaustive.
\tcblower
Let $F=$ ``studies French'' and $C=$ ``studies Computer Science''. Then
\[P(F)=\frac{25}{40}=\frac58,\qquad P(C)=\frac{18}{40}=\frac{9}{20},\qquad P(F\cap C)=\frac{10}{40}=\frac14.\]

\textbf{(a)} By the definition of conditional probability,
\[P(C\mid F)=\frac{P(C\cap F)}{P(F)}=\frac{\frac14}{\frac58}=\frac14\times\frac85=\frac{2}{5}=0.4.\]
Interpretation: among the $25$ French students, $10$ also study Computer Science, and $\frac{10}{25}=\frac25$. $\checkmark$

\textbf{(b)} Test independence: $P(C\mid F)=0.4$ while $P(C)=\frac{9}{20}=0.45$. Since $P(C\mid F)\ne P(C)$, the events are \textbf{not independent}. (Equivalently, $P(F)P(C)=\frac58\times\frac{9}{20}=\frac{9}{32}\ne\frac14=P(F\cap C)$.)

\textbf{(c)} Exhaustive events must cover the whole sample space, i.e.\ their union must have probability $1$:
\[P(F\cup C)=\frac58+\frac{9}{20}-\frac14=\frac{25+18-10}{40}=\frac{33}{40}\ne 1.\]
Since $7$ students study neither subject, $F$ and $C$ are \textbf{not exhaustive}.
\end{exoresolu}

\courssub{Skill 5: Tree Diagrams and Selection With or Without Replacement}

\begin{methode}[Solving with a tree diagram]
\begin{enumerate}
\item Identify the \cle{stages} of the experiment (first draw, second draw, \dots).
\item Draw branches for all possible outcomes at each stage and write the probability on each branch. Probabilities leaving any node must sum to $1$.
\item \textbf{With replacement:} branch probabilities stay the same at every stage. \textbf{Without replacement:} reduce both the number of favourable items and the total by $1$ after each draw.
\item Multiply probabilities \emph{along} a path to get the probability of that path.
\item Add the probabilities of all paths corresponding to the required event.
\item Check: the probabilities of \emph{all} complete paths must sum to $1$.
\end{enumerate}
\end{methode}

\begin{exoresolu}[Balls drawn from a bag]
A bag contains $5$ red and $3$ blue balls. Two balls are drawn at random, one after the other, \textbf{without replacement}. (a) Draw a tree diagram for the experiment. (b) Find the probability that both balls are red. (c) Find the probability that the two balls are of different colours. (d) Repeat (b) if the first ball \emph{is} replaced before the second draw.
\tcblower
\textbf{(a)} At the first draw: $P(R)=\frac58$, $P(B)=\frac38$. Since there is no replacement, after a red first draw the bag has $4$ red and $3$ blue out of $7$; after a blue first draw it has $5$ red and $2$ blue out of $7$.

\begin{popfigure}
\begin{tikzpicture}[grow=right, level distance=3.2cm,
level 1/.style={sibling distance=3cm},
level 2/.style={sibling distance=1.5cm},
every node/.style={font=\small}]
\node {}
child { node {$B$}
  child { node {$BB:\ \frac{3}{28}$} edge from parent node[above, sloped] {$\frac{2}{7}$} }
  child { node {$BR:\ \frac{15}{56}$} edge from parent node[above, sloped] {$\frac{5}{7}$} }
  edge from parent node[above, sloped] {$\frac{3}{8}$}
}
child { node {$R$}
  child { node {$RB:\ \frac{15}{56}$} edge from parent node[above, sloped] {$\frac{3}{7}$} }
  child { node {$RR:\ \frac{5}{14}$} edge from parent node[above, sloped] {$\frac{4}{7}$} }
  edge from parent node[above, sloped] {$\frac{5}{8}$}
};
\end{tikzpicture}
\legende{Tree diagram for two draws without replacement; path probabilities are obtained by multiplying along the branches.}
\end{popfigure}

\textbf{(b)} Path $R$ then $R$: multiply along the branches:
\[P(RR)=\frac58\times\frac47=\frac{20}{56}=\frac{5}{14}.\]

\textbf{(c)} Different colours means $RB$ or $BR$:
\begin{align*}
P(RB)&=\frac58\times\frac37=\frac{15}{56},\\
P(BR)&=\frac38\times\frac57=\frac{15}{56},\\
P(\text{different})&=\frac{15}{56}+\frac{15}{56}=\frac{30}{56}=\frac{15}{28}.
\end{align*}
Check of the whole tree: $\frac{20}{56}+\frac{15}{56}+\frac{15}{56}+\frac{6}{56}=\frac{56}{56}=1$. $\checkmark$

\textbf{(d)} With replacement, the bag is unchanged for the second draw, so the draws are independent:
\[P(RR)=\frac58\times\frac58=\frac{25}{64}.\]
Note that $\frac{25}{64}\approx 0.391 > \frac{5}{14}\approx 0.357$: replacement makes a repeated colour slightly more likely here.
\end{exoresolu}

\courssub{Skill 6: Using Permutations and Combinations in Probability}

\begin{methode}[Counting with ${}^nC_r$ and ${}^nP_r$]
\begin{enumerate}
\item Decide whether \cle{order matters}. Committee, hand of cards, selection $\Rightarrow$ order does not matter: use ${}^nC_r=\dfrac{n!}{r!(n-r)!}$. Arrangement, ranking, podium $\Rightarrow$ order matters: use ${}^nP_r=\dfrac{n!}{(n-r)!}$.
\item Count $n(S)$, the total number of possible selections/arrangements, using the same rule as for $n(A)$.
\item Count $n(A)$ by splitting the population into categories and multiplying the choices for each category.
\item Compute $P(A)=\dfrac{n(A)}{n(S)}$ and simplify.
\item For ``at least one'' questions, it is often faster to use the complement: $P(\text{at least one})=1-P(\text{none})$.
\end{enumerate}
\textbf{Reflex:} numerator and denominator must count the \emph{same kind} of object (both unordered or both ordered).
\end{methode}

\begin{exoresolu}[Choosing a delegation]
A school in Yaound\'e must send a delegation of $4$ students chosen at random from $7$ boys and $5$ girls. Find the probability that the delegation contains (a) exactly $2$ girls; (b) at least one girl; (c) all four students of the same sex.
\tcblower
There are $12$ students and the delegation is an \emph{unordered} selection, so
\[n(S)={}^{12}C_4=\frac{12!}{4!\,8!}=\frac{12\times11\times10\times9}{4\times3\times2\times1}=495.\]

\textbf{(a) Exactly 2 girls} means $2$ girls from $5$ \emph{and} $2$ boys from $7$:
\[n(A)={}^{5}C_2\times{}^{7}C_2=10\times21=210,\qquad
P(A)=\frac{210}{495}=\frac{14}{33}.\]

\textbf{(b) At least one girl.} Use the complement: ``no girl'' means $4$ boys from $7$:
\[P(\text{no girl})=\frac{{}^{7}C_4}{{}^{12}C_4}=\frac{35}{495}=\frac{7}{99},\]
\[P(\text{at least one girl})=1-\frac{7}{99}=\frac{92}{99}.\]

\textbf{(c) All of the same sex} means $4$ girls or $4$ boys (mutually exclusive):
\[n={}^{5}C_4+{}^{7}C_4=5+35=40,\qquad P=\frac{40}{495}=\frac{8}{99}.\]
Each answer lies in $[0,1]$; the direct counting in (c) confirms the complement figure used in (b).
\end{exoresolu}

\courssub{Skill 7: Total Probability and Bayes' Theorem}

\begin{methode}[Applying Bayes' theorem]
\begin{enumerate}
\item Identify the \cle{partition} of the sample space: events $B_1,B_2,\dots,B_n$ that are mutually exclusive and exhaustive (e.g.\ which box, which factory, which route).
\item Identify the event $A$ that is observed (e.g.\ ``defective item'', ``positive test'').
\item Write down $P(B_i)$ and the conditional probabilities $P(A\mid B_i)$.
\item Compute the \cle{total probability}: $P(A)=\displaystyle\sum_i P(B_i)\,P(A\mid B_i)$.
\item Reverse the conditioning with \cle{Bayes' theorem}:
\[P(B_k\mid A)=\frac{P(B_k)\,P(A\mid B_k)}{\sum_i P(B_i)\,P(A\mid B_i)}.\]
\item A tree diagram with the $B_i$ at the first stage makes the whole computation visual and safe.
\end{enumerate}
\end{methode}

\begin{exoresolu}[Two factories supplying a market]
A market in Buea is supplied with bottled groundnut oil by two factories $X$ and $Y$. Factory $X$ supplies $60\%$ of the bottles and $Y$ supplies $40\%$. It is known that $2\%$ of the bottles from $X$ and $5\%$ of those from $Y$ are defective. A bottle is bought at random. (a) Find the probability that it is defective. (b) Given that the bottle is defective, find the probability that it came from factory $Y$.
\tcblower
Let $D=$ ``bottle is defective''. The events $X$ and $Y$ form a partition: $P(X)=0.6$, $P(Y)=0.4$, with $P(D\mid X)=0.02$ and $P(D\mid Y)=0.05$.

\textbf{(a)} By the law of total probability,
\begin{align*}
P(D)&=P(X)P(D\mid X)+P(Y)P(D\mid Y)\\
&=(0.6)(0.02)+(0.4)(0.05)\\
&=0.012+0.020=0.032.
\end{align*}
So $3.2\%$ of all bottles on the market are defective.

\textbf{(b)} By Bayes' theorem,
\[P(Y\mid D)=\frac{P(Y)\,P(D\mid Y)}{P(D)}=\frac{0.020}{0.032}=\frac{20}{32}=\frac58=0.625.\]

\textbf{Interpretation.} Although factory $Y$ supplies only $40\%$ of the bottles, it produces $\frac58=62.5\%$ of the \emph{defective} ones, because its defect rate is higher. Note that $P(Y\mid D)\ne P(D\mid Y)$: confusing these two conditional probabilities is the classic Bayes trap. Check: $P(X\mid D)=\frac{0.012}{0.032}=0.375$, and $0.375+0.625=1$. $\checkmark$
\end{exoresolu}

\newpage

\coursec{Exercises}

\courssub{I check my knowledge}

\begin{exercice}[Multiple choice questions]
For each question, choose the single correct answer.
\begin{enumerate}
\item A fair die is thrown once. The probability of obtaining a prime number is:
\begin{enumerate}
\item[(a)] $\dfrac{1}{6}$ \quad (b) $\dfrac{1}{3}$ \quad (c) $\dfrac{1}{2}$ \quad (d) $\dfrac{2}{3}$
\end{enumerate}
\item Two events $A$ and $B$ are mutually exclusive, with $P(A)=0.3$ and $P(B)=0.5$. Then $P(A\cup B)$ equals:
\begin{enumerate}
\item[(a)] $0.15$ \quad (b) $0.8$ \quad (c) $0.65$ \quad (d) $0.2$
\end{enumerate}
\item If $A$ and $B$ are independent events with $P(A)=0.4$ and $P(B)=0.25$, then $P(A\cap B)$ equals:
\begin{enumerate}
\item[(a)] $0.65$ \quad (b) $0.1$ \quad (c) $0.15$ \quad (d) $0.55$
\end{enumerate}
\item A card is drawn from a pack of 52. The probability that it is a king or a heart is:
\begin{enumerate}
\item[(a)] $\dfrac{17}{52}$ \quad (b) $\dfrac{15}{52}$ \quad (c) $\dfrac{4}{13}$ \quad (d) $\dfrac{1}{13}$
\end{enumerate}
\end{enumerate}
\end{exercice}

\begin{exercice}[True or false]
State whether each statement is true or false, justifying your answer with a short argument or a counter-example.
\begin{enumerate}
\item If two events are mutually exclusive, then they are independent.
\item For any event $A$, $P(A)+P(A')=1$.
\item If $P(A\mid B)=P(A)$, then $A$ and $B$ are independent.
\item The empirical probability of an event always equals its theoretical probability.
\end{enumerate}
\end{exercice}

\begin{exercice}[Definitions]
Define each of the following terms precisely, giving one example of each:
\begin{enumerate}
\item a random experiment and its sample space;
\item mutually exclusive events;
\item exhaustive events;
\item independent events;
\item the conditional probability $P(A\mid B)$.
\end{enumerate}
\end{exercice}

\begin{exercice}[Direct calculations]
\begin{enumerate}
\item A letter is chosen at random from the word \textbf{CAMEROON}. Find the probability that it is a vowel.
\item Two fair coins are tossed. Write down the sample space and find the probability of obtaining exactly one head.
\item Given $P(A)=0.6$, $P(B)=0.5$ and $P(A\cap B)=0.2$, find $P(A\cup B)$ and $P(A\mid B)$.
\item A fair coin is tossed 200 times and lands heads 118 times. State the empirical probability of heads, compare it with the theoretical probability, and explain what you expect as the number of tosses increases.
\end{enumerate}
\end{exercice}

\courssub{I practise}

\begin{exercice}[Balls of the same colour]
A bag contains 6 red and 4 black balls. Two balls are drawn at random. Find the probability that they are of the same colour:
\begin{enumerate}
\item when the first ball is replaced before the second draw;
\item when the first ball is not replaced.
\end{enumerate}
Compare the two answers and comment.
\end{exercice}

\begin{exercice}[Physics and Chemistry students]
In a class of 40 students, 18 study Physics, 15 study Chemistry and 7 study both subjects. A student is chosen at random from the class.
\begin{enumerate}
\item Illustrate the situation with a Venn diagram.
\item Find the probability that the student studies Physics or Chemistry.
\item Find the probability that the student studies neither subject.
\item Given that the chosen student studies Physics, find the probability that he or she also studies Chemistry.
\end{enumerate}
\end{exercice}

\begin{exercice}[Committee selection]
A committee of 4 people is to be chosen at random from 5 men and 6 women. Using combinations, find the probability that the committee contains:
\begin{enumerate}
\item exactly 2 women;
\item at least one man.
\end{enumerate}
\end{exercice}

\begin{exercice}[Rainy weekend in Buea]
The probability that it rains on any day in Buea in July is $0.6$, independently of other days. Using a tree diagram for a 3-day weekend, find the probability that:
\begin{enumerate}
\item it rains on exactly one day;
\item it rains on at least one day.
\end{enumerate}
\end{exercice}

\courssub{I prepare for the GCE}

\begin{exercice}[Two dice and independence]
Two fair dice are thrown and the scores are added.
\begin{enumerate}
\item List the outcomes for which the sum is greater than 8, and show that $P(\text{sum}>8)=\dfrac{5}{18}$. \hfill (3 marks)
\item Find the probability that at least one die shows a 6. \hfill (2 marks)
\item Given that the sum is greater than 8, find the probability that at least one die shows a 6. \hfill (3 marks)
\item Are the events ``sum greater than 8'' and ``at least one 6'' independent? Justify your answer carefully. \hfill (2 marks)
\end{enumerate}
\end{exercice}

\begin{exercice}[Defective items and Bayes' theorem]
Three machines $A$, $B$ and $C$ produce $50\%$, $30\%$ and $20\%$ respectively of a factory's output. The percentages of defective items produced by $A$, $B$ and $C$ are $2\%$, $4\%$ and $5\%$ respectively.
\begin{enumerate}
\item Draw a fully labelled tree diagram representing this situation. \hfill (3 marks)
\item Show that the probability that a randomly chosen item is defective is $0.032$. \hfill (3 marks)
\item An item is chosen at random and found to be defective. Use Bayes' theorem to find the probability that it came from machine $A$. \hfill (3 marks)
\item From which machine is a defective item most likely to have come? Justify. \hfill (2 marks)
\end{enumerate}
\end{exercice}

\begin{exercice}[Drawing cards without replacement]
Cards are drawn one by one, at random and without replacement, from a full pack of 52 playing cards.
\begin{enumerate}
\item Find the probability that the first two cards drawn are both aces. \hfill (3 marks)
\item Find the probability that the first heart appears on the third draw. \hfill (4 marks)
\item Find the probability that the second card drawn is a heart, given that the first card drawn was a heart. \hfill (2 marks)
\item Using combinations, find the probability that a hand of 5 cards contains exactly 3 hearts. \hfill (3 marks)
\end{enumerate}
\end{exercice}

\newpage

\coursec{Solutions to the exercises}

\begin{corrige}[Exercise 1 — Multiple choice questions]
\begin{enumerate}
\item The prime numbers on a fair die are 2, 3, 5. There are 3 favourable outcomes out of 6 equally likely outcomes. Hence $P(\text{prime}) = \frac{3}{6} = \frac{1}{2}$. The correct answer is \textbf{(c)}.
\item For mutually exclusive events, $P(A \cup B) = P(A) + P(B) = 0.3 + 0.5 = 0.8$. The correct answer is \textbf{(b)}.
\item For independent events, $P(A \cap B) = P(A) \times P(B) = 0.4 \times 0.25 = 0.1$. The correct answer is \textbf{(b)}.
\item Let $K$ be the event ``king'' and $H$ be the event ``heart''. There are 4 kings and 13 hearts, but the king of hearts is counted twice. So $P(K \cup H) = \frac{4}{52} + \frac{13}{52} - \frac{1}{52} = \frac{16}{52} = \frac{4}{13}$. The correct answer is \textbf{(c)}.
\end{enumerate}
\end{corrige}

\begin{corrige}[Exercise 2 — True or false]
\begin{enumerate}
\item \textbf{False.} Mutually exclusive events cannot occur simultaneously, so $P(A \cap B) = 0$. For independence we would need $P(A \cap B) = P(A)P(B)$. If $P(A)>0$ and $P(B)>0$, then $P(A)P(B) > 0$, so they cannot be independent. Counter‑example: rolling a die, $A=\{1,2\}$, $B=\{3,4\}$ are mutually exclusive but not independent.
\item \textbf{True.} By definition of the complement, $A$ and $A'$ partition the sample space, so $P(A) + P(A') = 1$.
\item \textbf{True.} $P(A \mid B) = \frac{P(A \cap B)}{P(B)}$. If $P(A \mid B) = P(A)$, then $P(A \cap B) = P(A)P(B)$, which is exactly the definition of independence (provided $P(B)>0$).
\item \textbf{False.} Empirical probability is based on observed frequencies and varies from experiment to experiment. It approaches the theoretical probability as the number of trials increases (Law of Large Numbers), but it is not always equal.
\end{enumerate}
\end{corrige}

\begin{corrige}[Exercise 3 — Definitions]
\begin{enumerate}
\item \textbf{Random experiment:} An experiment whose outcome cannot be predicted with certainty before it is performed, but whose set of all possible outcomes is known. \textbf{Sample space ($S$):} The set of all possible outcomes of a random experiment. \emph{Example:} Tossing a fair coin. The experiment is random; the sample space is $S = \{H, T\}$.
\item \textbf{Mutually exclusive events:} Two events $A$ and $B$ are mutually exclusive (or disjoint) if they cannot occur at the same time, i.e. $A \cap B = \emptyset$. \emph{Example:} In a single die roll, $A = \{\text{even}\}$ and $B = \{\text{odd}\}$ are mutually exclusive.
\item \textbf{Exhaustive events:} A set of events $\{E_1, E_2, \dots, E_n\}$ is exhaustive if their union is the whole sample space, i.e. $E_1 \cup E_2 \cup \dots \cup E_n = S$. \emph{Example:} In a die roll, $E_1 = \{1,2,3\}$, $E_2 = \{4,5,6\}$ are exhaustive.
\item \textbf{Independent events:} Two events $A$ and $B$ are independent if the occurrence of one does not affect the probability of the other, i.e. $P(A \cap B) = P(A)P(B)$. \emph{Example:} Tossing two fair coins. Let $A$ = ``first coin is heads'', $B$ = ``second coin is heads''. Then $P(A \cap B) = \frac{1}{4} = \frac{1}{2} \times \frac{1}{2} = P(A)P(B)$.
\item \textbf{Conditional probability $P(A \mid B)$:} The probability of event $A$ occurring given that event $B$ has already occurred, defined as $P(A \mid B) = \frac{P(A \cap B)}{P(B)}$ for $P(B) > 0$. \emph{Example:} Draw a card from a standard deck. $A$ = ``card is a king'', $B$ = ``card is a face card''. $P(A \mid B) = \frac{4/52}{12/52} = \frac{1}{3}$.
\end{enumerate}
\end{corrige}

\begin{corrige}[Exercise 4 — Direct calculations]
\begin{enumerate}
\item The word \textbf{CAMEROON} has 8 letters: C, A, M, E, R, O, O, N. Vowels: A, E, O, O (4 vowels). $P(\text{vowel}) = \frac{4}{8} = \frac{1}{2}$.
\item Sample space for two fair coins: $S = \{HH, HT, TH, TT\}$. Exactly one head: $\{HT, TH\}$. $P(\text{exactly one head}) = \frac{2}{4} = \frac{1}{2}$.
\item $P(A \cup B) = P(A) + P(B) - P(A \cap B) = 0.6 + 0.5 - 0.2 = 0.9$. $P(A \mid B) = \frac{P(A \cap B)}{P(B)} = \frac{0.2}{0.5} = 0.4$.
\item Empirical probability of heads = $\frac{118}{200} = 0.59$. Theoretical probability = $0.5$. As the number of tosses increases, the empirical probability tends to approach the theoretical probability (Law of Large Numbers).
\end{enumerate}
\end{corrige}

\begin{corrige}[Exercise 5 — Balls of the same colour]
Bag: 6 red (R), 4 black (B). Total = 10.
\begin{enumerate}
\item \textbf{With replacement:} Draws are independent.
$P(\text{both red}) = \frac{6}{10} \times \frac{6}{10} = \frac{36}{100} = 0.36$.
$P(\text{both black}) = \frac{4}{10} \times \frac{4}{10} = \frac{16}{100} = 0.16$.
$P(\text{same colour}) = 0.36 + 0.16 = 0.52 = \frac{13}{25}$.
\item \textbf{Without replacement:} Draws are dependent.
$P(\text{both red}) = \frac{6}{10} \times \frac{5}{9} = \frac{30}{90} = \frac{1}{3}$.
$P(\text{both black}) = \frac{4}{10} \times \frac{3}{9} = \frac{12}{90} = \frac{2}{15}$.
$P(\text{same colour}) = \frac{1}{3} + \frac{2}{15} = \frac{5}{15} + \frac{2}{15} = \frac{7}{15} \approx 0.4667$.
\end{enumerate}
\textbf{Comment:} The probability is higher with replacement ($0.52$) than without ($\approx 0.467$). With replacement, the composition of the bag remains the same for the second draw, keeping the probability of drawing a red ball high. Without replacement, drawing a red ball first reduces the number of red balls left, making a second red less likely; similarly for black. The effect of ``without replacement'' reduces the chance of two balls of the same colour because the first draw changes the proportions.
\end{corrige}

\begin{corrige}[Exercise 6 — Physics and Chemistry students]
Let $P$ = studies Physics, $C$ = studies Chemistry. Total students = 40.
$n(P) = 18$, $n(C) = 15$, $n(P \cap C) = 7$.
\begin{enumerate}
\item \textbf{Venn diagram:} Two overlapping circles. Region $P \cap C$: 7. Region $P$ only: $18-7=11$. Region $C$ only: $15-7=8$. Outside both: $40 - (11+7+8) = 14$.
\begin{popfigure}
\begin{tikzpicture}
\draw (0,0) circle (1.5cm) node[left=1.8cm] {$P$};
\draw (3,0) circle (1.5cm) node[right=1.8cm] {$C$};
\node at (-0.5,0) {11};
\node at (1.5,0) {7};
\node at (3.5,0) {8};
\node at (1.5,-2) {14};
\end{tikzpicture}
\legende{Venn diagram for Physics and Chemistry students}
\end{popfigure}
\item $P(P \cup C) = \frac{n(P \cup C)}{40} = \frac{11+7+8}{40} = \frac{26}{40} = \frac{13}{20} = 0.65$.
\item $P(\text{neither}) = \frac{14}{40} = \frac{7}{20} = 0.35$.
\item $P(C \mid P) = \frac{P(C \cap P)}{P(P)} = \frac{7/40}{18/40} = \frac{7}{18}$.
\end{enumerate}
\end{corrige}

\begin{corrige}[Exercise 7 — Committee selection]
Total people = 11 (5 men, 6 women). Committee of 4 chosen at random. Total number of possible committees: $\binom{11}{4} = 330$.
\begin{enumerate}
\item Exactly 2 women $\Rightarrow$ 2 women and 2 men.
Number of ways = $\binom{6}{2} \times \binom{5}{2} = 15 \times 10 = 150$.
$P(\text{exactly 2 women}) = \frac{150}{330} = \frac{5}{11}$.
\item At least one man = complement of ``no men'' (i.e., all 4 women).
Number of all-women committees = $\binom{6}{4} = 15$.
$P(\text{at least one man}) = 1 - \frac{15}{330} = 1 - \frac{1}{22} = \frac{21}{22}$.
\end{enumerate}
\end{corrige}

\begin{corrige}[Exercise 8 — Rainy weekend in Buea]
$P(\text{rain on a day}) = 0.6$, $P(\text{no rain}) = 0.4$. Three independent days.
\begin{popfigure}
\begin{tikzpicture}[level distance=1.5cm, sibling distance=2cm]
\node {Start}
child {node {R (0.6)}
child {node {R (0.6)}
child {node {R (0.6)} edge from parent node[left] {R}}
child {node {N (0.4)} edge from parent node[right] {N}}
edge from parent node[left] {R}
}
child {node {N (0.4)}
child {node {R (0.6)} edge from parent node[left] {R}}
child {node {N (0.4)} edge from parent node[right] {N}}
edge from parent node[right] {N}
}
}
child {node {N (0.4)}
child {node {R (0.6)}
child {node {R (0.6)} edge from parent node[left] {R}}
child {node {N (0.4)} edge from parent node[right] {N}}
edge from parent node[left] {R}
}
child {node {N (0.4)}
child {node {R (0.6)} edge from parent node[left] {R}}
child {node {N (0.4)} edge from parent node[right] {N}}
edge from parent node[right] {N}
}
};
\end{tikzpicture}
\legende{Tree diagram for 3-day weekend (R = rain, N = no rain)}
\end{popfigure}
\begin{enumerate}
\item Exactly one day of rain: sequences RNN, NRN, NNR.
Each has probability $0.6 \times 0.4 \times 0.4 = 0.096$.
Total $P = 3 \times 0.096 = 0.288$.
\item At least one day of rain = complement of no rain at all (NNN).
$P(\text{NNN}) = 0.4^3 = 0.064$.
$P(\text{at least one rain}) = 1 - 0.064 = 0.936$.
\end{enumerate}
\end{corrige}

\begin{corrige}[Exercise 9 — Two dice and independence]
Two fair dice: sample space of 36 equally likely ordered pairs $(i,j)$.
\begin{enumerate}
\item Sum > 8 means sum = 9, 10, 11, 12.
Outcomes:
Sum 9: (3,6), (4,5), (5,4), (6,3) → 4 outcomes.
Sum 10: (4,6), (5,5), (6,4) → 3 outcomes.
Sum 11: (5,6), (6,5) → 2 outcomes.
Sum 12: (6,6) → 1 outcome.
Total favourable = 4+3+2+1 = 10.
$P(\text{sum}>8) = \frac{10}{36} = \frac{5}{18}$. \hfill (3 marks)
\item At least one die shows a 6.
Complement: no 6 on either die. Each die has 5 non-6 faces → $5 \times 5 = 25$ outcomes.
$P(\text{at least one 6}) = 1 - \frac{25}{36} = \frac{11}{36}$. \hfill (2 marks)
\item Given sum > 8, we restrict to the 10 outcomes listed above. Among these, which have at least one 6?
Sum 9: (3,6), (6,3) → 2 outcomes.
Sum 10: (4,6), (6,4) → 2 outcomes.
Sum 11: (5,6), (6,5) → 2 outcomes.
Sum 12: (6,6) → 1 outcome.
Total = 7 outcomes.
$P(\text{at least one 6} \mid \text{sum}>8) = \frac{7}{10} = 0.7$. \hfill (3 marks)
\item Check independence: $P(\text{sum}>8) = \frac{5}{18} \approx 0.2778$, $P(\text{at least one 6}) = \frac{11}{36} \approx 0.3056$.
Product = $\frac{5}{18} \times \frac{11}{36} = \frac{55}{648} \approx 0.0849$.
$P(\text{sum}>8 \cap \text{at least one 6}) = \frac{7}{36} \approx 0.1944$.
Since $\frac{7}{36} \neq \frac{55}{648}$, the events are \textbf{not independent}. \hfill (2 marks)
\end{enumerate}
\textbf{Mark scheme / Examiner expectations:}
\begin{itemize}
\item Part (i): List outcomes clearly (3 marks: 1 for listing, 1 for count, 1 for probability).
\item Part (ii): Use complement or direct count (2 marks).
\item Part (iii): Correct conditional probability using reduced sample space (3 marks).
\item Part (iv): State definition of independence, compute both sides, conclude not equal (2 marks).
\end{itemize}
\end{corrige}

\begin{corrige}[Exercise 10 — Defective items and Bayes' theorem]
Machines: $A$ (50\%), $B$ (30\%), $C$ (20\%). Defective rates: $P(D|A)=0.02$, $P(D|B)=0.04$, $P(D|C)=0.05$.
\begin{enumerate}
\item \textbf{Tree diagram:}
\begin{popfigure}
\begin{tikzpicture}[level distance=2cm, sibling distance=3cm]
\node {Start}
child {node {$A$ (0.5)}
child {node {$D$ (0.02)} edge from parent node[left] {$D$}}
child {node {$D'$ (0.98)} edge from parent node[right] {$D'$}}
edge from parent node[left] {$A$}
}
child {node {$B$ (0.3)}
child {node {$D$ (0.04)} edge from parent node[left] {$D$}}
child {node {$D'$ (0.96)} edge from parent node[right] {$D'$}}
edge from parent node[midway,above] {$B$}
}
child {node {$C$ (0.2)}
child {node {$D$ (0.05)} edge from parent node[left] {$D$}}
child {node {$D'$ (0.95)} edge from parent node[right] {$D'$}}
edge from parent node[right] {$C$}
};
\end{tikzpicture}
\legende{Tree diagram for machine output and defectives}
\end{popfigure}
\hfill (3 marks)
\item Total probability of defective:
$P(D) = P(A)P(D|A) + P(B)P(D|B) + P(C)P(D|C)$
$= 0.5 \times 0.02 + 0.3 \times 0.04 + 0.2 \times 0.05$
$= 0.01 + 0.012 + 0.01 = 0.032$. \hfill (3 marks)
\item Bayes' theorem for $P(A|D)$:
$P(A|D) = \frac{P(A)P(D|A)}{P(D)} = \frac{0.5 \times 0.02}{0.032} = \frac{0.01}{0.032} = \frac{10}{32} = \frac{5}{16} = 0.3125$. \hfill (3 marks)
\item Compare $P(A|D)=0.3125$, $P(B|D)=\frac{0.3 \times 0.04}{0.032} = \frac{0.012}{0.032} = 0.375$, $P(C|D)=\frac{0.2 \times 0.05}{0.032} = \frac{0.01}{0.032} = 0.3125$.
Machine $B$ has the highest posterior probability (0.375). So a defective item is most likely to have come from machine $B$. \hfill (2 marks)
\end{enumerate}
\textbf{Mark scheme / Examiner expectations:}
\begin{itemize}
\item Tree diagram: branches labelled with probabilities, all outcomes shown (3 marks).
\item Total probability: correct formula and arithmetic (3 marks).
\item Bayes' theorem: correct substitution and simplification (3 marks).
\item Comparison of posterior probabilities, clear conclusion (2 marks).
\end{itemize}
\end{corrige}

\begin{corrige}[Exercise 11 — Drawing cards without replacement]
Standard deck: 52 cards, 4 aces, 13 hearts.
\begin{enumerate}
\item First two cards both aces.
$P(\text{1st ace}) = \frac{4}{52}$, $P(\text{2nd ace} \mid \text{1st ace}) = \frac{3}{51}$.
$P(\text{both aces}) = \frac{4}{52} \times \frac{3}{51} = \frac{12}{2652} = \frac{1}{221}$. \hfill (3 marks)
\item First heart appears on the third draw.
This means: 1st not heart, 2nd not heart, 3rd heart.
$P(\text{1st not heart}) = \frac{39}{52} = \frac{3}{4}$.
$P(\text{2nd not heart} \mid \text{1st not heart}) = \frac{38}{51}$.
$P(\text{3rd heart} \mid \text{first two not hearts}) = \frac{13}{50}$.
Product: $\frac{39}{52} \times \frac{38}{51} \times \frac{13}{50} = \frac{39 \times 38 \times 13}{52 \times 51 \times 50}$.
Simplify: $\frac{39}{52} = \frac{3}{4}$, $\frac{38}{51} = \frac{38}{51}$, $\frac{13}{50} = \frac{13}{50}$.
$= \frac{3 \times 38 \times 13}{4 \times 51 \times 50} = \frac{1482}{10200} = \frac{247}{1700} \approx 0.1453$. \hfill (4 marks)
\item $P(\text{2nd heart} \mid \text{1st heart})$. After drawing one heart, 51 cards remain, 12 hearts.
$P = \frac{12}{51} = \frac{4}{17}$. \hfill (2 marks)
\item Hand of 5 cards, exactly 3 hearts.
Number of ways: choose 3 hearts from 13, and 2 non-hearts from 39.
$\binom{13}{3} \times \binom{39}{2} = 286 \times 741 = 211926$.
Total 5-card hands: $\binom{52}{5} = 2598960$.
$P = \frac{211926}{2598960} = \frac{\binom{13}{3}\binom{39}{2}}{\binom{52}{5}} \approx 0.0815$. \hfill (3 marks)
\end{enumerate}
\textbf{Mark scheme / Examiner expectations:}
\begin{itemize}
\item (i) Correct conditional probability multiplication (3 marks).
\item (ii) Sequence of three dependent events, correct fractions (4 marks).
\item (iii) Simple conditional probability after one heart removed (2 marks).
\item (iv) Correct use of combinations, numerator and denominator (3 marks).
\end{itemize}
\end{corrige}

\newpage

\coursec{Summary sheet}

\begin{popfigure}
\begin{tikzpicture}[
  grow cyclic,
  mindmap,
  concept color=popblue,
  every node/.style={concept, font=\small},
  root concept/.append style={concept color=popdark, font=\bfseries, text=white, minimum size=2.6cm},
  level 1 concept/.append style={level distance=3.4cm, sibling angle=51.5, font=\scriptsize, minimum size=2.1cm},
  level 2 concept/.append style={level distance=2.2cm, font=\tiny, minimum size=1.5cm},
]
\node[root concept]{Elementary Probability}
  child[concept color=popblue]{ node{Experiments \& Sample Spaces}
    child[concept color=popblue!70]{ node{outcome, $S$} }
    child[concept color=popblue!70]{ node{events, types} }
  }
  child[concept color=popgreen]{ node{Classical \& Empirical}
    child[concept color=popgreen!70]{ node{$P(E)=\frac{n(E)}{n(S)}$} }
    child[concept color=popgreen!70]{ node{relative frequency} }
  }
  child[concept color=poporange]{ node{Addition Laws}
    child[concept color=poporange!70]{ node{exclusive: $P(A)+P(B)$} }
    child[concept color=poporange!70]{ node{general: $-P(A\cap B)$} }
  }
  child[concept color=poppink]{ node{Multiplication \& Conditional}
    child[concept color=poppink!70]{ node{$P(A\cap B)=P(A)P(B\mid A)$} }
    child[concept color=poppink!70]{ node{independence} }
  }
  child[concept color=poppurple]{ node{Tree Diagrams}
    child[concept color=poppurple!70]{ node{with replacement} }
    child[concept color=poppurple!70]{ node{without replacement} }
  }
  child[concept color=popgold]{ node{Counting Techniques}
    child[concept color=popgold!70]{ node{permutations ${}^nP_r$} }
    child[concept color=popgold!70]{ node{combinations $\binom{n}{r}$} }
  }
  child[concept color=popink]{ node{Total Probability \& Bayes}
    child[concept color=popink!70]{ node{partition of $S$} }
    child[concept color=popink!70]{ node{$P(A_i\mid B)$} }
  };
\end{tikzpicture}
\legende{Mind map of the chapter Elementary Probability}
\end{popfigure}

\begin{bilan}
\textbf{Key definitions.}
\begin{itemize}
\item A \cle{random experiment} is a process whose outcome cannot be predicted with certainty; an \cle{outcome} is a single possible result; the \cle{sample space} $S$ is the set of all outcomes.
\item An \cle{event} is a subset of $S$. Events may be \emph{simple}, \emph{compound}, \emph{sure} ($S$), \emph{impossible} ($\varnothing$), \emph{complementary} ($A'$), \emph{mutually exclusive} ($A\cap B=\varnothing$), or \emph{exhaustive} (their union is $S$).
\item \cle{Classical probability}: for equally likely outcomes, $P(E)=\dfrac{n(E)}{n(S)}$.
\item \cle{Empirical probability}: $P(E)\approx \dfrac{\text{frequency of }E}{\text{number of trials}}$ for a large number of trials.
\item \cle{Conditional probability}: $P(A\mid B)=\dfrac{P(A\cap B)}{P(B)}$, $P(B)>0$.
\item $A$ and $B$ are \cle{independent} when $P(A\cap B)=P(A)\,P(B)$, equivalently $P(A\mid B)=P(A)$.
\end{itemize}

\textbf{Laws and formulas to know.}
\begin{itemize}
\item $0\le P(E)\le 1$; $P(S)=1$; $P(\varnothing)=0$; $P(A')=1-P(A)$.
\item Addition law (mutually exclusive): $P(A\cup B)=P(A)+P(B)$.
\item Addition law (general): $P(A\cup B)=P(A)+P(B)-P(A\cap B)$.
\item Multiplication law: $P(A\cap B)=P(A)\,P(B\mid A)=P(B)\,P(A\mid B)$; for independent events $P(A\cap B)=P(A)P(B)$.
\item Total probability: if $A_1,\dots,A_n$ partition $S$, then $P(B)=\sum_i P(A_i)\,P(B\mid A_i)$.
\item Bayes' theorem: $P(A_i\mid B)=\dfrac{P(A_i)\,P(B\mid A_i)}{\sum_j P(A_j)\,P(B\mid A_j)}$.
\item Counting: ${}^nP_r=\dfrac{n!}{(n-r)!}$ (order matters); $\binom{n}{r}=\dfrac{n!}{r!(n-r)!}$ (order does not matter); $P(E)=\dfrac{\text{favourable selections}}{\text{total selections}}$.
\end{itemize}

\textbf{Methods.}
\begin{itemize}
\item Always start by writing the sample space (or its size) and checking that outcomes are equally likely.
\item For sequential experiments, draw a \cle{tree diagram}: multiply along branches, add between branches ending in the required event.
\item \emph{With replacement}: probabilities stay constant along the tree (independent stages). \emph{Without replacement}: denominators decrease by one at each stage.
\item ``At least one'' problems: use the complement, $P(\text{at least one})=1-P(\text{none})$.
\item For selection problems, decide first whether order matters (permutation) or not (combination).
\end{itemize}

\textbf{Common mistakes.}
\begin{itemize}
\item Adding $P(A)+P(B)$ without checking that the events are mutually exclusive.
\item Confusing independence ($P(A\cap B)=P(A)P(B)$) with mutual exclusivity ($P(A\cap B)=0$): exclusive events with positive probabilities are never independent.
\item Forgetting to reduce the denominator in selection without replacement.
\item Inverting the condition: $P(A\mid B)\ne P(B\mid A)$ in general.
\item Probabilities on a tree branch not summing to $1$ at each node.
\end{itemize}

\textbf{GCE tips.}
\begin{itemize}
\item Give final answers as simplified fractions, or decimals to at least 3 s.f. when required.
\item In Bayes-type questions, define events clearly in words before computing; examiners award marks for correct notation.
\item Show the tree diagram: it earns method marks even if the final arithmetic slips.
\item Check your answer lies in $[0,1]$; a probability greater than $1$ signals an error.
\end{itemize}
\end{bilan}

\findecours

\end{document}
